\documentclass{article}
\usepackage{decomp}
\begin{document}
\sourcefile{2026.10.15.759299v1.full.pdf}
\decompnote{bioRxiv full-text PDF, version 1, 50 pages: text pp.\,1--34 (title and summary, introduction, results, discussion, figure legends, Extended Data legends, methods, acknowledgements, references) and figure pages 35--50. Every text page carries the footer ``bioRxiv preprint doi: https://doi.org/10.64898/2026.10.15.759299; this version posted October 16, 2026. The copyright holder for this preprint (which was not certified by peer review) is the author/funder, who has granted bioRxiv a license to display the preprint in perpetuity. It is made available under a CC-BY 4.0 International license.'' (quoted once here, left out below). Citations are printed as superscript numbers and are kept as such; the reference list keeps its original numbering.}
\pdfpage{1}
\section*{Rapid alternation of conscious and suppressed stimulus representations in the primate brain}

\noindent Janis K. Hesse\textsuperscript{1}, Frank F. Lanfranchi\textsuperscript{1,3,*,\ddag}, Conor McGrory\textsuperscript{2,\ddag}, Yuelin Shi\textsuperscript{1,3,\dag}, Tatiana A. Engel\textsuperscript{2}, Doris Y. Tsao\textsuperscript{1,4}


\textbf{Affiliations}

\textsuperscript{1} Department of Neuroscience, University of California, Berkeley, Berkeley, CA, USA\\
\textsuperscript{2} Princeton Neuroscience Institute, Princeton University, Princeton, NJ, USA\\
\textsuperscript{3} Division of Biology and Biological Engineering, California Institute of Technology, Pasadena, CA, USA\\
\textsuperscript{4} Howard Hughes Medical Institute, Berkeley, CA, USA

\textbf{Present addresses}

* Present address: Department of Neurobiology, Stanford University, Stanford, CA, USA\\
\dag{} Present address: University of California, San Francisco, San Francisco, CA, USA\\
\ddag: Frank F. Lanfranchi and Conor McGrory contributed equally.

Correspondence: Janis Hesse (\href{mailto:janishesse@googlemail.com}{janishesse@googlemail.com}), Doris Tsao (\href{mailto:dortsao@berkeley.edu}{dortsao@berkeley.edu})

\section*{Summary}

The neural correlates of consciousness remain elusive\textsuperscript{1}. Previous studies found that high-level visual cortex is weakly modulated by conscious perception\textsuperscript{2,3}, yet a pure signal of conscious perception has not been found anywhere in the brain. Here, we employed Neuropixels recordings across multiple face patches in inferotemporal and prefrontal cortex during a binocular flash suppression paradigm\textsuperscript{4}, in which one of two images is consciously perceived while the other, though continuously present, is suppressed from awareness. We found that, during a constant stimulus and percept, the brain rapidly alternates, about 15 times per second, between conscious and suppressed micro-epochs. During conscious micro-epochs, a pure signal of conscious perception was present, allowing accurate reconstruction of the conscious image from neural activity. During suppressed micro-epochs this reversed: the pure signal now tracked the suppressed stimulus, allowing accurate reconstruction of the suppressed image. At the single-cell level, neurons responded selectively to the conscious stimulus and not to the suppressed stimulus during conscious micro-epochs; conversely, during suppressed micro-epochs, neurons responded selectively to the suppressed stimulus. Despite this apparent symmetry, the two epochs could be reliably distinguished, with feature tuning amplified during conscious micro-epochs. Recordings in prefrontal cortex yielded qualitatively similar results. Together, these findings show that both inferotemporal and prefrontal cortex rapidly interleave conscious with suppressed content, indicating that conscious access requires temporal selection beyond localization.

\pdfpage{2}\section*{Introduction}

At any moment, only some of the brain's activity reaches conscious awareness, while the rest is processed without ever being experienced. How this division arises is among the most basic unsolved problems in neuroscience. We set aside the ``hard problem'' of why neural activity feels like anything, and ask: what distinguishes neural activity that is consciously accessible from activity that is not?

Answering this question requires a situation in which the physical input stays fixed while conscious perception changes, so that differences in neural activity can be attributed to perception rather than to the stimulus. Binocular rivalry is a paradigm that provides exactly this dissociation\textsuperscript{5,6}. When conflicting images are presented separately to the two eyes, observers consciously perceive only one image at a time, despite both images remaining physically present. In binocular flash suppression, the dissociation is made experimentally controllable\textsuperscript{3,4,7}: when one image is presented alone to one eye, and then a second image is abruptly introduced to the other eye, observers typically perceive the new image (\textbf{Fig. 1a}).

Rivalry studies have shown that neural activity in the primate brain is modulated by conscious perception: a cell that responds more to faces than houses will, under rivalry, fire more when the face is perceived than when the house is perceived, even though both images are continuously present. Yet this perceptual modulation poses a persistent puzzle: it has consistently been found to be weaker than the modulation the same cell shows when the face and house are simply presented one at a time\textsuperscript{2,3,8,9}. Moreover, neurons continue to carry information about the suppressed image\textsuperscript{2}. Thus, previous studies \textit{have not revealed a pure neural representation of conscious perception, i.e., one uncontaminated by suppressed representation and closely resembling the representation of a perceptually-equivalent unambiguous stimulus, in any recorded region of the brain during rivalry.}

This absence of a pure correlate has contributed to an ongoing debate over whether the neural correlates of consciousness reside in the front or the back of the brain, with frontal theories such as Global Neuronal Workspace and Higher-Order theories on one side and posterior theories such as Integrated Information Theory and Recurrent Processing Theory on the other\textsuperscript{10}. Recent large-scale efforts, including the COGITATE adversarial collaboration\textsuperscript{11}, have tested these competing predictions, yet no definitive anatomical locus has emerged.

Conscious perception is a dynamic process. A given sensory input can be interpreted in multiple ways, and perception actively resolves which interpretation we experience at any moment\textsuperscript{12}. Rivalry makes this especially stark: with the physical input held constant, perception settles into one of two interpretations and alternates over time. Understanding this process therefore requires asking \textit{how different stimulus representations are dynamically organized in time} within cortical populations. Previous studies using conventional single-unit recordings have generally required averaging across trials, because the activity of individual neurons is too sparse and variable to resolve moment-to-moment representational dynamics within a single trial. Addressing this requires tracking the coordinated activity of many neurons in real time while presenting enough stimuli to characterize what each momentary population state represents.

\pdfpage{3}Here, we examined how conscious and suppressed representations are organized in time, using large-scale Neuropixels recordings\textsuperscript{13} from macaque face patches in inferotemporal and prefrontal cortex during binocular flash suppression. Presenting a large set of face stimuli allowed us to reconstruct the specific face stimuli represented by the population\textsuperscript{14}, while the use of Neuropixels probes allowed us to track changes in representations within individual trials, at a resolution not accessible with single-neuron recordings. We found that even during a stable percept and constant stimulus, inferotemporal population activity rapidly alternated between two states, one purely representing the consciously perceived face and the other purely representing the suppressed face. Conscious versus suppressed face epochs could be reliably decoded, with higher gain for representation of the conscious face during conscious epochs than for the suppressed face during suppressed epochs. Prefrontal cortex showed the same interleaved structure. These findings show that conscious and suppressed representations are temporally multiplexed within cortical population activity; both inferotemporal and prefrontal cortex carry a rich representation of suppressed stimuli that are not consciously perceived, indicating that selective access to the conscious percept requires a temporal selection mechanism.

\section*{Results}

\subsection*{Binocular flash suppression dissociates conscious perception from the physical stimulus}

To dissociate conscious perception from the physical stimulus, we used binocular flash suppression (\textbf{Fig. 1a}). A priming image is shown to one eye, and after an extended period (1 or 2 s) a second image is flashed to the other eye; the flashed image dominates perception while the priming image is suppressed from awareness.

We first verified the effectiveness of this paradigm in human participants. Participants reported by button press whether they perceived the image presented to the left eye, the image presented to the right eye, or a mixed percept. Shortly after the flash, participants predominantly reported perceiving the newly introduced stimulus, with few reports of the priming or mixed percept ($p = 4 \times 10^{-7}$, t-test on proportion of flashed versus priming percept across 14 participants, \textbf{Fig. 1b}, see Methods). For all subsequent analyses, we focus on the period 300 to 1000 ms after the flash, where the flashed stimulus dominated perception and the priming stimulus was suppressed.

\subsection*{Average neural activity is only weakly modulated by conscious perception}

We performed large-scale electrophysiological recordings while rhesus macaque monkeys viewed the same flash suppression displays verified in humans (\textbf{Fig. 1c}). NHP Neuropixels probes were inserted into multiple nodes of the macaque brain using a custom insertion mechanism (\textbf{Fig. 1d}): face patches ML and AM in inferotemporal cortex, and the orbitofrontal (PO) and ventrolateral (PVL) regions of prefrontal cortex. Data were recorded from two monkeys; throughout the paper, we show data pooled across the two animals in the main figures; data for individual animals are shown in the extended data. Face patches were localized with fMRI using a faces-versus-objects contrast (\textbf{Extended Data Fig. 1}). A majority of cells in ML and AM were face-selective (\textbf{Fig. 1f}, 91\% and 71\% of visually responsive cells for the example session shown; \pdfpage{4}faces vs. objects, t-test, $p < 0.05$). A single trial of simultaneously recorded ML and AM activity is shown in \textbf{Fig. 1g}.

We first focus on inferotemporal cortex. We asked whether single neurons there purely represented the consciously perceived stimulus. For each cell, we compared its responses to the two images of a stimulus pair, A and B, using a modulation index $MI = \frac{R_A - R_B}{R_A + R_B}$, where $R_A$ and $R_B$ are the cell's average firing rates to stimulus A and stimulus B. We compared modulation during rivalry with modulation during an unambiguous alternation condition. During the unambiguous condition, stimuli were presented alone after priming rather than in binocular conflict, providing a measure of the cell's visual selectivity. During binocular flash suppression the same two stimuli were presented in conflict, and we computed the index with respect to the consciously dominant stimulus, so that $R_A$ and $R_B$ now reflect responses when stimulus A or stimulus B, respectively, was perceived. If activity during flash suppression purely reflected the conscious percept, modulation in the binocular condition would match modulation during unambiguous priming. Instead, modulation during binocular flash suppression was much weaker than during unambiguous alternation (\textbf{Fig. 1h}, $p = 8 \times 10^{-133}$, paired t-test over 5550 cell-stimulus pairs). Virtually all neurons that strongly discriminated the two stimuli during unambiguous presentation showed much smaller differences when the same stimuli were presented in binocular conflict, with nearly all points falling below the identity diagonal (slope = 0.23). This weak modulation is puzzling: perceptually, one image is clearly seen while the other is absent from awareness, yet the neural responses only weakly distinguished the two.

A clue to the origin of this weak modulation comes from an earlier finding that the suppressed stimulus, though absent from awareness, can still be decoded from inferotemporal activity during rivalry\textsuperscript{2}. In that study, however, only three stimuli were used (and previous studies generally contrasted only two rival alternatives\textsuperscript{3,8,9,15}), leaving opaque what, if any, details of either the conscious or the suppressed face are represented during rivalry. We therefore set out to characterize the detailed content of conscious and suppressed representations. The face patch system is well suited to this: face identity can be described in a parametric feature space to which cells are linearly tuned, so that face features can be reconstructed from neural responses by linear regression\textsuperscript{14}. To precisely map the feature code of each cell, we presented 1500 faces in flash suppression. This allowed us to ask not merely whether a conscious or suppressed face was decodable, but precisely which features of each were represented at each moment.

\subsection*{Single-trial population activity can be partitioned into rapidly alternating states}

Having access to the represented content at each moment in time, we first asked whether there are epochs in which only the conscious or only the suppressed face is represented. First, we computed template vectors for the conscious and suppressed faces by taking the average neural responses during priming trials when each was shown alone. For each 10 ms time bin during rivalry, we then asked which template the observed spike counts more probably arose from (a Poisson log-likelihood ratio (LLR) decoder, see Methods), labelling the bin state 1 when the conscious face template was more likely and state 2 when the suppressed face template was more likely (\textbf{Fig. 2a}). This produced a time-resolved partition of each trial into two states. During \pdfpage{5}priming, when only one face was present, activity was assigned mostly to the state corresponding to the priming image (\textbf{Fig. 2b}). After the flash, however, the population rapidly alternated between the two states, and state 1 and state 2 were occupied about equally (\textbf{Fig. 2c}). To quantify how quickly the population alternated, we measured the autocorrelation of the moment-by-moment log-likelihood-ratio signal (\textbf{Fig. 2d}, black). It decayed with a time constant of 67 ms, whereas the same signal computed after shuffling time points decayed within a single bin (\textbf{Fig. 2d}, grey). Thus, the state transitions occurred more than an order of magnitude faster than the perceptual switches reported during binocular rivalry (see Methods and Results section ``Dynamics of state switching'' below for more details).

For stimulus pairs where each of the two faces was preferred by a distinct set of neurons, the alternation was visible both in the population (\textbf{Fig. 2e}) and in single-cell rasters (\textbf{Fig. 2f, g}).

\subsection*{Reconstructions from inferred states recover conscious and suppressed stimuli}

After partitioning the population activity into epochs whose representations more closely resemble either the conscious or the suppressed stimulus, we next asked what face features were represented during each inferred state. Exploiting the linear tuning of face cells to features of a parametric face space\textsuperscript{14}, we trained regression weights to predict face features from responses to priming stimuli (i.e., in the absence of binocular conflict) and then applied these weights to activity during state 1 and state 2 to attempt to reconstruct the face represented in each. We inferred states individually for each of 9 sessions (from two monkeys) and averaged responses of each neuron over all state 1 and, separately, over all state 2 epochs from all trials. We then pooled these state-averaged neural responses from ML and AM neurons across sessions to predict features. The predicted features were transformed into images for visualization.

Strikingly, reconstructions computed from activity during state 1 closely resembled the conscious face, whereas reconstructions computed from state 2 closely resembled the suppressed face (\textbf{Fig. 2h}). Thus, even though the suppressed face was not consciously perceived, a clear representation of it was present in IT cells during distinct epochs in time. These example reconstructions suggest that state 1 and state 2 carry separate representations of the conscious and suppressed face. By contrast, a time-averaged null model, in which features were predicted from activity averaged across the binocular period without state segmentation, produced reconstructions that matched neither the conscious nor the suppressed face, but rather their average (\textbf{Extended Data Fig. 2}).

To quantify reconstruction accuracy, we computed explained variance of each face feature, defined as the percentage of variance of the true values of that feature across faces explained by its reconstructed values (\textbf{Fig. 2i}). During state 1, prediction of the consciously perceived face reached over 80\% explained variance for some features. This was only slightly lower than the upper bound for this analysis, obtained by reconstructing faces from average priming responses, which reached over 90\% explained variance. The upper bound is informative because priming uses the same regression weights but without binocular conflict, so it shows how well faces can be reconstructed from this population under ideal conditions; state 1 approaches it despite the competing suppressed face. By contrast, explained variance for the suppressed face during state \pdfpage{6}1 was extremely low, comparable to chance level when that face was not present (\textbf{Fig. 2i}, orange dashed trace; compare \textbf{Extended Data Fig. 2b}, black dashed trace). Thus, during state 1 there is strong information about the consciously perceived face and little to none about the suppressed face.

During state 2, the pattern reversed: explained variance for the suppressed stimulus was now high, albeit slightly lower than for the conscious face during state 1 (\textbf{Fig. 2i, purple}). Explained variance for the conscious stimulus during state 2 fell to chance level as above. Thus, state 1 and state 2 separately encoded the conscious and suppressed faces ($p = 3 \times 10^{-29}$, t-test, see Methods).

As a complementary measure, we asked whether each reconstruction was accurate enough to identify the presented face among a number of randomly selected faces, following the decoding approach used previously to read out facial identity from face cell responses\textsuperscript{14}. For each reconstructed face, we compared its predicted features to the features of the correct face and to those of randomly chosen distractor faces, and counted the reconstruction as correctly decoded when it was closest to the correct face (\textbf{Fig. 2j}). During state 1, decoding accuracy was high for the conscious face and low for the suppressed face; during state 2, this reversed, with high accuracy for the suppressed face and low accuracy for the conscious face ($p = 10^{-5}$, permutation test, see Methods).

We also asked whether the state-resolved approach could help with decoding of stimulus identity on \textit{single trials} rather than state-resolved activity averaged across trials in a separate set of 10 experiments where we presented 10 stimuli per experiment. We chose the 10 stimuli from a screening set of 49 stimuli that maximized decodability based on a screening session that was conducted and analysed right before each main session. We presented 25 paired combinations of the 10 stimuli in binocular flash suppression. We then trained a 10-class linear decoder on time-averaged responses on individual priming trials to discriminate the 10 identities. We applied this decoder to single-trial binocular responses: within a 200-ms sliding window on each trial, we decoded separately from the time points assigned to state 1, to state 2, or from both states combined, and asked whether the predicted identity matched the conscious or suppressed stimulus (\textbf{Fig. 2k}; see Methods). Across sessions, decoding using average activity during the binocular condition was about equal for the conscious and suppressed stimuli (40\% and 34\%, respectively). During state 1, on the other hand, the conscious stimulus was decoded far more accurately than the suppressed stimulus (over 61\% versus 4\%; \textbf{Fig. 2l}, $p = 2 \times 10^{-6}$, paired t-test). During state 2, the pattern reversed ($p = 3 \times 10^{-7}$, paired t-test).

Together, these results indicate that population activity during rivalry can be partitioned into temporal epochs whose average activity represents either the conscious or the suppressed face. Because the regression weights were trained only on unambiguous priming responses, the accuracy of both reconstructions means that the average activity within each set of epochs reinstates the same feature code evoked when the face is presented alone.

However, because the same neural activity was used both to define the state partition and to evaluate representations within these states, the analyses above are subject to a circular \pdfpage{7}conditioning bias: the circular re-use of the same data samples allows each neuron's independent spiking fluctuations to influence both state assignment and its subsequent state-conditioned response. The resulting selection bias systematically sorts these fluctuations along the state-decoder direction, artifactually increasing the separation between state-conditioned population averages and inflating reconstruction accuracy. The state-conditioned responses therefore reflect a mixture of genuine coordinated population fluctuations and a bias introduced by the analysis. Consequently, finding a partition in which the state-conditioned average responses separately reconstruct the conscious and suppressed faces does not establish that the population undergoes coordinated switching between these two representations.

\subsection*{Coordinated population states alternate between conscious and suppressed stimulus representations}

We therefore tested whether population activity exhibits coordinated transitions between representations of conscious and suppressed stimuli. We asked first whether a two-state model captured population activity better than the best-fitting single-state model, obtained by maximum likelihood estimation (see Methods). The two-state model assigned higher likelihood to the observed population activity than the best single-state model for all 19 sessions (\textbf{Fig. 3a}, $p = 10^{-4}$, t-test). This indicates that the observed switching between two states reflects coordinated population dynamics rather than independent fluctuations of individual neurons.

We next tested whether the representational content alternated in a coordinated manner across the population. We split neurons into two random halves, using one half of the population to infer the state sequence, and the other half to predict face features. If the inferred state sequence captured a genuine coordinated alternation, the neurons held out for prediction should change their represented stimulus when alternating between inferred states.

Rasters of populations of held-out neurons and individual held-out neurons showed the same changes in activity at inferred state transitions (\textbf{Fig. 3b}--\textbf{d}). Most importantly, the held-out population reconstructed the conscious face during epochs classified as state 1 by the other half of the neurons, and the suppressed face during epochs classified as state 2 (\textbf{Fig. 3e--i}). Reconstructions using the held-out half of neurons were modestly less accurate than those obtained from the full population (as expected given the circularity bias described above) but continued to clearly separate conscious and suppressed face representations during state 1 and state 2 ($p = 6 \times 10^{-17}$, t-test, see Methods). The cross-validated results held true on an individual monkey level (\textbf{Extended Data Fig. 3}) and an individual session level (\textbf{Extended Data Fig. 4}). Together, these results demonstrate that the population genuinely alternates between conscious and suppressed representations.

We confirmed that the lower reconstruction accuracy for the held-out half was not due to using fewer neurons for state inference and face reconstructions. To do so, we inferred states from all but one neuron, used these states to compute the state-conditioned responses of the held-out neuron, and then pooled these responses across all held-out neurons to generate reconstructions (leave-one-neuron-out). This analysis removes circular conditioning without sacrificing population size. The leave-one-neuron-out validation increased face explained variance by only around 2 \pdfpage{8}percentage points relative to cross-validation using half of neurons, showing that state inference and reconstruction accuracy were not limited by neuron count (\textbf{Extended Data Fig. 5a--c}). On the other hand, restricting the split-half analysis to high-confidence time points increased explained variance by 10 percentage points (\textbf{Extended Data Fig. 5d--f}), showing that these moments contained high-fidelity representations of conscious and suppressed stimuli.

\subsection*{Control analyses confirm coordinated within-trial state alternation}

Having established that the population exhibits coordinated dynamics, we next controlled for two alternative hypotheses: the first alternative is that the states reflect differences between trials rather than within-trial switching; the second alternative is that interleaved representations of conscious and suppressed stimuli are imposed by the inference procedure itself.

We first asked whether the observed alternation between conscious and suppressed stimuli reflected genuine \textit{within-trial} switching or rather differences between trials. Even though the analysis was restricted to the 700 ms period in which human observers reliably perceive the flashed face, the two states could in principle arise from some trials being dominated by the conscious face and others by the suppressed face, rather than from switching within each trial. To rule this out, we subsampled time points within each trial so that every trial contributed equal numbers of samples to state 1 and state 2. This ensured that the two state-resolved reconstructions were based on matched numbers of samples from the same trials, eliminating differences between trials as a source of the state separation. Despite the fact that this procedure discarded many time points, state-resolved reconstructions remained robust (\textbf{Extended Data Figure 6}). Thus, the separation between conscious and suppressed representations cannot be explained by different stimuli being represented on different trials; instead, both representations are expressed within single trials in alternating micro-epochs.

Second, we randomly shuffled the response templates used to define the two states and repeated the reconstruction analysis. This tested whether our state-inference procedure could manufacture accurate reconstructions of any stimulus simply by sorting time points to whatever templates it was given. If it could, this procedure should yield reconstructions that match those shuffled templates even though they were not presented while recording the activity used for reconstruction. To the contrary, reconstructions obtained with randomly re-assigned templates did not resemble the templates, and their explained variance was largely below zero (\textbf{Extended Data Fig. 7a--c}). We also compared reconstructions obtained with shuffled templates to the faces that were actually presented on each trial (\textbf{Extended Data Fig. 7d--f}). These reconstructions were much worse than those obtained with the correct templates, indicating that accurate reconstruction depends on the template used for state inference matching the true face pair, and does not arise from arbitrary partitions of neural activity.

\subsection*{Single cells have higher gain of tuning during conscious compared to suppressed epochs}

The reconstruction results, together with the control analyses above, establish that the population alternates between representing the conscious and suppressed face. This raises the question of how the two representations differ at the single-cell level, and whether the conscious \pdfpage{9}representation carries any signature that distinguishes it from the suppressed one. We therefore asked how individual face cells were tuned to each face, and how that tuning changed between the two states. As in \textbf{Fig. 3}, we inferred states based on one half of neurons, and measured single-cell tuning in the held-out half of neurons.

Face cells exhibit axis tuning: they respond as a linear function of a face's position along a preferred axis in shape-appearance space\textsuperscript{14}. If this tuning is preserved during rivalry, a cell's responses should ramp along its preferred axis as the represented face varies. To test this, we ordered the 1500 binocular stimuli by the response each cell was predicted to give to one of the two faces, using weights fit to priming responses, and plotted the cell's actual binocular responses in that order; a cell representing that face should produce responses that ramp along the ordering of stimuli. We performed analyses separately for orderings by the conscious and the suppressed face, and separately for time points assigned to state 1 and state 2.

In individual cells, the ramp appeared for the conscious face during state 1 and for the suppressed face during state 2 (example cell in \textbf{Fig. 4a}). During state 1, a cell's responses increased steadily as the conscious face moved along its preferred axis, while ordering the same responses by the suppressed face produced no gradient. During state 2 this reversed: the same cell now ramped with the suppressed face and showed no gradient for the conscious face.

This pattern held across the population (\textbf{Fig. 4b}). A clear gradient was present for the conscious face during state 1 and for the suppressed face during state 2, with little gradient for the off-target face in each state. Rivalry therefore preserved each cell's underlying feature tuning but expressed it selectively for the face represented by the population at a given moment. Confirming this, the cell's preferred tuning axis computed for the rivalry responses to the conscious face during state 1 matched the tuning axis computed from priming responses, barely lower than noise ceiling (\textbf{Extended Data Figure 8}; $p = 0.01$, paired t-test). Similarly, the tuning axis from responses to the suppressed face during state 2 was indistinguishable from the priming tuning axis relative to noise ceiling ($p = 0.35$). On the other hand, responses to the suppressed face during state 1, and the conscious face during state 2 did not yield tuning axes consistent with the priming tuning axis (state 1 suppressed: $p = 10^{-87}$, state 2 conscious: $p = 6 \times 10^{-69}$).

Comparing gain of tuning, measured by the slope of the response against the preferred feature axis, confirmed that during state 1, tuning to the conscious face was significantly steeper than tuning to the suppressed face (\textbf{Fig. 4c}, top; $p = 8 \times 10^{-95}$, paired t-test); during state 2, the reverse held, with tuning to the suppressed face significantly steeper than to the conscious face (\textbf{Fig. 4c}, middle; $p = 6 \times 10^{-60}$, paired t-test).

Despite this apparent symmetry, the two states were not mirror images of each other. The gain of tuning to the conscious face during state 1 was higher than that to the suppressed face during state 2 (\textbf{Fig. 4c}, bottom; $p = 3 \times 10^{-40}$, paired t-test; mean ratio between conscious and suppressed tuning gain = 1.21). Across cells, tuning gain for the conscious face during state 1 was strongly correlated with gain for the suppressed face during state 2 (\textbf{Fig. 4d}, Pearson's $r = 0.90$), indicating that the same neurons encoded both faces. We also compared the mean responses across all stimuli during state 1 and state 2, and found that, unlike tuning gain, the \pdfpage{10}mean firing rate was weakly but significantly higher during state 2 for most neurons (\textbf{Fig. 4e}, $p = 2 \times 10^{-15}$, paired t-test). The two states therefore differ in two respects: a tuning-gain difference (steeper in the conscious state) and a small mean-rate offset (higher in state 2). A linear decoder distinguished state-1 from state-2 responses with 66\% accuracy (\textbf{Fig. 4f}, black, see Methods).

\subsection*{Temporal multiplexing explains weak average modulation during binocular rivalry}

We next revisited the original puzzle, why perceptual modulation during binocular rivalry appears weak when averaging across time (\textbf{Fig. 4g}, top left). One possibility is that during state 1, cells are strongly modulated by the conscious stimulus, and during state 2 cells are strongly modulated by the suppressed stimulus, and averaging across both states yields apparent weak modulation. To test this, we computed modulation indices separately for time points assigned to state 1 and state 2, inferred from a held-out half of neurons as in \textbf{Fig. 3}.

Strikingly, we found that modulation by conscious percept in the binocular condition during state 1 was approximately as strong as modulation for average unambiguous responses (\textbf{Fig. 4g}, top right, Pearson's $r = 0.66$). Modulation indices clustered around the diagonal (slope = 0.93), indicating that during state 1, cells responded as if the conscious stimulus was presented unambiguously. Conversely, during state 2, the pattern reversed, and modulation by conscious percept in the binocular condition was strongly anticorrelated with average unambiguous responses (\textbf{Fig. 4g}, bottom left, Pearson's $r = -0.44$). Modulation indices clustered around the antidiagonal (slope = $-0.76$), indicating that during state 2, cells responded as if the suppressed stimulus was presented unambiguously.

These results suggest that the weak average modulation was observed because averaging collapses responses across the two opposing states: state 1, which represents the conscious stimulus, and state 2, which represents the suppressed stimulus. However, this explanation raises a further question: \textit{why does any positive modulation remain after averaging across the two states?} In the previous section we found that the gain of tuning to conscious stimuli during state 1 was slightly higher than the gain to suppressed stimuli during state 2. We next asked whether this asymmetry could account for the residual positive correlation between binocular and unambiguous modulation indices.

To test this, we re-computed modulation indices for responses where we equalized the tuning gain for state 1 and state 2, and averaged the gain-equalized responses, weighted by the occupancy of state 1 and state 2 for each stimulus pair (\textbf{Fig. 4g}, bottom right; see Methods). Interestingly, after this correction the residual positive correlation disappeared (Pearson's $r = -0.02$, $p = 0.31$; slope = 0.07). Thus, weak average modulation during binocular rivalry can be explained by averaging across two temporally interleaved states with opposite stimulus selectivity, combined with slightly stronger tuning in state 1.

\subsection*{Dynamics of state switching}

We next asked how fast the population switches between the two states, and how long it dwells in each. Our earlier analysis labelled each time bin independently, ignoring any dependence \pdfpage{11}between successive bins. To ask whether the states persist over time, and if so, for how long, we turned to a Hidden Markov Model (HMM), fit separately for each stimulus pair. As in the earlier LLR analysis, the two states emitted spikes from fixed templates (the pair's priming-derived firing rates for the conscious and suppressed face), with the additional assumption that transitions between states were governed by a Markov process; only the $2 \times 2$ transition probability matrix was free, and it was estimated from the data (\textbf{Fig. 5a}). Fitting the model to that pair's binocular trials inferred the transition probabilities and the state sequence that maximized the probability of the observed spike counts, thereby partitioning each trial into two states over time (\textbf{Fig. 5b}). To determine the characteristic time scale of state alternations, we measured the median dwell times of the inferred states. Because dwell time cannot be resolved below the width of a single bin, it is expected to increase with bin size once the bin width becomes sufficiently large, and this dependence must be accounted for before the estimate can be interpreted. We therefore fit the HMM across a range of bin sizes and, at each bin size, asked how well the inferred state sequence explained the spikes of held-out neurons, relative to a one-state model (normalized conditional log-likelihood ratio; Methods). This measure peaked at small bin sizes, and over this range median dwell time plateaued around ${\sim}60$ ms (\textbf{Fig. 5c}), indicating that the estimate reflects the data rather than the temporal resolution of the analysis. Specifically, for the peak-likelihood bin size of 5 ms, median dwell times for state 1 and state 2 were very similar, 64 ms and 59 ms, respectively (\textbf{Fig. 5d}). This is consistent with the 67 ms time constant of likelihood autocorrelation reported earlier (\textbf{Fig. 2d}).

\subsection*{On and Off states}

Single-trial population activity showed strong fluctuations in overall firing, resembling previously described On and Off states\textsuperscript{16} (\textbf{Fig. 1g}). Details regarding these On-Off dynamics within face patches will be the focus of a separate forthcoming manuscript (C.M., J.K.H., F.L.L., D.Y.T., T.A.E., in preparation). Here, we address the specific concern that our two states might track this gain rhythm (high versus low excitability) rather than which stimulus is represented (conscious versus suppressed). To test this directly, we used a modified version of our log-likelihood ratio decoder that used two template vectors for each stimulus, one for the high-gain On periods and another for the low-gain Off periods. This four-state LLR decoder, like the original two-state LLR decoder, labelled each time point with the template vector that maximized the likelihood of the observed spikes. The On-Off axis absorbed the large excitability fluctuations, yet the conscious-suppressed axis remained independent of it: (i) conscious and suppressed states were equally represented within both gain regimes (49\%/51\% during On, 50\%/50\% during Off), (ii) the conscious/suppressed partition agreed with a partition by total firing rate (separating bins into higher/lower than median firing rate) no better than chance (Cohen's $\kappa \approx 0.0$), and (iii) the conscious/suppressed states recovered by the four-state model agreed closely with those from the original two-state model (\textbf{Fig. 2a}; 80\% agreement), leaving the state-resolved reconstructions and single-cell tuning asymmetry largely unchanged (\textbf{Extended Data Figure 9}). The On-Off fluctuations therefore did not substantially affect the inference of conscious/suppressed states.

\pdfpage{12}\subsection*{Alternations of conscious and suppressed states are weakly coordinated across face patches}

At the perceptual level, rivalry is unified: only one image is consciously seen at a time, and the corresponding representation is thought to be coherent across the brain. It is not clear, however, that this unity extends to the fast micro-epochs uncovered here. The alternation between conscious and suppressed states could be a coordinated, network-wide event, with regions switching together, or each region could alternate on its own clock, arriving at the same percept through independent selection. To distinguish these possibilities, we inferred the state sequence separately from the simultaneously recorded ML and AM populations, each with its own neurons and its own templates (\textbf{Fig. 2a}), and then predicted face features using neurons in the other patch (\textbf{Fig. 5e}). Explained variance of features was relatively low compared to \textbf{Fig. 3}, and for AM$\rightarrow$ML, explained variance for the conscious face in state 2 was higher than that for the suppressed face; combining across conditions, we did find a significant separation overall between conscious representations in state 1 and suppressed representations in state 2 (ML$\rightarrow$AM: $p = 3 \times 10^{-4}$; AM$\rightarrow$ML: $p = 9 \times 10^{-6}$, t-test, see Methods). However, our results do not provide strong evidence for coordination across patches in generating conscious and suppressed micro-epochs.

\subsection*{Comparison of two-state model to alternative models}

Several alternative mechanisms could in principle reproduce one or another signature of our analyses so far without reflecting a genuine, population-level temporal multiplexing of conscious and suppressed representations. We therefore considered a broad set of alternative models and subjected each, together with the recorded data, to a battery of complementary control analyses (\textbf{Extended Data Fig. 10a}, model $\times$ test matrix). Only a two-state model, in which the same neurons coherently alternate, within single trials, between faithfully representing the conscious and the suppressed face, reproduced every signature observed in the real data; each alternative failed at least three of the tests.

\textit{Alternative models}. We simulated four generative models of single-trial neural activity, each built from the same priming-derived, per-neuron response templates for the two faces of a trial ($r_A$ conscious, $r_B$ suppressed) and matched to the recorded trial structure of the real data (see Methods): (i) a two-state model, in which the entire population coherently alternates between the two faces at the ${\sim}60$ ms timescale estimated from the data; (ii) a conscious-only model in which only the conscious face is represented throughout the trial; (iii) an averaging model, in which every neuron fires at a fixed rate equal to the average of its two priming templates, with no alternation; and (iv) an uncoordinated model, in which every neuron alternates between the conscious and suppressed face with the same ${\sim}60$ ms dwell time as the two-state model but on its own, independent of the rest of the population. In every model, spike counts were generated as Poisson draws from the model's per-neuron, per-time-bin firing rate.

\textit{A battery of tests.} We evaluated every model, and the real data, on four tests (matrix rows in \textbf{Extended Data Fig. 10a}): (i) reconstruction of the conscious and suppressed faces from state-resolved activity using circular analysis (as in \textbf{Fig. 2h,i}); (ii) coordinated reconstruction, the same reconstruction after the state is inferred from one half of the population and the faces are \pdfpage{13}reconstructed from the held-out half (as in \textbf{Fig. 3}); (iii) reversal, the same tuning reversal measured in held-out neurons (as in \textbf{Fig. 4g}); (iv) model comparison of whether a two-state description of the activity fits held-out data better than a one-state alternative. In all cases, we used an HMM for state partitioning.

The conscious-only model cannot reconstruct the suppressed face for obvious reasons, and therefore is clearly incompatible with our experimental findings; we therefore focus on the three models that represent both faces. The two-state, averaging, and uncoordinated models could all reconstruct the conscious face in inferred state 1 and the suppressed face in inferred state 2 above chance using circular analysis (\textbf{Extended Data Fig. 10b}). This may at first seem surprising: how could a model with no genuine alternation, e.g., an averaging model that never represents either face on its own, be partitioned into two sets of time points from which noise alone reconstructs each pure face? Yet this is actually expected on deeper examination. In the averaging model every neuron's response sits at the mean of its two responses, $\frac{1}{2}(r_A + r_B)$, and only Poisson noise makes the activity vary; the HMM assigns the moments that happen to deviate toward $r_A$ to state 1 and those toward $r_B$ to state 2. Since the same data samples used for state inference are also used for computing state-conditioned responses, reconstructing from these selected moments then recovers $r_A$ and $r_B$ due to circular conditioning. Any model that contains both faces' responses---averaged together (averaging model), or switched independently (uncoordinated)---passes this test for the same reason. However, note that the averaging model, though seemingly simple, is actually quite unnatural: representing two competing stimuli as the average of their individual responses is a linear-superposition property that a strongly nonlinear system such as the ventral visual stream would not generally exhibit. Importantly, inferred states switched much more frequently for the averaging model data than for the real data, allowing for exacerbated circular conditioning; when matching dwell times (see Methods), real data yielded significantly better reconstructions than the averaging model ($\Delta = 8.2\%$ explained variance at ${\sim}80$ ms dwell, $p = 0.007$, bootstrap over n = 60 stimulus pairs, \textbf{Extended Data Fig. 10c}). Still, above-chance state-resolved reconstruction is expected even of models with no genuine alternation, and on its own is not diagnostic of a genuine, coordinated alternation.

The decisive question is whether these fluctuations are coordinated across neurons rather than private to each cell. When the state was inferred from one half of the population and the conscious versus suppressed face features reconstructed from the held-out half (\textbf{Extended Data Fig. 10d}), reconstruction held in the real data and the two-state model but collapsed for the averaging and uncoordinated models, whose apparent alternation arises from independent single-neuron fluctuations that are not shared across neurons.

The lower cross-validated explained variance in the real data compared with the idealized two-state simulation (\textbf{Extended Data Fig. 10b,d}) may partly reflect the fact that transitions in biological neural populations are unlikely to be instantaneous. While the simulated activity jumps directly between two discrete states, real population activity is expected to take a finite time to move between locations in neural state space, producing intermediate activity patterns. Forcing these transition periods into one of two states may reduce state-specific reconstruction accuracy and hence explained variance.

\pdfpage{14}Single-cell tuning showed the same pattern as above: When inferring the state from one half of the population and measuring tuning in the held-out half, only the real data and the two-state model retained a reversal, whereas the averaging and uncoordinated models collapsed to zero (\textbf{Extended Data Fig. 10e}).

As a complement to the preceding tests, we asked whether the population activity is statistically better described by two states than by one. For the real data and for each model, we fit two-state and single-state HMMs and compared their conditional likelihoods on held-out data (\textbf{Extended Data Fig. 10f}). For the single-state model the one rate was a weighted average of the two faces' priming rates; we scanned the averaging weight from 0 to 1 and kept the maximum-likelihood value. Two states fit the real data substantially better than one, establishing that the activity is genuinely organized around two population states; the coordinated two-state model showed the same advantage. For the averaging and uncoordinated models, by contrast, a single state fit better than two: adding a second state captured only noise and reduced the held-out likelihood.

\subsection*{Evidence for similar interleaved representations in prefrontal cortex}

The results so far show that inferotemporal cortex interleaves representations of the conscious and suppressed face. We next asked whether prefrontal cortex carries a pure representation of the conscious face, or an interleaved representation similar to inferotemporal cortex. We repeated the flash-suppression experiments in prefrontal cortex, recording from the lateral orbitofrontal sulcus and the infraprincipal dimple (\textbf{Fig. 1e, right}).

Reconstructions were of lower quality than in inferotemporal cortex (\textbf{Fig. 6a}): the best possible explained variance, reconstructing faces from priming responses, was below 80\%, compared to over 90\% for inferotemporal cortex, and particularly low (less than 30\%) for shape features (\textbf{Fig. 6b}). This may reflect a lower proportion of face cells or a different feature code in prefrontal cortex. Despite this, the same pattern of state-dependent representation of conscious and suppressed stimuli held in prefrontal cortex: reconstructions during state 1 more closely resembled the conscious face, and those during state 2 the suppressed face, confirmed by both explained variance and decoding using the same data samples as for state inference (\textbf{Fig. 6b-c}, $p = 10^{-14}$, t-test, see Methods). Importantly, as in inferotemporal cortex, the effect held under a cross-neuron validation: when the states were inferred from one half of the prefrontal population and the faces reconstructed from the held-out half, the state-dependent separation persisted (\textbf{Fig. 6d-e}, $p = 3 \times 10^{-4}$, t-test, see Methods), indicating a genuine population-level alternation.

Modulation by conscious percept was again much weaker compared to unambiguous alternation for average activity (\textbf{Fig. 6f}, $p = 4 \times 10^{-10}$, paired t-test; slope = 0.32), suggesting that, just like inferotemporal cortex, prefrontal cortex also does not have a pure representation of conscious perception for average responses. During state 1, however, modulation resembled the conscious stimulus (slope = 0.91) and during state 2 modulation resembled the suppressed stimulus (slope = $-0.52$). Temporal multiplexing of conscious and suppressed representations may therefore be a broad cortical phenomenon rather than a feature unique to inferotemporal cortex.

\pdfpage{15}\section*{Discussion}

The findings presented above reveal a temporally multiplexed code for conscious and suppressed visual representations. State-resolved reconstructions show that inferotemporal cortex alternates between complementary representational epochs: one in which the population represents the consciously perceived face, and another in which it represents the suppressed face. These representational switches in neural activity happen over an order of magnitude faster (15 Hz) than perceptual switches (${\sim}0.5$ Hz)\textsuperscript{6,17}. Single-cell analyses show that this alternation is not produced by separate populations of conscious- and suppressed-selective neurons, but by the same neurons changing which stimulus they represent across states. Importantly, during the conscious state, single-cell tuning to the conscious stimulus has stronger gain than tuning to the suppressed stimulus in the suppressed state, and a linear decoder can discriminate state 1 and state 2 across stimulus pairs. These states are weakly coordinated across simultaneously recorded face patches. Finally, prefrontal cortex shows a qualitatively similar interleaved structure. Together, these findings show that time provides a critical dimension for separating conscious and suppressed representations within cortical population activity.

Our findings expose a limitation of searching for the neural correlates of consciousness within the average activity of specific brain regions, whether toward the front or back of the brain. Rather than finding a pure representation of conscious perception in high-level areas, we found that average rivalry responses remained less modulated than responses to unambiguous stimulus alternation, even in prefrontal cortex; this is consistent with previous studies\textsuperscript{7}. The apparent weak modulation can now be understood as the consequence of averaging across temporally interleaved conscious and suppressed representations. The discovery of this temporal structure highlights the importance of large-scale simultaneous recordings for resolving coordinated population dynamics within single trials rather than relying on responses averaged across trials.

Previous studies noted that responses during rivalry did not fully collapse to the conscious stimulus response, leading the authors to suspect that suppressed stimuli might still leave residual information in neural activity\textsuperscript{7}. This possibility was later supported by the finding that suppressed stimuli could be decoded from inferotemporal responses\textsuperscript{2}. The present work extends these findings by showing that neural activity not only allows discriminating between two different suppressed stimuli but allows reconstruction of the suppressed stimulus with high fidelity. Thus, the suppressed stimulus is not represented as a weak residual trace but occupies the same high-dimensional face space as the conscious stimulus during complementary temporal epochs, and both epochs rely on a shared mapping between face features and single-cell responses. The presence of a suppressed stimulus representation even at high levels of the cortical hierarchy may help explain how unconscious stimuli are able to influence behaviour\textsuperscript{18}, attention\textsuperscript{19}, and affective responses\textsuperscript{20}.

The temporal multiplexing reported here may reflect a broader coding principle by which neural populations organize information over time. In attention, behavioural sensitivity and frontoparietal activity fluctuate rhythmically, suggesting that visual information is sampled in alternating temporal windows rather than through a continuously stable spotlight\textsuperscript{21-23}. In working memory, stimulus information can be expressed in brief beta- and gamma-band events, and recent work suggests \pdfpage{16}that frontal theta phase can modulate which remembered location is most accessible for readout, indicating that temporal structure can regulate when stored representations are expressed or accessed\textsuperscript{24-28}. In value-based decision-making, neural populations in orbitofrontal cortex alternate on a single-trial basis between states representing the values of different choice options\textsuperscript{29}. In spatial navigation, abrupt changes in context can produce theta-paced flickering between distinct hippocampal place-cell maps\textsuperscript{30}. Together with our findings, these studies suggest that temporal structure may be a general organizing principle for encoding multiple representations within overlapping neural populations.

The presence of rich representations of the suppressed stimulus raises a question for future work: why do representations in some micro-epochs reach awareness, whereas others do not? The fact that conscious and suppressed states can be discriminated by a linear decoder suggests that neural activity contains information not only about stimulus content, but also about whether the content represented in a given temporal epoch is destined for awareness or remains unconscious. One possibility is that this signature determines whether a sensory representation is dynamically routed to downstream systems involved in report, introspection, and behavioural control. Such routing could be controlled by recurrent interactions within visual cortex or by feedback from higher-order regions such as parietal cortex, the pulvinar, or other networks capable of coordinating perceptual states. Speculatively, suppressed epochs may allow the brain to explore alternative interpretations of the visual input, without being broadcast or influencing behaviour, whereas conscious epochs may express a working hypothesis that, through dynamic routing, becomes accessible to circuits that use it for introspection or actions such as report. Identifying the mechanisms that orchestrate these transitions, and determining why only some temporal segments become available to awareness, will be essential for understanding how temporal population dynamics give rise to conscious perception.




\section*{Figure legends}

\noindent\textbf{Figure 1. Large-scale recordings across cortical regions during a binocular flash suppression paradigm.} [Note: All human faces in this figure have been replaced by synthetically generated faces due to bioRxiv policy on displaying human faces; they are not photographs.] (\textbf{a}) Schematic of binocular flash suppression paradigm. First, a priming stimulus is presented monocularly in one eye. Then, without a blank period, a different stimulus is flashed to the other eye, and both are presented binocularly. Following the flash, the flashed stimulus is consciously perceived, while the priming stimulus is perceptually suppressed. (\textbf{b}) Perceptual reports in human participants verify perceptual suppression. Human participants were instructed to report whether they were seeing the stimulus in the left eye, right eye or mixed. The priming stimulus was presented at t=0 seconds in the left eye (left) and right eye (right), respectively. The dominant stimulus was flashed in the other eye at 2 seconds, causing the conscious percept to flip on most trials. Curves indicate average proportion of reported percepts of the stimulus in the left eye (red), right eye (blue), or a mixed percept (yellow). We compared the difference in report probability between the flashed and the primed stimulus, averaged over left eye and right eye primed trials, 300--1000 ms + report latency after flash onset with a two-tailed t-test against 0 ($t(13) = 9.4$, \pdfpage{17}$p = 4 \times 10^{-7}$, $n = 14$ participants, see Methods). (\textbf{c}) Schematic of recording setup. Monkeys viewed flash suppression stimuli on a screen through anaglyph goggles to separate input to the left and right eyes while two Neuropixels probes were inserted into different nodes of the visual cortical hierarchy. (\textbf{d}) CAD drawing of Neuropixels setup. Neuropixels probes were advanced using a custom insertion mechanism designed to target the desired regions. (\textbf{e}) Schematic of macaque brain with four regions targeted for electrophysiological recordings. (\textbf{f}) Response profiles of visually responsive cells (rows) to 96 visual stimuli (columns) from six categories, including faces and other objects. Shown are example sessions from ML (left), and AM (right). (\textbf{g}) Raster plot showing an example flash suppression trial. Black ticks represent spikes of 429 simultaneously recorded cells (y-axis) from ML and AM. Vertical bars and icons show when the priming stimulus and the rivalry stimulus were presented, respectively. (\textbf{h}) Scatter plot showing modulation indices for the binocular (y-axis) versus unambiguous (x-axis) conditions. Each dot corresponds to a cell and stimulus pair where the responses were significantly selective for the unambiguous stimulus. Modulation during the unambiguous condition was compared to modulation during the binocular condition using a two-sided paired t-test ($t(5549) = 25.2$, $p = 8 \times 10^{-133}$, $n = 5550$ cell-stimulus pairs).



\noindent\textbf{Figure 2. Partitioning neural activity into conscious and suppressed states.} (\textbf{a}) Schematic of log-likelihood ratio (LLR) state inference method. The LLR finds a partition of the time course into two states. (\textbf{b}) Example raster plot and state partition of a trial consisting of a priming stimulus, followed by binocular flash suppression. Orange shading indicates state 1, which corresponds to the flashed, conscious stimulus; purple shading indicates state 2, which corresponds to the suppressed priming stimulus. Vertical line indicates the time of the flash. Horizontal line demarcates 157 selective cells recorded from ML and 254 selective cells recorded from AM. (\textbf{c}) Proportion of time points inferred to be state 1 (orange) and state 2 (purple) during priming and binocular condition. Shaded regions indicate standard deviation across sessions. (\textbf{d}) Auto-correlation of the log-likelihood ratio for real data (black) and time-shuffled data (grey). Shaded regions indicate standard deviation across sessions. (\textbf{e}) Raster plot of neural activity of neurons from ML and AM, aligned to 22 transitions to state 1 (left) and 17 transitions to state 2 (right) during the binocular condition, for an example stimulus pair from the same session as (b). Cells are sorted by their modulation index of the priming responses to the two stimuli. Only cells that significantly preferred one of the stimuli were included. Cells above the solid horizontal line prefer the conscious face, and cells below prefer the suppressed face. Only state transitions where at least 20 ms of one state were immediately followed by at least 20 ms of the other were included. (\textbf{f}) Spike raster of an example neuron from ML that preferred the conscious face, from the same session and pair as (e). (\textbf{g}) Same as (f) for a simultaneously recorded example neuron from AM that preferred the suppressed face. (\textbf{h}) Example reconstructions of face images from neural activity during state 1 and state 2. The first and third columns show the ground truth conscious and suppressed faces presented to the monkey. The second and fourth columns show reconstructions obtained through linear regression of face features from neural activity during priming, and applied to neural activity in the binocular condition during state 1 and state 2, respectively. (\textbf{i}) Explained variance of conscious and suppressed face features. Explained variance (y-axis) was computed for the 20 principal shape features and 30 principal appearance features by comparing ground truth \pdfpage{18}features to features predicted by linear regression. Orange solid line: Conscious stimulus features from state 1. Orange dashed line: Suppressed stimulus features from state 1. Purple solid line: Conscious stimulus features from state 2. Purple dashed line: Suppressed stimulus features from state 2. Black solid line: Priming stimulus features predicted from average activity during priming. For each stimulus pair, we compared how much closer the reconstruction was to the conscious than the suppressed face between state 1 and state 2, and tested this difference against zero with a paired t-test ($t(29) = 48.5$, $p = 3 \times 10^{-29}$, $n = 30$ stimulus pairs, see Methods). (\textbf{j}) Decoding accuracy of conscious and suppressed faces during state 1 (orange) and state 2 (purple). Euclidean distance was computed between predicted face features and ground truth face features of the conscious/suppressed face and a given number of randomly selected distractor faces (x-axis). A predicted face was defined as decoded correctly when the distance was smaller to the conscious (solid) or suppressed (dashed) face compared to the distractors. The double dissociation of state-resolved decoding accuracy was tested per rivalry pair with a two-sided sign-flip permutation ($p = 10^{-5}$, $n = 30$ stimulus pairs). (\textbf{k}) Single-trial decoding for an example session. A linear multi-class decoder was trained on priming responses and applied to binocular rivalry responses to discriminate between 10 stimulus identities that were shown in a separate set of experiments. During binocular rivalry, within a 200-ms sliding window on each trial, the decoder used average activity across all time points in the window (top left), or across only the time points in the window assigned to state 1 (bottom left) or state 2 (bottom right), and the decoded identity was compared to the ground truth conscious (solid) or suppressed (dashed) stimulus. (\textbf{l}) Single-trial decoding summary across sessions. Activity was averaged across the entire binocular trial (grey), state 1 (orange), or state 2 (purple), and the decoded identity was compared to the ground truth conscious (solid) or suppressed (dashed) stimulus. Thin lines indicate individual sessions. Thick lines indicate average across the ten sessions. Chance performance is 10\%. The difference in decoding accuracy for conscious versus suppressed was evaluated using a \pdfpage{19}two-sided paired t-test (state 1: $t(9) = 11.0$, $p = 2 \times 10^{-6}$, $n = 10$ sessions; state 2: $t(9) = -13.7$, $p = 3 \times 10^{-7}$, $n = 10$ sessions).



\noindent\textbf{Figure 3. Coordinated switching of the population between conscious and suppressed states.} (\textbf{a}) Difference of marginal log-likelihoods of two-state and one-state models. We computed per-trial, per-neuron gain in log-likelihood of the two-state model over the best-fitting single-state model, whose maximum-likelihood static firing rate is given by the average firing rate during binocular rivalry. Boxes, median and interquartile range across the 19 sessions; whiskers, range; points, individual sessions; horizontal dashed line indicates where neither model has an advantage. We compared the log-likelihood difference against zero using a two-sided t-test ($t(18) = 4.84$, $p = 10^{-4}$, $n = 19$ sessions). (\textbf{b}) Same as Fig. 2e but for a held-out randomly selected half of neurons; the other half was used to infer states. (\textbf{c})-(\textbf{d}) Same example cells as in Fig. 2f, g but all other neurons were used for inferring states. (\textbf{e}) Reconstructions as in Fig. 2h but using held-out half of neurons. (\textbf{f}) Explained variance as in Fig. 2i but using held-out half of neurons (t-test, $t(29) = 17.5$, $p = 6 \times 10^{-17}$, $n = 30$ stimulus pairs). (\textbf{g}) Decoding accuracy as in Fig. 2j but using held-out half of neurons (two-sided sign-flip permutation test, $p = 10^{-5}$, $n = 30$ stimulus pairs). (\textbf{h}) Single-trial decoding for example session as in Fig. 2k but using held-out half of neurons. (\textbf{i}) Single-trial decoding summary as in Fig. 2l but using held-out half of neurons (two-sided paired t-test; state 1: $t(9) = 7.6$, $p = 3 \times 10^{-5}$, $n = 10$ sessions; state 2: $t(9) = -6.5$, $p = 1 \times 10^{-4}$, $n = 10$ sessions).



\noindent\textbf{Figure 4. Single cells switch between conscious and suppressed representations.} (\textbf{a}) Ramp-shaped tuning to features of conscious and suppressed faces during state 1 and state 2 for an example cell. Each dot represents one of the 1500 images, projected onto the cell's preferred axis (x-axis) and principal orthogonal axis (y-axis). The colour of each dot represents the cell's response to the image. Dot positions are determined by the features of the conscious face (left) or the suppressed face (right). Responses were computed during state 1 (top) or state 2 (bottom). The center inset shows marginal distributions for the conscious face in state 1 (orange) and suppressed face in state 2 (purple) as a function of the projection onto the preferred axis. Since red and cyan faces were presented through the anaglyph goggles, we removed the first principal component of features to remove this colour component for (a)-(f). The results without removing this component were qualitatively the same. (\textbf{b}) Population tuning across all cells to conscious face features (left) and suppressed face features (right) during state 1 (top) and during state 2 (bottom). For each cell (rows of each subpanel), colour indicates responses to each of the 1500 stimuli (columns), sorted by the projection onto the preferred axis of each cell (inferred by linear regression during priming). Cells are sorted by the slope of their response gradient, which we refer to as tuning gain. (\textbf{c}) Histogram of tuning gain of single cells for the features of the conscious face during state 1 (orange), features of the suppressed face during state 2 (purple), features of the suppressed face during state 1 (bright orange), and features of the conscious face during state 2 (bright purple). We compared gains for conscious and suppressed using a two-sided paired t-test (state 1: $t(542) = 25.5$, $p = 8 \times 10^{-95}$, $n = 543$ cells; state 2: $t(542) = -18.6$, $p = 6 \times 10^{-60}$, $n = 543$ cells). (\textbf{d}) Sharper tuning for conscious features during state 1. Each dot represents one cell. X-axis corresponds to gains for features of the conscious face during state 1. Y-axis corresponds to gains for features of the suppressed face during state 2. Dashed diagonal line indicates where the two gains are equal. We evaluated the difference between the gains with a paired t-test ($t(542) = 14.4$, $p = 3 \times 10^{-40}$, $n = 543$ cells) and computed the correlation between gains (Pearson's $r = 0.90$, $p = 5 \times 10^{-192}$, $n = 543$ cells). (\textbf{e}) Slightly stronger mean responses during state 2. Same conventions as (d), but for mean response across stimuli instead of gain. We evaluated the difference with a two-sided paired t-test ($t(542) = -8.0$, $p = 2 \times 10^{-15}$, $n = 543$ cells). (\textbf{f}) Decoding accuracy for discriminating state 1 vs state 2 across all 1500 stimuli. A linear decoder used the firing rate of all neurons. Shaded regions indicate standard deviation across 100 decoding iterations with randomly selected subset of cells and random 50/50 split of stimulus pairs into train and test set. Dashed line indicates 50\% chance level. (\textbf{g}) Unambiguous modulation vs perceptual modulation during average activity (top left; Pearson's $r = 0.45$, $p = 7 \times 10^{-147}$, $n = 2885$ cell-stimulus pairs), state 1 (top right; Pearson's $r = 0.66$, $p = 2 \times 10^{-355}$, $n = 2885$ cell-stimulus pairs), state 2 (bottom left; Pearson's $r = -0.44$, $p = 5 \times 10^{-139}$, $n = 2885$ cell-stimulus pairs), and activity averaged across gain-normalized state 1 and state 2 responses (bottom right; Pearson's $r = -0.02$, $p = 0.3$, $n = 2885$ cell-stimulus pairs; see Methods). As in Fig. 3, states were inferred using one half of neurons, and analysis was performed using the held-out half of neurons.



\noindent\textbf{\pdfpage{20}Figure 5. Dynamics of state alternations.} (\textbf{a}) Schematic of Hidden Markov Model (HMM). The HMM finds the most likely partition of the time course into two states given inferred state transition probabilities and the observed spike counts. (\textbf{b}) State partitions of all trials (rows), sorted by proportion of state 1, for an example session. (\textbf{c}) Fit accuracy (top) and median state dwell times (bottom) for state 1 (orange) and state 2 (purple) as a function of bin size parameter. Fit accuracy is quantified using a conditional log-likelihood ratio of a two-state HMM and a one-state HMM, normalized by its peak value across bin sizes per session. The emission vector in the one-state HMM was fixed at the mean of the two templates. Shaded regions indicate s.e.m. across sessions. (\textbf{d}) Fit accuracy (y-axis) as a function of median dwell time (x-axis, binned). Shaded regions indicate s.e.m. across sessions. (\textbf{e}) Explained variance of face-feature reconstruction when the two HMM states are inferred in one area and the features are reconstructed from the other. Left, states inferred from ML, reconstructed from AM; right, states inferred from AM, reconstructed from ML (paired t-test; ML$\rightarrow$AM: $t(29) = 4.1$, $p = 3 \times 10^{-4}$, $n = 30$ stimulus pairs; AM$\rightarrow$ML: $t(29) = 5.4$, $p = 9 \times 10^{-6}$, $n = 30$ stimulus pairs). Same conventions as Fig. 2i. Only 7 of the 9 sessions were included because cross-region analyses required simultaneous recordings from both ML and AM.



\noindent\textbf{Figure 6. Temporal multiplexing in prefrontal cortex.} (\textbf{a}) Reconstructions of faces during state 1 and state 2 for prefrontal cortex. Same conventions as Fig. 2h but LLR templates were computed and features predicted from neural activity in prefrontal patches PO and PVL. (\textbf{b}) Explained variance of features. Same as Fig. 2i but for PFC (t-test, $t(29) = 14.2$, $p = 10^{-14}$, $n = 30$ stimulus pairs). (\textbf{c}) Decoding accuracy. Same as Fig. 2j but for PFC (two-sided sign-flip permutation test, $p = 10^{-5}$, $n = 30$ stimulus pairs). (\textbf{d}) Explained variance of features using held-out half of neurons as in Fig. 3f, but for PFC (paired t-test, $t(29) = 4.1$, $p = 3 \times 10^{-4}$, $n = 30$ stimulus pairs). (\textbf{e}) Decoding accuracy, using held-out half of neurons as in Fig. 3g, but for PFC (two-sided sign-flip permutation test, $p = 0.03$, $n = 30$ stimulus pairs). (\textbf{f}) Unambiguous modulation vs perceptual modulation. Same as Fig. 4g but for PFC. We compared average perceptual modulation to unambiguous modulation using a two-sided paired t-test ($t(282) = -3.5$, $p = 5 \times 10^{-4}$, $n = 283$ cell-stimulus pairs). Correlation between unambiguous and perceptual modulation indices was quantified using Pearson's correlation ($r = 0.56$, $p = 3 \times 10^{-24}$; $n = 283$ cell-stimulus pairs for average, $r = 0.63$, $p = 1 \times 10^{-32}$; $n = 283$ cell-stimulus pairs for state 1, $r = -0.32$, $p = 5 \times 10^{-8}$; $n = 283$ cell-stimulus pairs for state 2, $r = 0.24$, $p = 0.005$; $n = 283$ cell-stimulus pairs for gain-equalized).


\section*{Extended Data Figures}


\noindent\textbf{Extended Data Figure 1. Recording targets.} MRI anatomical scans showing targeted face patches in inferotemporal cortex (\textbf{a}) ML and (\textbf{b}) AM, as well as in prefrontal cortex (\textbf{c}) PO and (\textbf{d}) PVL. Colour overlay shows fMRI contrasts for visual activation to faces versus objects, which was used to guide targeting.



\noindent\textbf{\pdfpage{21}Extended Data Figure 2. Weak reconstructions from average responses.} (\textbf{a}) Reconstructions as in Fig. 2h but using average activity rather than state-resolved activity. (\textbf{b}) Explained variance of predicting features of the conscious (solid) or suppressed (dashed) face as in Fig. 2i but using average activity (grey). The black dashed line shows the similarity between the face reconstructed from priming activity and the other face of the stimulus pair, i.e., the face that will subsequently be presented as the flashed stimulus, but was not yet present during the analyzed priming period. As the face has not appeared yet during the analyzed period, this is a measure for chance performance. (\textbf{c}) Euclidean-distance decoding based on predicted features as in Fig. 2j but using average activity.



\noindent\textbf{Extended Data Figure 3. Interleaved representations in individual monkeys.} (\textbf{a})-(\textbf{c}) Same as Fig. 3e-g but pooling only across sessions from monkey A. (\textbf{d})-(\textbf{f}) Same as Fig. 3e-g but pooling only across sessions from monkey M. As in Fig. 3, states were inferred from one half of neurons and features were predicted using the held-out half of neurons.



\noindent\textbf{Extended Data Figure 4. Interleaved representations in individual sessions.} Explained variance (mean of the shape and appearance explained variances, each image-variance-weighted across its features) for nine individual sessions. Boxes, median and interquartile range across the 9 sessions; whiskers, range; points, individual sessions. As in Fig. 3, states were inferred from one half of neurons and explained variance was computed using the held-out half of neurons.



\noindent\textbf{Extended Data Figure 5. Interleaved representations using leave-one-neuron-out cross-validation and high-confidence time points.} (\textbf{a})-(\textbf{c}) Leave-one-neuron-out cross-validation. Same conventions as Fig. 3e-g but for each neuron the state inference was computed using all other neurons, and then responses of each neuron were pooled to reconstruct faces. (\textbf{d})-(\textbf{f}) Reconstructions from high-confidence time points. Same conventions as Fig. 3e--g, but instead of assigning all time points to either state 1 or state 2, for each state the analysis was restricted to the 50\% of time points with the largest absolute-value log-likelihood ratio (i.e., the most confidently assigned time points).



\noindent\textbf{Extended Data Figure 6. Balance-trial control.} (\textbf{a})-(\textbf{e}) Same conventions as Fig. 3e-i but for subsampling time points so that for each trial there is the same number of state 1 and state 2 time points. As in Fig. 3, states were inferred from one half of neurons and reconstruction and decoding analyses performed on the held-out half.



\noindent\textbf{Extended Data Figure 7. Shuffle template control.} (\textbf{a})-(\textbf{c}): Same conventions as Fig. 2h-j but response templates for inferring states were randomly reassigned. For each presented stimulus pair (A, B), the templates were replaced by those of two randomly selected stimuli (A$'$, B$'$). The same stimulus reassignment was used across neurons, trials and sessions. Reconstructions were then compared to these newly assigned templates (A$'$, B$'$). (\textbf{d})-(\textbf{f}) Same as (a)-(c) but reconstructions were compared to the actually presented faces (A, B). Reconstructions for this control were obtained using circular reuse of the data for state estimation (double-dipping, as in Fig. 2) rather than cross-validation, so as to yield the highest-quality reconstructions and thus the most stringent control.



\noindent\textbf{\pdfpage{22}Extended Data Figure 8. Single-cell feature-tuning axes within each state align with the priming axis.} For each cell with priming-reconstruction explained variance $> 0.05$ the feature-tuning axis was estimated either from its priming responses (as in Fig. 4) or from its state-resolved binocular responses to either the conscious or suppressed face. Alignment between a state axis and the priming axis is quantified as the cosine similarity between the two weight vectors (intercept excluded); histograms show its distribution across 548 selected cells. Left, state 1; right, state 2. Within each panel, the axes estimated when the cell's preferred face was conscious and when it was suppressed are shown separately with different colours, see legend. We also computed the noise ceiling of the priming tuning axis itself, as the cosine similarity between priming axes estimated from two independent halves of the trials (black distribution). The axis computed from the conscious face during state 1 was only slightly less aligned with the priming axis than the noise ceiling (two-sided paired t-test, $t(547) = -2.5$, $p = 0.01$, $n = 548$ cells). The same was true for the state 2 suppressed face (two-sided paired t-test, $t(547) = -0.9$, $p = 0.35$, $n = 548$ cells). In comparison, the alignment was much weaker for state 1 suppressed compared to priming noise ceiling (two-sided paired t-test, $t(547) = -24.0$, $p = 10^{-87}$, $n = 548$ cells) and for state 2 conscious (two-sided paired t-test, $t(547) = -20.3$, $p = 6 \times 10^{-69}$, $n = 548$ cells). The first principal component of the face features, corresponding to stimulus colour, was removed, as in Fig. 4a-f.



\noindent\textbf{Extended Data Figure 9. On/Off states are separate from conscious/suppressed states.} (\textbf{a}) Example raster as in Fig. 1g. Orange and purple overlays correspond to conscious and suppressed states. Bright and dark overlays correspond to On and Off states. (\textbf{b})-(\textbf{d}) Same as Fig. 3e-g but for a four-state LLR that included On/Off x conscious/suppressed states. State 1 = conscious On+Off states, State 2 = suppressed On+Off states. As in Fig. 3, states were inferred from one half of neurons and reconstructions obtained using the held-out half. (\textbf{e})-(\textbf{h}) Same as Fig. 4c-g but for four-state LLR.



\noindent\textbf{Extended Data Figure 10. Comparison of the two-state model to alternative models.} Across all panels, models are abbreviated: Real, recorded data; 2-state, two-state; C-only, conscious-only; Average, averaging; Uncoord, uncoordinated. (\textbf{a}) Model $\times$ test matrix. Columns, the recorded data and the four generative models; rows, the four tests. Green check, pass; red cross, fail. Pass criteria: Reconstruction/\allowbreak Reconstruction\textsuperscript{coord}, in state 1, the explained variance of the conscious face exceeds both zero and that of the suppressed face; in state 2, the explained variance of the suppressed face exceeds both zero and that of the conscious face (bootstrap for each comparison, $p < 0.05$). The superscript ``coord'' denotes the cross-validated version of reconstruction, in which the state is inferred from one half of the neurons and reconstructions are evaluated in the held-out half. Reversal, state 2 held-out tuning correlation significantly negative (bootstrap, $p < 0.05$). Log-likelihood, held-out two-state log-likelihood exceeding one-state. (\textbf{b}) Reconstruction accuracy. Explained variance for the conscious and suppressed face features, reconstructed from state-resolved activity (mean of the shape and appearance explained variances, each image-variance-weighted across its features), for the recorded data and each model. Bars, mean; error bars, bootstrap standard deviation over stimulus pairs (2000 samples). (\textbf{c}) Conscious-face reconstruction explained variance in state 1 as a function of the achieved mean state dwell time, for the recorded data (red) and the averaging model (blue). Each point \pdfpage{23}corresponds to one imposed HMM self-transition probability $p$; $p$ was varied from 0.5 to 0.99 and for each $p$ the empirical dwell time was measured for both real and synthetic data (x-axis). (\textbf{d}) Cross-validated reconstruction accuracy. As in b, but with the state inferred from one half of the neurons and the face features reconstructed from the held-out half. (\textbf{e}) Held-out tuning reversal. State-1 (orange) and state-2 (purple) single-cell tuning correlation $r$ for the recorded data and each model. States were inferred from one half of the neurons and tuning measured in the held-out half. Bars, mean; error bars, standard deviation of 2000 bootstrap resamples across cell-stimulus pairs. (\textbf{f}) Model comparison. Difference in cross-neuron validated conditional log-likelihood between a two-state and a one-state Hidden Markov model (per cell) for the recorded data and each model; values above zero indicate that two states fit better. Error bars, s.e.m. across sessions.



\section*{Methods}

All animal procedures in this study complied with local and National Institutes of Health guidelines including the US National Institutes of Health Guide for the Care and Use of Laboratory Animals and were performed with the approval of the California Institute of Technology Institutional Animal Care and Use Committee (IACUC) and the University of California Berkeley Animal Care and Use Committee (ACUC). Human participant experiments were performed with approval of the California Institute of Technology Institutional Review Board (IRB).

\paragraph{Targeting.} Two male rhesus macaques were implanted with head posts and acrylic implants. We targeted face patches ML and AM in IT cortex, as well as ventrolateral prefrontal and orbitofrontal cortex for electrophysiological recordings. ML and AM were identified using fMRI. Monkeys were scanned in a 3T scanner (Siemens), as described previously\textsuperscript{31}. During functional scans, MION contrast agent was injected to increase signal-to-noise ratio and monkeys passively viewed blocks of faces and blocks of other objects to identify face-selective patches in the brain, using a custom eight-channel coil (MGH) at isotropic 1 mm resolution. Recording chambers were implanted over ML and AM, and custom grids and lattices were 3D printed (3D Systems) and inserted into the chamber to target the desired areas. Both chamber placement and grid design were planned with the software Planner\textsuperscript{32}. We inserted tungsten electrodes and then performed anatomical MRI scans to confirm correct targeting of the desired location.

\paragraph{Electrophysiological recording.} Single tungsten electrodes were used for scout recordings in the targeted areas prior to Neuropixels recordings. An Omniplex DHP system (Plexon) and custom analysis software written in MATLAB (MathWorks) were used for these initial recordings. Tungsten electrodes with an impedance of 100 k$\Omega$ (FHC) were inserted through plastic guide tubes (WPI), that extended 2 mm past the dura. Guide tubes were inserted through a custom 3D printed grid with holes designed to be in the direction of the target. After tungsten electrode recordings confirmed the presence of cells at the desired target, custom 3D-printed lattices were designed for Neuropixels recordings using the same trajectory. Subsequently, Neuropixels 1.0 NHP probes (45 mm long probes with 4416 contacts along the shaft, of which 384 are selectable at any time) were inserted using a custom-machined insertion mechanism, as part of which either a MDO-974A microdrive (Narishige) or an ultrasonic piezo motor drive (XLA-1 series, Xeryon) \pdfpage{24}slowly advanced the Neuropixels probe through a stainless steel cannula (A-M systems) along a linear ball bearing rolling guide rail (IKO). Our insertion mechanism is described in more detail in 13. All Neuropixels data were acquired using SpikeGLX (Bill Karsh, HHMI/Janelia Research Campus) and spike sorted using Kilosort3\textsuperscript{33}. Both single-units and multi-units were used for analyses.
\decompnote{The citation number 13 in the sentence above is printed as plain text in the PDF, not as a superscript; kept as printed.}

\paragraph{Stimuli and behavioural task.} Monkeys were head-fixed and viewed an LCD screen (Asus ROG Swift PG43UQ) of 60 visual degree width and 44 degree height in a dark room. Stimuli were displayed on the screen with a 7 visual degree size while red/cyan anaglyph goggles were placed in front of the left and right eye of the monkey. Anaglyph goggles were custom made with filters (LEE filters) to match the red and green/blue emission spectrum of the screen, respectively, so that inputs to left and right eyes could be controlled independently. Emission spectra were measured using a PR-650 SpectraScan colorimeter (Photo Research). Gaze position was monitored using an infrared camera eye tracking system (Eyelink) sampled at 1000 Hz. The camera recorded one or both eyes through the red/cyan anaglyph filter. In the first phase of the experiment, each monkey passively viewed two sets of 96 and 49 screening stimuli, in pseudorandom order (150 ms ON time, 150 ms OFF time) with a fixation spot of 0.2$^\circ$ diameter in the center of the screen. Screening stimuli consisted of 16 images of faces and 80 non-face images in the first screening set, with ${\sim}8$ repeats per image, and 40 images of faces and 9 non-face images in the second screening set, with ${\sim}30$ repeats per image. Monkeys received a juice reward for maintaining fixation for at least 3 s within a 0.5$^\circ$ rectangle centered on the fixation spot. Stimulus progression was contingent on fixation. For all experiments, only trials where monkeys maintained fixation for at least 90\% of the trial duration were included for analyses.

After the screening phase, one of two types of flash suppression stimulus sets was presented. In flash suppression, a priming stimulus is presented, followed immediately, without a blank period, by a binocular stimulus or another type of control stimulus (\textbf{Fig. 1a}). For the first type of stimulus set, ten face/object stimuli (five stimuli for left eye, and five stimuli for right eye) were selected from the 49 stimuli screening set that maximized the mean pairwise Euclidean distance between stimuli in neural space. The 25 combinations of the five left eye stimuli and five right eye stimuli were presented as stimulus pairs. For each stimulus pair, one of the two stimuli was presented monocularly as a priming stimulus for two seconds. The priming stimulus was followed by one of the two types of test stimuli, presented for four seconds: (1) a binocular stimulus, consisting of the priming stimulus and the other stimulus of the stimulus pair presented to the two eyes, respectively, referred to as the `binocular' context, (2) a monocular stimulus, consisting of only the other stimulus of the pair that was not presented during priming, referred to as the `unambiguous' context, used for the modulation on the x-axis of \textbf{Fig. 1h} and \textbf{Fig. 4g}. For each stimulus pair, both contexts were presented, and either stimulus was presented as priming stimulus in different trials. The new stimulus that was flashed during the post-priming test stimulus is referred to as the dominant stimulus, whereas the stimulus that was already present during priming is referred to as the suppressed stimulus. The test stimulus was followed by a 500 ms blank period containing only the fixation spot on a grey screen. Four stimulus pairs were presented with ${\sim}30$ trials per condition, and the rest of the 25 stimulus pairs were presented with ${\sim}5$ trials \pdfpage{25}per condition (for image pair $\times$ which image was dominant $\times$ unambiguous / binocular = 100 conditions). This experiment was performed in seven sessions in monkey A and three sessions in monkey M and analysed for \textbf{Fig. 1h, Fig. 2c,d,k,l, Fig. 3a,h,i, Fig. 4g, Fig. 5c,d} and \textbf{Fig. 6f}.

For the second type of flash suppression stimulus set, 1500 unique faces, randomly sampled from face feature space, were presented. The same 1500 faces were presented in all experiments using this second type of stimulus. To maximize the number of unique stimuli, we did not present all combinations; instead, this set consisted of 750 stimulus pairs, with each of the 1500 faces being part of only one stimulus pair. Moreover, the priming stimulus duration and post-priming test stimulus duration were both reduced to one second, and the blank period after the test stimulus was reduced to 250 ms. As in the first experiment, both stimuli of the stimulus pair were presented as priming stimulus in different trials. This type of experiment did not include an unambiguous alternation condition. A subset of 60 faces was presented for a higher number of repeats (${\sim}20$ trials per condition for high repeat stimulus pairs, ${\sim}2$ trials per condition for the rest of the 750 stimulus pairs). Reconstructions were generated for this subset of 60 faces, while the linear regression was trained on all 1500 faces. This experiment was performed in a separate set of five sessions in monkey A and four sessions in monkey M, and analysed for \textbf{Fig. 1g, Fig. 2b-j, Fig. 3a-g, Fig. 5} and \textbf{Fig. 6a-e}.

\paragraph{Data analysis.} For fitting LLR and HMM, Python version 3.11 was used. All other analyses were performed using MATLAB 2023b and 2025b (MathWorks). For \textbf{Extended Data Fig. 5} and \textbf{Extended Data Fig. 10} Claude Code was used in order to reduce computational runtime, and all generated code was reviewed for correctness by J.K.H.

\textit{Human reports.} Because button reports lag stimulus presentation, we first estimated report latency from the priming period, independently of the flash: for each observer, the time at which the probability of reporting the primed colour reached half its rise from the pre-onset baseline to its maximum (0.86 $\pm$ 0.05 s, mean latency $\pm$ s.e.m., \textbf{Fig. 1b}). The analysis window for determining significance of suppression was 300--1000 ms after flash onset (the same window as used for neural analysis, see below), shifted by this latency. Flash suppression was quantified per observer as a suppression index: the mean difference in report probability between the flashed and the priming stimulus within this window. We performed a two-sided t-test of this suppression index against the null hypothesis of it being 0 across all 14 human participants. The flashed percept was reported on 82\% of binocular rivalry trials, only 5 percentage points below the 87\% correct reporting of the unambiguous priming stimulus. The 13-percentage-point deviation from ceiling during priming provides an estimate of baseline report errors human participants made even when the percept was unambiguous. The small additional reduction during rivalry may be partially due to a refractory period in reporting rather than failures of perceptual suppression: rivalry reports followed shortly after the preceding priming report, whereas the priming report was made after a substantially longer response-free interval. Closely spaced responses are known to interfere with one another, producing a psychological refractory period or dual-task cost that reduces accuracy of the second response\textsuperscript{34}.

\textit{\pdfpage{26}Log-likelihood ratio (LLR).} For each binocular trial we computed a continuous single-trial index of the momentary state as the population log-likelihood ratio (LLR) between the two states of a Poisson mixture model with equal mixture weights (\textbf{Fig. 2a}). In consecutive 10-ms bins, the spike counts of the visually responsive IT population were evaluated under the two emission templates, which were computed from the mean firing rates of each neuron during the corresponding single-face priming trials. The LLR was computed as the difference of the population Poisson log-likelihoods,
\begin{equation*}
\ell(t) = \sum_i n_i(t) \cdot \log\!\left(\lambda_i^{c} / \lambda_i^{s}\right) - \left(\lambda_i^{c} - \lambda_i^{s}\right) \cdot \Delta t ,
\end{equation*}
where $n_i(t)$ is the spike count of neuron $i$, $\lambda_i^{c}$ and $\lambda_i^{s}$ its template firing rates derived from priming activity when the conscious/suppressed face was presented unambiguously, and $\Delta t = 10$ ms. $\ell(t)$ is positive when the population evidence favours the conscious state and negative when it favours the suppressed state, and states were assigned based on this sign. To determine whether rivalry responses were better described by two states than by a single static response, we compared the marginal likelihood under the two-state Poisson mixture model (used for LLR decoding) with that under the best-fitting one-state model. The best-fitting one-state model was obtained by maximum-likelihood estimation, yielding a static firing-rate vector equal to the average response during binocular rivalry trials (\textbf{Fig. 3a}). For the two-state model, the two states were assigned equal prior probability and the marginal likelihood of the spike-count vector in each time bin was computed by averaging its Poisson likelihood under each of the two priming templates, thereby marginalizing over state assignments. The likelihood of a trial was the product of these marginal likelihoods across time bins. To compute the likelihood of the static model, we split the rivalry trials into halves, computed the average rivalry response from one half, and then calculated the Poisson likelihood of the spikes in the other half using this estimate as the underlying firing-rate vector. We then repeated this procedure with the two halves reversed. The two-state templates were estimated from independent priming trials and therefore did not require a trial split.

\textit{Autocorrelation and LLR timescale.} For each session, we computed the autocorrelation of the LLR at 1 ms lag resolution up to 500 ms, pooling over binocular trials. We then averaged this curve across 19 sessions (\textbf{Fig. 2d}). As a null, we also recomputed the autocorrelation of the LLR after permuting time bins of neural data within trials, independently for each neuron, which abolished all temporal structure. Both real and shuffled autocorrelation curves were normalized by their zero-lag value. The real autocorrelation has two regimes: below 10 ms it reflects the triangular autocorrelation of the 10-ms analysis window ($\ell$ is constant within a bin) and carries no dynamical information, whereas beyond 10 ms it decays with the persistence of the LLR signal. The shuffled autocorrelation goes to zero after the 10 ms window, as expected. Modelling the state as a symmetric two-state Markov chain, whose autocorrelation is $\rho \cdot e^{-2\tau/D}$ with mean dwell time $D$, we fit this form to the session-averaged autocorrelation over lags 10--90 ms (from the bin edge to the first zero-crossing) by least squares, with the amplitude $\rho$ free. This mean dwell time of $D = 67$ ms corresponds to a median dwell of $D \cdot \ln 2 \approx 46$ ms for exponentially distributed dwell times, approximately matching the HMM-inferred time scale described below.

\textit{Reconstructions.} Face features were predicted from population activity by linear regression, $F = W \times R$, where $F$ is the vector of 50 face features, $R$ is the population response projected into principal component (PC) space, and $W$ is a matrix of regression weights. Responses of 3501 \pdfpage{27}visually responsive neurons across sessions were computed from spike counts averaged across trials and across time points from 300 ms to 1000 ms after stimulus onset and baseline-normalized using the spike counts from 50 ms after blank onset until the end of the blank period that followed each binocular flash suppression stimulus. Visual responsiveness of cells was independently determined based on whether responses from 70 to 200 ms after stimulus onset were significantly different from baseline firing 0 ms to 70 ms after stimulus onset during the first stimulus set with 96 stimuli (t-test, $p < 0.05$). Because the number of neurons greatly exceeded the number of stimuli available for fitting, regressing features directly onto single-neuron responses would have overfitted the data. To prevent overfitting, we first projected neural responses into a 500-dimensional space using principal component analysis (PCA) before performing linear regression (MATLAB lscov, ordinary least squares, no regularization); 300 components were used when reconstructing from the 661 prefrontal neurons, reflecting the smaller population. Coefficients for PCA were computed from priming responses, and subsequently these coefficients were used to project all responses into the same space. Linear regression weights were also computed using priming responses, and then those same weights were used when predicting face features in the binocular context. During the binocular context, only time points during the same time window 300 to 1000 ms that were assigned to state 1, or state 2, respectively, were used to compute responses, whereas for priming responses all time points from 300 ms to 1000 ms were averaged. For both priming and state-resolved binocular rivalry, predictions were cross-validated by excluding the predicted face from the training set. After averaging state-dependent responses for each neuron, responses were concatenated across the neuron dimension across sessions. Reconstructions of face images were computed from face features using the active appearance model used in \textsuperscript{35}, rendered in greyscale and contrast-normalized for visualization purposes (\textbf{Fig. 2h}).

\textit{Explained variance.} For the explained variance of a given feature, we computed one minus the ratio of the sum of squared residuals and the sum of squared mean-subtracted features (\textbf{Fig. 2i}):
\begin{equation*}
R^2 = 1 - \frac{\sum_i (y_i - \hat{y}_i)^2}{\sum_i (y_i - \bar{y})^2} ,
\end{equation*}
where $y_i$ is the actual feature of face $i$, $\hat{y}_i$ is the predicted feature of face $i$, and $\bar{y}$ is the mean of the feature across faces. To test whether the two states encode the conscious and suppressed faces separately, we reconstructed the face features in each state and, for each stimulus, computed a leaning index: the distance (in feature space) from the state's reconstruction to the suppressed face minus its distance to the conscious face (positive = closer to the conscious face; conscious = reported/dominant face, suppressed = its pair partner). We then took the difference between states, leaning(state 1) $-$ leaning(state 2), which is positive when state 1 leans toward the conscious face and state 2 toward the suppressed face. This difference was averaged over each pair's two dominance conditions and tested against zero with a one-sample t-test across high-repeat stimulus pairs. To aggregate the explained variance of individual face-model features into a single metric, we computed the explained variance separately for the shape features and the appearance features and averaged the two. Within each set, features were weighted by the variance they contribute to the face image, namely the variance in landmark coordinates for shape, and in the standard-shape pixel intensities for appearance.

\textit{\pdfpage{28}Euclidean decoding.} Euclidean-distance-based decoding was performed on the predicted face features (\textbf{Fig. 2j}). We picked the 50 principal z-scored shape and appearance features that were predicted for a given face, and then compared the Euclidean distance of the selected features to the ground truth features of the selected face and the ground truth features of a given number of randomly chosen distractor faces. The selected face was labelled as decoded correctly if the distance of the predicted features was closer to its ground truth features than features of all the other faces. This was repeated for all 60 predicted faces to yield a decoding accuracy. To test for statistical significance of the Euclidean-distance-based decoding we computed for each high-repeat face identity the 10-way Euclidean nearest-neighbour decoding accuracy of the state-resolved reconstruction against the true face-feature vectors (averaged over random 10-identity subsets), separately for state-1 (dominant/conscious) and state-2 (suppressed) responses and for the dominant and the suppressed face. Per identity we formed the interaction contrast, $(\mathrm{dominant}_{\mathrm{state\,1}} - \mathrm{suppressed}_{\mathrm{state\,1}}) - (\mathrm{dominant}_{\mathrm{state\,2}} - \mathrm{suppressed}_{\mathrm{state\,2}})$, averaged the two members of each rivalry pair into a single value (n = 30 pairs, the unit of inference), and tested it against zero with a two-sided sign-flip permutation test ($10^5$ permutations, each randomly flipping the sign of whole pairs).

\textit{Single-trial decoding.} A cross-validated multi-class classifier was trained to read out face identity from the population response and applied to individual binocular trials. The classifier (a linear support-vector machine, 8-fold stratified cross-validation) was trained on the priming trials, using as features each neuron's mean firing rate over the priming response window; class labels were the identities of the 10 presented faces. The priming-trained classifier was then applied, without retraining, to binocular trials, in sliding 200 ms windows stepped by 10 ms across the trial (\textbf{Fig. 2k}) or over the full stimulus window (\textbf{Fig. 2l}). For every trial and window, the population response was averaged over the window and classified into a single predicted face. We computed conscious decoding accuracy, the fraction of trials on which the predicted face matched the consciously perceived face, and suppressed decoding accuracy, the fraction matching the suppressed face. We repeated the readout on state-restricted responses: within each window the population response was averaged only over time points assigned to a given state (state 1 or state 2), and trials with no time points in that state in the window were dropped.

\textit{Axis tuning.} For each neuron we defined its preferred axis in face-feature space from a cross-validated (leave-one-face-out) linear encoding model fit to average priming responses, predicting the neuron's response to each face from face features; the fitted weight vector is the neuron's preferred axis, and the predicted response is each face's projection onto it. For the example cell (\textbf{Fig. 4a}), tuning was visualized in two feature dimensions: the preferred axis (x, cross-validated predicted response) and an orthogonal axis (y, first principal component of the features after orthogonalizing to the preferred axis; the PCA was computed on stimuli with the same colour so that it did not collapse onto the colour axis). Each point corresponds to one face image, positioned by its projection onto the two axes and coloured by the neuron's measured response (red: high response, blue: low response) during the binocular condition in a given state/percept; accompanying marginal plots show the neuron's binned response along the preferred axis for state 1/conscious (solid) versus state 2/suppressed (dashed). \pdfpage{29}For population tuning maps (\textbf{Fig. 4b}), each neuron's state-resolved binocular response was plotted across its axis-ranked stimuli as a neurons $\times$ stimuli image, separately for the four conditions --- state 1 / conscious face, state 2 / conscious face, state 1 / suppressed face, state 2 / suppressed face. We only included neurons whose priming responses could be explained by features (cross-validated $R^2 > 0.05$). Axis tuning gain was quantified per neuron as the slope of an ordinary-least-squares regression of the measured response on stimulus position along the preferred axis (intercept + slope), computed separately in each state and percept.

\textit{Decoding state 1 vs state 2.} To test whether IT population activity allows distinguishing the two inferred states, generalizable beyond stimulus-specific tuning, we trained a linear decoder to discriminate the ordering of paired state 1 and state 2 responses (\textbf{Fig. 4f}). For each stimulus we took the population response vectors to that stimulus being conscious during state 1 and to that stimulus being suppressed during state 2. We then constructed a balanced two-class decoding problem: for every stimulus, we created two types of population vector concatenations: the state-1 response, followed by the state-2 responses, was labelled class 1, and the reversed concatenation, state 2, followed by state 1, was labelled class 2. Because each stimulus contributes identical responses to both classes and differs only in concatenation order, stimulus identity carries no information about the label; the only cue is whether the state-1 response is the first half or second half of the vector. The decoder can therefore only discriminate labels correctly if the state-1 and state-2 responses to the same stimulus are systematically distinguishable, beyond stimulus identity. We trained a linear SVM on a random half of the stimulus pairs and tested on the held-out half with both members of a pair kept in the same set, so that no stimulus appeared in both training and test. Accuracy was computed as a function of number of neurons by subsampling neurons, over 100 random split-and-subsample iterations per size.

\textit{Modulation.} Modulation indices were computed using the average response for each cell and stimulus pair from 300 ms to 1000 ms after stimulus onset, where human participants reported only perceiving the flashed stimulus for most trials after accounting for report delay. We chose this time window conservatively: 300 ms avoids visual transients caused by the flash, as visual response latencies are around 100 ms in face patches. The end of the time window, 1000 ms, was determined by the stimulus duration in the main experiment. We computed modulation indices using the formula $MI = \frac{r_A - r_B}{r_A + r_B}$, where $r_A$ and $r_B$ are the responses to stimulus A and B, respectively. For the x-axis, the responses to unambiguous stimuli (A, B) of a given pair were used, while for the y-axis, the responses to binocular stimuli were used, when stimulus A was dominant for $r_A$ and when stimulus B was dominant for $r_B$. For \textbf{Fig. 4g}, the responses for modulation indices on the y-axis were computed over time points from 300 ms to 1000 ms that were assigned to state 1 or state 2, respectively, whereas modulation indices on the x-axis were computed across all timepoints. We inferred states using one half of neurons and computed modulation indices only for the held-out set of neurons. Only cells that (1) were significantly selective to A vs B during priming responses (t-test over responses 300 ms to 1000 ms after priming onset, with threshold $p < 0.05$) and (2) had at least one spike for both stimuli of the pair and both states were included. These selection criteria included a total of 2885 cell-stimulus pairs. For gain-equalization, we computed tuning gain for state 1 conscious and state 2 suppressed, \pdfpage{30}and modified responses during state 1 to equalize the gain:
\begin{equation*}
\hat{r}^{s1}_i = \left(r^{s1}_i - \bar{r}^{s1}_i\right) \times \frac{\mathrm{gain}^{s2}_i}{\mathrm{gain}^{s1}_i} + \bar{r}^{s1}_i ,
\end{equation*}
where $\hat{r}^{s1}_i$ is the gain-normalized response of neuron $i$ to the conscious face in state 1, $\mathrm{gain}^{s1}_i$ and $\mathrm{gain}^{s2}_i$ are the tuning gain of neuron $i$ to the conscious face in state 1 and suppressed face in state 2, respectively, and $\bar{r}^{s1}_i$ is the mean response of neuron $i$ during state 1, which is kept constant during the gain equalization. Neurons whose two gains had opposite sign, for which no positive rescaling exists, were excluded from this control. The gain-equalized average response to binocular rivalry was then computed as $\hat{r}^{\mathrm{binoc\_avg}}_i = p^{s1}\,\hat{r}^{s1}_i + p^{s2}\,r^{s2}_i$ where $r^{s2}_i$ is the non-gain-equalized response to the suppressed face during state 2, and $p^{s1}$ and $p^{s2}$ are the proportions of inferred states 1 and 2 for the given stimulus pair. The relationship between binocular and unambiguous face selectivity was quantified as the slope of each neuron's binocular modulation index against its unambiguous modulation index, estimated by total least-squares regression. Because both modulation indices are noisy single-neuron estimates, we used total least squares, which is symmetric in the two variables rather than ordinary least squares which would be attenuated toward zero by measurement noise.

\textit{Hidden Markov Model (HMM).} HMMs with two latent states were fit to neural recordings to partition each trial into segments representing the conscious and suppressed stimulus, respectively (\textbf{Fig. 5a-d}). For each stimulus pair (A, B), average stimulus response vectors $r^A$ and $r^B$ were computed from priming trials so that $r_i^A$ is equal to the average firing rate of unit $i$ across all priming trials where stimulus A was shown. To account for stimulus onset time, only time bins 70-1000 ms after stimulus onset were included in this average. These average response vectors were then used as fixed firing rates in a two-state Poisson HMM for identifying micro-epochs in this stimulus pair. The hidden state $z_t \in \{1, 2\}$ for this model follows a Markov chain with transition probabilities determined by a $2 \times 2$ matrix $T$: $p(z_{t+1} = j \mid z_t = i) = T_{ij}$, where $p(z_0 = 1) = p(z_0 = 2) = 0.5$. The vector of unit spike counts $y_t \in \mathbb{R}^N$ for time $t$ with a bin size of $\Delta t$ is assumed to follow a Poisson distribution for each unit with mean determined by the state $z_t$:
\begin{equation*}
p(y_t \mid z_t) = \prod_{i=1}^{N} \frac{1}{y_i!} \left(\lambda_i(z_t)\,\Delta t\right)^{y_i} e^{-\lambda_i(z_t)\,\Delta t} ,
\qquad \text{where } \lambda_i(z_t) = \begin{cases} r_i^A & \text{if } z_t = 1 \\ r_i^B & \text{if } z_t = 2 \end{cases}
\end{equation*}
The transition matrix $T$ is the only free parameter in this HMM, and was fit by maximizing the likelihood of binned spike counts observed during binocular trials for each stimulus pair, using the Baum-Welch algorithm. After the transition matrix was estimated, the Viterbi algorithm was used to compute the most likely state sequence $\{z_{k1}, z_{k2}, \ldots, z_{kT}\}$ for each binocular trial $k$ for all stimulus pairs. The HMM states $\{z_{kt}\}$ corresponding to the stimuli were converted into states $\{c_{kt}\}$ indicating whether the dominant or suppressed stimulus was being represented at time $t$:
\begin{equation*}
c_{kt} = \begin{cases} z_{kt} & \text{if stimulus A is dominant in trial } k \\ 3 - z_{kt} & \text{if stimulus B is dominant in trial } k \end{cases}
\end{equation*}
To compute conditional likelihood, we split units into two random halves. We used the first half to infer states with the HMM, and then calculated the likelihood of observing the spikes from the second half of units given the state assignments inferred from the first half. To compare to a single-state model, we then calculated the likelihood of observing the same spikes from the second half under a static rate, namely, the average of the priming templates. For each session, \pdfpage{31}we took the trial-averaged difference between the HMM model and the static model and normalized it by the number of units used to compute it. All code for HMM fitting and prediction was written in Python using the \texttt{hmmlearn} package.

\textit{Synthetic models.} We simulated four generative models of the single-trial population activity, each constructed per session from the priming-derived response templates of the two faces of a trial ($r_{\mathrm{conscious}}$, $r_{\mathrm{suppressed}}$) and matched to the recorded number of trials and time bins of each stimulus pair (\textbf{Extended Data Fig. 10}). In every model, spike counts in each bin were independent Poisson draws based on the mean firing rate for that bin. The only difference between models was how that mean firing rate was set for each time bin and each neuron. Two-state (temporal multiplexing): the entire population followed a single hidden two-state Markov chain with a self-transition probability giving a mean dwell of 60 ms, held identical across sessions, emitting from the conscious template in one state and the suppressed template in the other. The same time partition was used for all neurons so that they switched together. Uncoordinated: identical to the two-state model except that each neuron followed its own independent two-state partition (same 60 ms dwell), so switching was independent for each neuron rather than coordinated. Averaging: every neuron fired at the fixed constant mean of its two templates, $\frac{1}{2}(r_{\mathrm{conscious}} + r_{\mathrm{suppressed}})$, with no switching; all fluctuations in activity arose from Poisson noise. Conscious-only: every neuron fired from the conscious template, with no switching. The recorded data and all four models were then passed through the identical HMM state-inference, reconstruction, cross-neuron-validated reconstruction and reversal, and likelihood analyses, and each model was scored on whether it reproduced four signatures of the data, according to the criteria stated below. The criteria for passing each test were defined as follows. Reconstruction: in state 1, the explained variance of the conscious face significantly exceeds both zero and that of the suppressed face; in state 2, the explained variance of the suppressed face significantly exceeds both zero and that of the conscious face from the full population. Coordinated reconstruction: reconstruction survived cross-neuron validation where states were inferred from one random half of the neurons and features were predicted from the other half. Tuning reversal: the suppressed-state modulation indices were significantly negatively correlated with unambiguous modulation. As for coordinated reconstruction, states were inferred from one random half and modulation indices computed for the other half. Likelihood: in per-neuron held-out log-likelihood, a two-rate description, where each bin emitted from either the conscious or the suppressed template, exceeded the best single static rate, defined as the optimal fixed weighted average $p \times r_{\mathrm{conscious}} + (1 - p) \times r_{\mathrm{suppressed}}$ of the two templates ($p$ chosen by grid search of 41 uniformly spaced values from 0 to 1).

\subsection*{Data availability}

Electrophysiological data and associated metadata of this study will be made publicly available through the DANDI Archive upon publication.

\subsection*{Code availability}

Custom code used for data processing, analysis and generation of the figures will be made publicly available on GitHub upon publication.

\pdfpage{32}\section*{Acknowledgements}

We thank Mo Mao, Chang-In Sohn and Yaozu Li for technical assistance, Daniel Wagenaar for help with the Neuropixels insertion design, and members of the Tsao lab for discussions.

\subsection*{Funding}

This work was supported by the Howard Hughes Medical Institute, the NIH (grant 1R01NS138157-01 to D.Y.T. and T.A.E. and grant T32MH065214 to C.M.).

\subsection*{Author contributions}

J.K.H. and D.Y.T. conceived the project and designed the study. J.K.H., F.F.L., and Y.S. collected the data. J.K.H., C.M., T.A.E. and D.Y.T. designed the analysis methods, analysed the data, and interpreted the results. J.K.H. prepared the figures. J.K.H. and D.Y.T. wrote the initial draft, which was subsequently developed into the final manuscript by J.K.H., C.M., T.A.E. and D.Y.T.

\subsection*{Competing interests}

The authors declare no competing interests.

\subsection*{Additional Information}

Correspondence and requests for materials should be addressed to \href{mailto:janishesse@googlemail.com}{janishesse@googlemail.com} or \href{mailto:dortsao@berkeley.edu}{dortsao@berkeley.edu}.

\section*{References}

\begin{hangpars}
1 Koch, C., Massimini, M., Boly, M. \& Tononi, G. Neural correlates of consciousness: progress and problems. \textit{Nature reviews neuroscience} \textbf{17}, 307-321 (2016).

2 Hesse, J. K. \& Tsao, D. Y. A new no-report paradigm reveals that face cells encode both consciously perceived and suppressed stimuli. \textit{Elife} \textbf{9}, e58360 (2020).

3 Sheinberg, D. L. \& Logothetis, N. K. The role of temporal cortical areas in perceptual organization. \textit{Proceedings of the National Academy of Sciences} \textbf{94}, 3408-3413 (1997).

4 Wolfe, J. M. Reversing ocular dominance and suppression in a single flash. \textit{Vision research} \textbf{24}, 471-478 (1984).

5 Tong, F., Meng, M. \& Blake, R. Neural bases of binocular rivalry. \textit{Trends in cognitive sciences} \textbf{10}, 502-511 (2006).

6 Blake, R. \& Logothetis, N. K. Visual competition. \textit{Nature Reviews Neuroscience} \textbf{3}, 13-21 (2002).

7 Panagiotaropoulos, T. I., Deco, G., Kapoor, V. \& Logothetis, N. K. Neuronal discharges and gamma oscillations explicitly reflect visual consciousness in the lateral prefrontal cortex. \textit{Neuron} \textbf{74}, 924-935 (2012).

8 Leopold, D. A. \& Logothetis, N. K. Activity changes in early visual cortex reflect monkeys' percepts during binocular rivalry. \textit{Nature} \textbf{379}, 549-553 (1996).

9 Kapoor, V. \textit{et al.} Decoding internally generated transitions of conscious contents in the prefrontal cortex without subjective reports. \textit{Nature Communications} \textbf{13}, 1535 (2022).

10 Seth, A. K. \& Bayne, T. Theories of consciousness. \textit{Nature reviews neuroscience} \textbf{23}, 439-452 (2022).

\pdfpage{33}11 Consortium, C. \textit{et al.} Adversarial testing of global neuronal workspace and integrated information theories of consciousness. \textit{Nature} \textbf{642}, 133-142 (2025).

12 Lee, T. S. \& Mumford, D. Hierarchical Bayesian inference in the visual cortex. \textit{Journal of the Optical Society of America A} \textbf{20}, 1434-1448 (2003).

13 Trautmann, E. M. \textit{et al.} Large-scale high-density brain-wide neural recording in nonhuman primates. \textit{Nature Neuroscience} \textbf{28}, 1562-1575 (2025).

14 Chang, L. \& Tsao, D. Y. The code for facial identity in the primate brain. \textit{Cell} \textbf{169}, 1013-1028. e1014 (2017).

15 Logothetis, N. K. \& Schall, J. D. Neuronal correlates of subjective visual perception. \textit{Science} \textbf{245}, 761-763 (1989).

16 Engel, T. A. \textit{et al.} Selective modulation of cortical state during spatial attention. \textit{Science} \textbf{354}, 1140-1144 (2016).

17 Logothetis, N. K., Leopold, D. A. \& Sheinberg, D. L. What is rivalling during binocular rivalry? \textit{Nature} \textbf{380}, 621-624 (1996).

18 Dehaene, S. \textit{et al.} Imaging unconscious semantic priming. \textit{Nature} \textbf{395}, 597-600 (1998).

19 McCormick, P. A. Orienting attention without awareness. \textit{Journal of Experimental Psychology: Human Perception and Performance} \textbf{23}, 168 (1997).

20 Whalen, P. J. \textit{et al.} Masked presentations of emotional facial expressions modulate amygdala activity without explicit knowledge. \textit{The Journal of neuroscience} \textbf{18}, 411-418 (1998).

21 Fiebelkorn, I. C. \& Kastner, S. A rhythmic theory of attention. \textit{Trends in cognitive sciences} \textbf{23}, 87-101 (2019).

22 Fiebelkorn, I. C., Pinsk, M. A. \& Kastner, S. A dynamic interplay within the frontoparietal network underlies rhythmic spatial attention. \textit{Neuron} \textbf{99}, 842-853. e848 (2018).

23 Helfrich, R. F. \textit{et al.} Neural mechanisms of sustained attention are rhythmic. \textit{Neuron} \textbf{99}, 854-865. e855 (2018).

24 Lundqvist, M. \textit{et al.} Gamma and beta bursts underlie working memory. \textit{Neuron} \textbf{90}, 152-164 (2016).

25 Lundqvist, M., Herman, P., Warden, M. R., Brincat, S. L. \& Miller, E. K. Gamma and beta bursts during working memory readout suggest roles in its volitional control. \textit{Nature communications} \textbf{9}, 394 (2018).

26 Abdalaziz, M., Redding, Z. V. \& Fiebelkorn, I. C. Rhythmic temporal coordination of neural activity prevents representational conflict during working memory. \textit{Current Biology} \textbf{33}, 1855-1863. e1853 (2023).

27 Han, H.-B., Brincat, S. L., Buschman, T. J. \& Miller, E. K. Working memory readout varies with frontal theta rhythms. \textit{Neuron} \textbf{114}, 159-166. e152 (2026).

28 Panichello, M. F. \textit{et al.} Intermittent rate coding and cue-specific ensembles support working memory. \textit{Nature} \textbf{636}, 422-429 (2024).

29 Rich, E. L. \& Wallis, J. D. Decoding subjective decisions from orbitofrontal cortex. \textit{Nature neuroscience} \textbf{19}, 973-980 (2016).

30 Jezek, K., Henriksen, E. J., Treves, A., Moser, E. I. \& Moser, M.-B. Theta-paced flickering between place-cell maps in the hippocampus. \textit{Nature} \textbf{478}, 246-249 (2011).

31 Tsao, D. Y., Freiwald, W. A., Tootell, R. B. \& Livingstone, M. S. A cortical region consisting entirely of face-selective cells. \textit{Science} \textbf{311}, 670-674 (2006).

32 Ohayon, S., Freiwald, W. A. \& Tsao, D. Y. What makes a cell face selective? The importance of contrast. \textit{Neuron} \textbf{74}, 567-581 (2012).

33 Pachitariu, M., Sridhar, S., Pennington, J. \& Stringer, C. Spike sorting with Kilosort4. \textit{Nature methods} \textbf{21}, 914-921 (2024).

34 Pashler, H. Dual-task interference in simple tasks: data and theory. \textit{Psychological bulletin} \textbf{116}, 220 (1994).

\pdfpage{34}35 She, L., Benna, M. K., Shi, Y., Fusi, S. \& Tsao, D. Y. Temporal multiplexing of perception and memory codes in IT cortex. \textit{Nature} \textbf{629}, 861-868 (2024).
\end{hangpars}


\decompnote{Pages 35 to 50 of the PDF are the figure pages, one figure per page without running footer: Figures 1 to 6, then Extended Data Figures 1 to 10. Their legends are printed on pp.\,16--23 (above). Each page is included below as a 200-dpi render with a panel-by-panel description; the text printed inside the figures (panel letters, axis labels, keys, titles) is in the descriptions, not in the running text.}

\pdfpage{35}
\begin{center}
\includegraphics[width=0.95\linewidth,height=0.8\textheight,keepaspectratio]{figures/fig1.png}
\end{center}
\noindent\textbf{Figure 1}
\alttext{Figure 1 (PDF page 35).  (a) Schematic of the binocular flash suppression paradigm, titled Binocular flash suppression. A row labelled Stimulus shows two grey rounded boxes, each split into a Left eye half and a Right eye half (labelled below). In the first box, earlier in time, only the left eye receives a stimulus: a square greyscale photograph of a macaque monkey face on a visual-noise background; the right-eye half is empty grey. In the second box both eyes receive a stimulus: the left eye still sees the same monkey face and the right eye now sees a different macaque face (rounder, lighter fur) on a noise background. A row labelled Percept shows one face image under each box: under the first box the percept is the first monkey face (the monocular priming stimulus); under the second box the percept is the second, newly flashed face, i.e. the flashed stimulus is perceived and the priming stimulus is suppressed. A horizontal arrow pointing to the right along the bottom is labelled Time. (b) Icon strip above the two plots of panel b: for each plot two grey rounded boxes with curly brackets below them mark the stimulus epochs. Left plot: a box with a monkey face in the left half and an empty right half spans the time before the flash; a box with two different monkey faces (one per eye) spans the time after the flash. Right plot: the same with the single face in the right half of the first box. Two line plots side by side titled Left eye suppressed (left) and Right eye suppressed (right). Y axis Proportion percept from 0 to 100 percent (ticks every 20), x axis Time from flash (seconds) from -2 to 6 (ticks -2, 0, 2, 4, 6). Each plot has three curves with shaded error bands: red = left eye, blue = right eye, yellow = mixed percept; a dashed vertical line marks the flash at t = 0. Left plot (priming stimulus in the left eye): the left-eye report (red) rises from 42 percent at -2 s to 89 percent just before the flash, then falls to 17 percent at 1 s and a minimum of 8 percent at 1.7 s, recovering to 39 percent at 6 s; the right-eye report (blue) is 4 percent at the flash, jumps to 76 percent at 1 s and peaks at 82 percent at 1.5 s, then decays to 30 percent at 6 s; mixed (yellow) is 4 percent around the flash and rises to 28 percent at 6 s. Right plot (priming stimulus in the right eye): mirror image, blue peaks at 95-96 percent around the flash and falls to 8 percent at 2 s, red rises from 3 percent to 84 percent at 1.6 s and decays to 41 percent at 6 s, mixed rises from 1 percent to about 20 percent. Error bands are roughly +/- 4 percent wide. (c) Colour cartoon of the recording setup. On the left, a side view of a macaque head drawn semi-transparent so that the pink brain with its sulci is visible inside the skull. Two Neuropixels probes, thin shafts with small green headstage circuit boards and grey cables that lead up and out of the picture, enter the brain from above at two different positions, one more posterior and one more anterior. The monkey wears anaglyph goggles: the lens over the visible eye is cyan and the other lens is red. On the right stands a black-framed computer monitor on a stand; its screen shows a square anaglyph image in which a cyan-tinted monkey face is superimposed on a red-tinted house, so that the two eyes receive different images through the coloured filters. (d) Rendered CAD drawing of the custom Neuropixels insertion mechanism in an isometric view, with parts in bright colours. An octagonal grey base with a layered gold-coloured underside carries two slender probe holders that rise at different angles, one tilted left (red and magenta parts) and one tilted right (green and orange parts). Each holder is a rail with a long probe shaft, a light-green cylindrical micro-drive motor at its top, and orange and purple clamp pieces; the two probes point downward and converge towards the base so that each can be advanced independently to its target. (e) Black line drawing of a macaque brain seen from the left side, anterior to the right, with the major sulci drawn. Four yellow filled ellipses mark the recording targets, each with a text label: ML, a larger ellipse on the temporal lobe in the middle of the superior temporal sulcus; AM, a smaller ellipse at the anterior tip of the temporal lobe; PVL, a small ellipse in the lateral prefrontal cortex; and PO, the smallest ellipse just anterior and ventral to PVL near the orbital surface of the frontal lobe. (f) Two heatmaps side by side titled ML (left, y label 145 visual cells) and AM (right, y label 215 visual cells). Rows are cells, columns are 96 stimuli (x label 96 stimuli); the colour gives the normalised response, from -1 (blue) through 0 (white) to +1 (orange and dark red), colourbar at the far right labelled Response (norm.) with ticks -1 and 1. Below each heatmap a strip of six thumbnails marks the six stimulus categories of 16 images each: a human face, a crouching human body, a bunch of grapes, a camera, a thumbs-up hand, and a scrambled mosaic pattern, each on a visual-noise background. In both regions the first 16 columns (faces) are strongly orange to red for almost every cell (mean normalised response 0.47 in ML, 0.42 in AM), while the remaining 80 columns are mostly pale (category means 0.08, 0.13, 0.05, -0.02, -0.02 in ML and 0.13, 0.14, 0.13, 0.09, 0.08 in AM). ML contains several horizontal bands of light-blue (negative) responses to non-face stimuli; AM is more uniformly pale yellow outside the face block. (g) Icon strip above the raster: a grey box with a single monkey face in its left half and a curly bracket spanning -1 to 0 s (priming stimulus to one eye), and a grey box with two monkey faces and a bracket spanning 0 to 1 s (both stimuli shown, one per eye). Spike raster of one example flash suppression trial. X axis Time from flash (seconds) from -1 to 1 with ticks at -1, -0.5, 0, 0.5 and 1; y axis labelled 429 cells, one row per cell. A thin horizontal line about 38 percent of the way down separates the ML cells (upper block, labelled ML at the right) from the AM cells (lower block, labelled AM). Each spike is a short black vertical tick; 5258 spikes are drawn (2284 in ML, 2974 in AM). A dashed vertical line marks the flash at t = 0. Firing in ML is clustered into near-synchronous bursts that appear as vertical dark stripes, for example around -0.95 to -0.9, -0.75, 0.1, 0.35 and 0.55 s, and many ML rows are empty; AM firing is denser and more evenly spread, with its strongest burst at -0.9 to -0.8 s. There is no obvious change in overall firing at the flash. (h) Scatter plot of small black dots, x axis Modulation index unambiguous from -1 to 1 and y axis Modulation index binocular from -1 to 1, ticks at -1, 0 and 1 on both axes. A dashed diagonal identity line runs from (-1,-1) to (1,1), and dashed lines mark x = 0 and y = 0. The 3498 dots form a cloud elongated along a shallow positive slope: x spans the full range (5th to 95th percentile -0.56 to 0.78, median 0.17), while y stays between about -0.65 and 0.65 (5th to 95th percentile -0.23 to 0.35, median 0.04). For positive x the dots lie clearly below the identity line (92 percent of them), i.e. the binocular modulation is smaller than the unambiguous modulation; for negative x most dots lie above the line. The legend states n = 5550 cell-stimulus pairs.}
\figuredata{b}{figures/fig1b_left_eye_suppressed.csv}{figures/fig1b.py}{vector paths in the PDF (tier 2, exact to the PDF's precision), axes calibrated on the tick marks whose values were read off the rendered panel; details in the CSV header}
\figuredata{b}{figures/fig1b_right_eye_suppressed.csv}{figures/fig1b.py}{vector paths in the PDF (tier 2, exact to the PDF's precision), axes calibrated on the tick marks whose values were read off the rendered panel; details in the CSV header}
\figuredata{f}{figures/fig1f_AM.csv}{figures/fig1f.py}{embedded raster image in the PDF at native resolution (tier 3), pixel values mapped through the colour bar; details and uncertainty in the CSV header}
\figuredata{f}{figures/fig1f_ML.csv}{figures/fig1f.py}{embedded raster image in the PDF at native resolution (tier 3), pixel values mapped through the colour bar; details and uncertainty in the CSV header}
\figuredata{f}{figures/fig1f_colourbar.csv}{figures/fig1f.py}{embedded raster image in the PDF at native resolution (tier 3), pixel values mapped through the colour bar; details and uncertainty in the CSV header}
\figuredata{g}{figures/fig1g_layout.csv}{figures/fig1g.py}{vector paths in the PDF (tier 2, exact to the PDF's precision), axes calibrated on the tick marks whose values were read off the rendered panel; details in the CSV header}
\figuredata{g}{figures/fig1g_spikes.csv}{figures/fig1g.py}{vector paths in the PDF (tier 2, exact to the PDF's precision), axes calibrated on the tick marks whose values were read off the rendered panel; details in the CSV header}
\figuredata{h}{figures/fig1h.csv}{figures/fig1h.py}{vector paths in the PDF (tier 2, exact to the PDF's precision), axes calibrated on the tick marks whose values were read off the rendered panel; details in the CSV header}

\pdfpage{36}
\begin{center}
\includegraphics[width=0.95\linewidth,height=0.8\textheight,keepaspectratio]{figures/fig2.png}
\end{center}
\noindent\textbf{Figure 2}
\alttext{Figure 2 (PDF page 36).  (a) Schematic of the log-likelihood-ratio state inference. Left: a column of 10 small grey squares of varying darkness labelled 1 and a second column labelled 2 (two template firing-rate vectors, one per state, with an orange bar under column 1 and a purple bar under column 2) and the word Poisson with a drawing of a die below. Right: a spike raster labelled Neurons with 10 rows of grey tick marks over time, a light-blue vertical band highlighting one 10 ms time bin, and below the raster a strip of orange and purple blocks showing the inferred state per bin, with one bin outlined by a dotted box. An arrow leads from the highlighted bin to the formula: log of P(x given S1) divided by P(x given S2). If this ratio is greater than 0 the arrow points to State 1 (orange), if it is less than 0 to State 2 (purple). (b) Spike raster of one example trial. X axis: time from flash in seconds, from -1 to 1, with a vertical black line at 0 (the flash). Y axis: 411 cells, with a thin horizontal line separating the 157 ML cells (upper block, labelled ML) from the 254 AM cells (lower block, labelled AM). Each of the 6201 spikes is a small dark tick. The background is shaded by the inferred state in each time bin (104 epochs): purple = state 2, orange = state 1. Before the flash the whole second is purple except 21 ms of orange; after the flash orange and purple epochs alternate rapidly (median epoch 8 ms), orange covering 0.28 s of the second after the flash. Above the raster, brackets mark the priming period (-1 to 0 s, with an icon of the single priming face) and the binocular period (0 to 1 s, with icons of the two faces shown to the two eyes). Spikes in both blocks become denser and more patterned after the flash. (c) Line plot of the proportion of time points assigned to each state versus time from flash (seconds, -1 to 1), with labels Priming (before 0) and Binocular (after 0), a dashed vertical line at 0 and a dashed horizontal line at 50 percent. Y axis: Proportion state (percent), 0 to 100. Orange line = state 1, purple line = state 2, shaded bands = +/- one standard deviation across sessions (about +/- 11 points). During priming state 2 is at about 76 percent and state 1 at about 24 percent; within 0.1 to 0.2 s after the flash the two lines cross and settle near 50 percent (state 2 about 54 percent and state 1 about 46 percent at 0.9 s). (d) Autocorrelation of the log-likelihood ratio versus lag in ms (0 to 200). Y axis: log-likelihood autocorrelation, 0 to 1. Black line = real data, grey line = time-shuffled data, grey bands = +/- one standard deviation across sessions. The data autocorrelation falls from 1 at lag 0 to 0.36 at 10 ms, 0.13 at 50 ms, about -0.02 at 100 ms and -0.07 at 200 ms. The shuffled autocorrelation drops to about -0.01 already at 10 ms and stays flat. A dashed vertical grey line marks a lag of 67 ms. (e) Two dense spike rasters aligned to state transitions, x axis time from transition in ms (-50 to 50) with a thin vertical line at 0. Left raster, titled State 2 to State 1 in purple and orange: 2178 cell-transitions (rows) with 2456 spikes; a solid horizontal line at row 810 separates cells preferring the conscious face (above) from cells preferring the suppressed face (below). Above the line spikes are sparse before the transition (259) and dense after it (685); below the line the opposite (926 before, 586 after). Right raster, titled State 1 to State 2: 1936 cell-transitions, 2257 spikes, separator at row 722; cells preferring the conscious face fire more before the transition (623 versus 339 after) and cells preferring the suppressed face fire more after it (875 versus 420 before). (f) Two spike rasters of an example ML neuron that prefers the conscious face, aligned to state transitions; x axis time from transition in ms (-50 to 50) with a dashed vertical line at 0; each spike is a short vertical tick. Left, State 2 to State 1: 22 transitions (rows), 110 spikes, 27 before and 83 after the transition, i.e. the neuron starts firing right after the switch into state 1, with a cluster of spikes in the first 20 ms. Right, State 1 to State 2: 17 transitions, 87 spikes, 52 before and 35 after, i.e. firing drops after the switch into state 2. (g) Same layout as panel f for a simultaneously recorded AM neuron that prefers the suppressed face. Left, State 2 to State 1: 22 transitions, 50 spikes, 36 before and 14 after the transition (the neuron falls silent after the switch into state 1). Right, State 1 to State 2: 17 transitions, 29 spikes, 9 before and 20 after (it starts firing after the switch into state 2). (h) Row 1 ground-truth conscious face: grayscale frontal macaque face on a uniform grey background, dark fur on the crown and brow, pale muzzle. Row 1 reconstruction from state 1 activity: a smoother, slightly blurred macaque face that resembles the row's conscious face (dark brow, pale muzzle). Row 1 ground-truth suppressed face: a different macaque with lighter, fluffier fur framing a broad face. Row 1 reconstruction from state 2 activity: a smoother face resembling the row's suppressed face (light fur, broad cheeks). Row 2 ground-truth conscious face: macaque face with a light crown, dark framing fur and pale eye region. Row 2 reconstruction from state 1 activity: smoother version resembling the row's conscious face. Row 2 ground-truth suppressed face: wider macaque face with a dark crown and brow and a bright muzzle. Row 2 reconstruction from state 2 activity: smoother version resembling the row's suppressed face. Row 3 ground-truth conscious face: broad macaque face with a pale brow, dark eyes and darker lower face. Row 3 reconstruction from state 1 activity: smoother version resembling the row's conscious face. Row 3 ground-truth suppressed face: narrower macaque face with a dark fur cap and pale muzzle. Row 3 reconstruction from state 2 activity: smoother version resembling the row's suppressed face. Row 4 ground-truth conscious face: pale-furred macaque face with a light brow and dark eyes. Row 4 reconstruction from state 1 activity: smoother version resembling the row's conscious face. Row 4 ground-truth suppressed face: macaque face with pale cheeks and dark fur on the sides of the head. Row 4 reconstruction from state 2 activity: smoother version resembling the row's suppressed face. Row 5 ground-truth conscious face: macaque face with a darker ruff of fur and a pale brow and muzzle. Row 5 reconstruction from state 1 activity: smoother version resembling the row's conscious face. Row 5 ground-truth suppressed face: bright, pale-furred macaque face with white fur around the face. Row 5 reconstruction from state 2 activity: smoother version resembling the row's suppressed face. (i) Explained variance of face features (percent, -100 to 100) versus feature number 1 to 50; a bracket labelled Shape spans features 1 to 20 and a bracket labelled Appearance spans 21 to 50. Five lines: black solid = priming stimulus features predicted from average priming activity (mean 37 percent over shape features, 79 percent over appearance features, up to 98); orange solid = conscious features from state 1 (means 26 and 62 percent); purple dashed = suppressed features from state 2 (means 25 and 55 percent); these three track each other closely. Purple solid = conscious features from state 2 and orange dashed = suppressed features from state 1 are negative: around -20 to -60 percent over the shape features and mostly at or below -100 percent (clipped at the bottom of the axis) over the appearance features, with brief upward excursions to about -20 percent at features 24 and 40. (j) Decoding accuracy (percent, 0 to 100) versus number of faces (1 to 60, the target plus distractors), with shaded +/- SD bands. All five series start at 100 percent for 1 face. Black solid (priming average) decreases to 89 percent at 10 faces and 70 percent at 60; orange solid (binocular conscious, state 1) to 79 percent at 10 and 58 percent at 60; purple dashed (binocular suppressed, state 2) to 75 percent at 10 and 50 percent at 60. Purple solid (binocular conscious, state 2) and orange dashed (binocular suppressed, state 1) collapse immediately: about 28 and 24 percent at 2 faces and 0 percent from about 10 faces on. Bands are about +/- 5 to 10 points wide for the three high series. (k) Three small plots of single-trial decoding accuracy (percent, 0 to 100) versus time from flash (ticks 0, 1, 2; the axis title says ms but the values are seconds) for one example session, 10 ms steps from -0.1 to 2 s. Solid line = decoded identity compared with the conscious stimulus, dashed = with the suppressed stimulus. Top left, Average (grey, all time points in a 200 ms window): both lines settle around 36 to 38 percent after 0.5 s; before the flash the suppressed (priming) line is high, about 63 to 77 percent, and the conscious line is 0. Bottom left, State 1 (orange): conscious rises to about 60 percent and suppressed falls to about 5 percent. Bottom right, State 2 (purple): suppressed stays high at about 64 percent while conscious is about 6 percent. (l) Single-trial decoding accuracy (percent, 0 to 100) for three conditions on the x axis: Average, State 1, State 2. Thin lines = 10 individual sessions, thick black lines = mean; solid lines and filled circles = decoded identity compared with the conscious stimulus, dashed lines and open circles = compared with the suppressed stimulus; circles are grey for Average, orange for State 1, purple for State 2. Means: conscious 40 percent (Average), 61 percent (State 1), 5 percent (State 2); suppressed 34 percent, 4 percent, 60 percent. The two mean lines form a cross: conscious peaks in state 1, suppressed peaks in state 2. Individual sessions range from 41 to 83 percent for conscious in state 1 and from 40 to 78 percent for suppressed in state 2. Chance is 10 percent.}
\figuredata{b}{figures/fig2b_spikes.csv}{figures/fig2b.py}{vector paths in the PDF (tier 2, exact to the PDF's precision), axes calibrated on the tick marks whose values were read off the rendered panel; details in the CSV header}
\figuredata{b}{figures/fig2b_states.csv}{figures/fig2b.py}{vector paths in the PDF (tier 2, exact to the PDF's precision), axes calibrated on the tick marks whose values were read off the rendered panel; details in the CSV header}
\figuredata{c}{figures/fig2c.csv}{figures/fig2c.py}{vector paths in the PDF (tier 2, exact to the PDF's precision), axes calibrated on the tick marks whose values were read off the rendered panel; details in the CSV header}
\figuredata{d}{figures/fig2d.csv}{figures/fig2d.py}{vector paths in the PDF (tier 2, exact to the PDF's precision), axes calibrated on the tick marks whose values were read off the rendered panel; details in the CSV header}
\figuredata{e}{figures/fig2e_state1_to_state2.csv}{figures/fig2e.py}{vector paths in the PDF (tier 2, exact to the PDF's precision), axes calibrated on the tick marks whose values were read off the rendered panel; details in the CSV header}
\figuredata{e}{figures/fig2e_state2_to_state1.csv}{figures/fig2e.py}{vector paths in the PDF (tier 2, exact to the PDF's precision), axes calibrated on the tick marks whose values were read off the rendered panel; details in the CSV header}
\figuredata{f}{figures/fig2f_state1_to_state2.csv}{figures/fig2fg.py}{vector paths in the PDF (tier 2, exact to the PDF's precision), axes calibrated on the tick marks whose values were read off the rendered panel; details in the CSV header}
\figuredata{f}{figures/fig2f_state2_to_state1.csv}{figures/fig2fg.py}{vector paths in the PDF (tier 2, exact to the PDF's precision), axes calibrated on the tick marks whose values were read off the rendered panel; details in the CSV header}
\figuredata{g}{figures/fig2g_state1_to_state2.csv}{figures/fig2fg.py}{vector paths in the PDF (tier 2, exact to the PDF's precision), axes calibrated on the tick marks whose values were read off the rendered panel; details in the CSV header}
\figuredata{g}{figures/fig2g_state2_to_state1.csv}{figures/fig2fg.py}{vector paths in the PDF (tier 2, exact to the PDF's precision), axes calibrated on the tick marks whose values were read off the rendered panel; details in the CSV header}
\figuredata{i}{figures/fig2i.csv}{figures/fig2i.py}{vector paths in the PDF (tier 2, exact to the PDF's precision), axes calibrated on the tick marks whose values were read off the rendered panel; details in the CSV header}
\figuredata{j}{figures/fig2j.csv}{figures/fig2j.py}{vector paths in the PDF (tier 2, exact to the PDF's precision), axes calibrated on the tick marks whose values were read off the rendered panel; details in the CSV header}
\figuredata{k}{figures/fig2k.csv}{figures/fig2k.py}{vector paths in the PDF (tier 2, exact to the PDF's precision), axes calibrated on the tick marks whose values were read off the rendered panel; details in the CSV header}
\figuredata{l}{figures/fig2l.csv}{figures/fig2l.py}{vector paths in the PDF (tier 2, exact to the PDF's precision), axes calibrated on the tick marks whose values were read off the rendered panel; details in the CSV header}

\pdfpage{37}
\begin{center}
\includegraphics[width=0.95\linewidth,height=0.8\textheight,keepaspectratio]{figures/fig3.png}
\end{center}
\noindent\textbf{Figure 3}
\alttext{Figure 3 (PDF page 37).  (a) Box plot with individual points. A single box in the centre; y axis Log-likelihood (two-state minus one-state) from 0 to 0.8 with ticks every 0.2; no x axis. Nineteen dark points, one per session, range from 0.008 to 0.765 and all lie above the grey dashed horizontal line at 0. The box shows the median and interquartile range: first quartile 0.047, median 0.162, third quartile 0.409; whiskers give the range, 0.008 to 0.765. The two-state model fits better than the one-state model in every session. (b) Two dense spike rasters side by side. x axis Time from transition (ms) from -50 to 50 with a dashed vertical line at 0, the state transition. Left raster, titled State 2 to State 1 (State 2 in purple, State 1 in orange): 1196 cell-transitions as rows and 1207 spikes drawn as small dots. Right raster, State 1 to State 2: 988 rows, 1007 spikes. A solid horizontal line divides each raster; rows above it (368 on the left, 304 on the right) are cells preferring the conscious face, rows below are cells preferring the suppressed face. Left: conscious-preferring cells fire more after the transition (99 spikes before vs 258 after 0 ms) while suppressed-preferring cells fire less (503 before vs 347 after), visible as the upper part getting denser and the lower part sparser right of the dashed line. Right: the reverse (conscious-preferring 205 before vs 116 after, suppressed-preferring 226 before vs 460 after). Spike times are binned at 1 ms. (c) Two small spike rasters of one example ML cell that prefers the conscious face, with states inferred from all other neurons. x axis Time from transition (ms) from -50 to 50, dashed vertical line at 0. Left, titled State 2 to State 1: 21 transitions as rows, 95 spikes drawn as short vertical ticks, 26 before and 69 after the transition, so the cell starts firing when state 1 begins. Right, State 1 to State 2: 18 transitions, 100 spikes, 57 before and 43 after; firing thins out after the switch to state 2. (d) Two small spike rasters, same layout as panel c, for an example AM cell that prefers the suppressed face. Left, State 2 to State 1: 26 transitions, 57 spikes, 43 before and 14 after the transition, so the cell stops firing when state 1 begins. Right, State 1 to State 2: 20 transitions, 35 spikes, 10 before and 25 after; firing resumes in state 2. x axis Time from transition (ms) from -50 to 50 with a dashed line at 0. (e) Ground truth conscious face, row 1: grey-scale frontal macaque face on a uniform mid-grey square; dark fur cap and dark bushy cheeks framing a bright pale muzzle and inner face, dark eyes. Reconstruction from state 1 activity, row 1: smooth, slightly blurred macaque face with a dark upper head and pale lower face, resembling the conscious face to its left. Ground truth suppressed face, row 1: wider macaque face with lighter grey fur all around, a pale face, dark close-set eyes and dark nostrils. Reconstruction from state 2 activity, row 1: smooth blurred pale macaque face with light fur, resembling the suppressed face to its left. Ground truth conscious face, row 2: macaque face with a dark ring of fur around a pale face, dark eyes and a narrow muzzle. Reconstruction from state 1 activity, row 2: smooth pale macaque face with dark eyes and darker fur above, similar to the conscious face. Ground truth suppressed face, row 2: lighter macaque face with prominent brow ridges, pale cheeks with whiskers and dark eyes. Reconstruction from state 2 activity, row 2: smooth broad pale macaque face with dark eyes, similar to the suppressed face. Ground truth conscious face, row 3: darker macaque face with large dark eye regions and a lighter muzzle, head slightly tilted. Reconstruction from state 1 activity, row 3: smooth pale macaque face with dark eyes and a darker crown. Ground truth suppressed face, row 3: macaque with very fluffy light fur around a pale face and dark eyes. Reconstruction from state 2 activity, row 3: smooth pale macaque face with light fur and dark eyes, similar to the suppressed face. Ground truth conscious face, row 4: macaque face with pale grey fur, a pale face and dark eyes. Reconstruction from state 1 activity, row 4: smooth pale macaque face with dark eyes. Ground truth suppressed face, row 4: the brightest face of the grid, nearly white fur and face with strongly contrasting dark eyes and nostrils. Reconstruction from state 2 activity, row 4: smooth bright pale macaque face with dark eyes, similar to the bright suppressed face. Ground truth conscious face, row 5: macaque face with dark fur and a dark face, pale muzzle and dark eyes. Reconstruction from state 1 activity, row 5: smooth pale macaque face with dark eyes and a dark crown. Ground truth suppressed face, row 5: pale macaque face with dark eye patches and scruffy fur, gaze slightly downward. Reconstruction from state 2 activity, row 5: smooth pale macaque face with dark eyes, similar to the suppressed face. (f) Line plot. x axis Feature number from 1 to 50 (ticks 10 to 50); brackets above the plot mark features 1-20 as Shape and 21-50 as Appearance. y axis Explained variance of features (percent) from -100 to 100, ticks every 50. Five series: black solid line, priming average, highest (mean 37 percent over the shape features and 79 percent over the appearance features, up to 98); orange solid, binocular conscious state 1 (mean 16 shape, 39 appearance); purple dashed, binocular suppressed state 2 (mean 10 shape, 22 appearance); purple solid, binocular conscious state 2 (mean -9 shape, -27 appearance, down to -58); orange dashed, binocular suppressed state 1, lowest (mean -16 shape, -60 appearance, touching the -100 axis limit at feature 25). All lines jump up at feature 21. Thus state 1 activity explains the conscious face features but not the suppressed ones, and state 2 activity explains the suppressed face features better than the conscious ones. Legend below the plot. (g) Line plot with shaded error bands. x axis Number of faces from 0 to 60 (ticks every 20); y axis Decoding accuracy (percent) from 0 to 100 (ticks every 20). All five series start at 100 percent with one face and fall as faces are added: black priming average to 70 percent at 60 faces (89 at 10 faces); orange solid, conscious face decoded from state 1, to 32 percent (59 at 10); purple dashed, suppressed face decoded from state 2, to 12 percent (39 at 10); purple solid, conscious face from state 2, to 1.7 percent (12 at 10); orange dashed, suppressed face from state 1, to about 0 percent (3.7 at 10). Light bands of the same colour are symmetric about each line, for example about +/- 6 percent around the black line and +/- 7 percent around the orange solid line at 30 faces, wider (+/- 8) around the purple dashed line. Legend below, shared with panel f. (h) Three small line plots of decoding accuracy (percent, y axis 0 to 100 with ticks 0, 50, 100) against time from flash (x axis -0.1 to 2 with ticks 0, 1, 2; labelled ms in the figure, but the range is seconds) for one example session: top left Average in grey, bottom left State 1 in orange, bottom right State 2 in purple. Solid lines are the conscious face, dashed lines the suppressed face (legend top right). In all three plots the solid line starts near 0-7 percent before the flash and the dashed line near 55-65 percent, and they cross within about 0.2 s. After 0.5 s: Average, conscious about 35 percent (range 24-44) and suppressed about 29; State 1, conscious about 39 and suppressed about 17; State 2, conscious about 21 and suppressed about 34. So state 1 activity decodes the conscious face and state 2 activity the suppressed face. (i) Paired line plot. x axis with three categories Average, State 1, State 2; y axis Decoding accuracy (percent) from 0 to 100, ticks every 20. Ten thin grey lines per condition (one per session): solid for the conscious face, dashed for the suppressed face. Thick black lines and large circles are the means: filled circles conscious, open circles suppressed, coloured grey for Average, orange for State 1 and purple for State 2 (legend top left). Means: conscious 39.5 percent for Average, 45.4 for State 1, 20.0 for State 2; suppressed 29.6, 13.5 and 37.1. Session values range 31-52 (conscious, Average), 35-65 (conscious, State 1), 6-34 (conscious, State 2), 18-49 (suppressed, Average), 1-44 (suppressed, State 1) and 30-48 (suppressed, State 2). The two mean lines cross between State 1 and State 2: conscious decoding is best from state 1, suppressed decoding is best from state 2.}
\figuredata{a}{figures/fig3a_box.csv}{figures/fig3a.py}{vector paths in the PDF (tier 2, exact to the PDF's precision), axes calibrated on the tick marks whose values were read off the rendered panel; details in the CSV header}
\figuredata{a}{figures/fig3a_points.csv}{figures/fig3a.py}{vector paths in the PDF (tier 2, exact to the PDF's precision), axes calibrated on the tick marks whose values were read off the rendered panel; details in the CSV header}
\figuredata{b}{figures/fig3b_meta.csv}{figures/fig3b.py}{vector paths in the PDF (tier 2, exact to the PDF's precision), axes calibrated on the tick marks whose values were read off the rendered panel; details in the CSV header}
\figuredata{b}{figures/fig3b_state1_to_state2.csv}{figures/fig3b.py}{vector paths in the PDF (tier 2, exact to the PDF's precision), axes calibrated on the tick marks whose values were read off the rendered panel; details in the CSV header}
\figuredata{b}{figures/fig3b_state2_to_state1.csv}{figures/fig3b.py}{vector paths in the PDF (tier 2, exact to the PDF's precision), axes calibrated on the tick marks whose values were read off the rendered panel; details in the CSV header}
\figuredata{c}{figures/fig3c_state1_to_state2.csv}{figures/fig3c.py}{vector paths in the PDF (tier 2, exact to the PDF's precision), axes calibrated on the tick marks whose values were read off the rendered panel; details in the CSV header}
\figuredata{c}{figures/fig3c_state2_to_state1.csv}{figures/fig3c.py}{vector paths in the PDF (tier 2, exact to the PDF's precision), axes calibrated on the tick marks whose values were read off the rendered panel; details in the CSV header}
\figuredata{d}{figures/fig3d_state1_to_state2.csv}{figures/fig3d.py}{vector paths in the PDF (tier 2, exact to the PDF's precision), axes calibrated on the tick marks whose values were read off the rendered panel; details in the CSV header}
\figuredata{d}{figures/fig3d_state2_to_state1.csv}{figures/fig3d.py}{vector paths in the PDF (tier 2, exact to the PDF's precision), axes calibrated on the tick marks whose values were read off the rendered panel; details in the CSV header}
\figuredata{f}{figures/fig3f.csv}{figures/fig3f.py}{vector paths in the PDF (tier 2, exact to the PDF's precision), axes calibrated on the tick marks whose values were read off the rendered panel; details in the CSV header}
\figuredata{g}{figures/fig3g.csv}{figures/fig3g.py}{vector paths in the PDF (tier 2, exact to the PDF's precision), axes calibrated on the tick marks whose values were read off the rendered panel; details in the CSV header}
\figuredata{h}{figures/fig3h_average.csv}{figures/fig3h.py}{vector paths in the PDF (tier 2, exact to the PDF's precision), axes calibrated on the tick marks whose values were read off the rendered panel; details in the CSV header}
\figuredata{h}{figures/fig3h_state1.csv}{figures/fig3h.py}{vector paths in the PDF (tier 2, exact to the PDF's precision), axes calibrated on the tick marks whose values were read off the rendered panel; details in the CSV header}
\figuredata{h}{figures/fig3h_state2.csv}{figures/fig3h.py}{vector paths in the PDF (tier 2, exact to the PDF's precision), axes calibrated on the tick marks whose values were read off the rendered panel; details in the CSV header}
\figuredata{i}{figures/fig3i.csv}{figures/fig3i.py}{vector paths in the PDF (tier 2, exact to the PDF's precision), axes calibrated on the tick marks whose values were read off the rendered panel; details in the CSV header}

\pdfpage{38}
\begin{center}
\includegraphics[width=0.95\linewidth,height=0.8\textheight,keepaspectratio]{figures/fig4.png}
\end{center}
\noindent\textbf{Figure 4}
\alttext{Figure 4 (PDF page 38).  (a) Four scatter plots in a 2 by 2 grid (columns: Conscious face, Suppressed face; rows: State 1, State 2) for one example cell. Each plot has 1500 dots, one per image, placed at the image's projection onto the cell's preferred axis (x, -1 to 1) and onto the principal orthogonal axis (y, -1 to 1), forming a round cloud. Dot colour is the cell's z-scored response on a blue-white-red scale from -3 to 6 (colour bar at the top centre; white is 1.5). In Conscious/State 1 and Suppressed/State 2 the dots on the right (high projection) are mostly red and those on the left blue, a ramp along the preferred axis (correlation of response with x about 0.5 and 0.4); in Conscious/State 2 and Suppressed/State 1 the colours are mixed with no gradient (correlations 0.1 and -0.05). Centre inset: mean response (y, 0 to 9) in six bins of the preferred-axis projection from -0.5 to 0.8 with error bars, rising from about 0.6 to 7.3 for the conscious face in state 1 (orange) and from 1.1 to 5.3 for the suppressed face in state 2 (purple). (b) Four heatmaps in a 2 by 2 grid (columns: Conscious, Suppressed; rows: State 1, State 2), each labelled 554 cells (rows, sorted by tuning gain) by 1500 stimuli (columns, sorted by the projection onto each cell's preferred axis). Colour is the z-scored response on a blue-white-red scale from -3 to 6 (colour bar at the top centre). Conscious/State 1 and Suppressed/State 2 show a left-to-right gradient from blue to red that is strongest in the top rows (high-gain cells: mean about 0.4 on the left versus 2.2 on the right in the top rows), while Suppressed/State 1 and Conscious/State 2 look uniformly speckled pink and blue with only a weak gradient. A few horizontal blue stripes mark cells with low responses to all stimuli. The embedded images are 169 by 169 pixels, i.e. heavily downsampled from the stated 554 by 1500 grid. (c) Three stacked histograms of single-cell tuning gain (x from -0.5 to 2.5, ticks at 0 and 2; y is the number of neurons, ticks at 0 and 100), drawn as translucent bars with stair-step outlines, with a legend at the right. Top: state 1 conscious (orange, bins 0.2 wide) peaks at gain 0.6 to 0.8 with 115 neurons and has a tail to 2.6; state 1 suppressed (light orange) peaks at 0.0 to 0.2 with 201 neurons and is concentrated between -0.2 and 0.6. Middle: state 2 suppressed (dark purple, 0.2 bins) peaks at 0.4 to 0.6 with 134 neurons; state 2 conscious (light purple, 0.1 bins) peaks at 0.2 to 0.3 with 86 neurons and is shifted toward lower gains. Bottom: state 1 conscious (orange) and state 2 suppressed (purple) overlaid; the two distributions nearly coincide, the orange one slightly wider toward higher gains. (d) Scatter plot of 531 small black dots (one per cell): x is the gain for the conscious face in state 1, y the gain for the suppressed face in state 2, both axes from -0.5 to 2.5 with ticks at 0 and 2; a dashed black identity diagonal and dashed lines at x = 0 and y = 0. The dots form a tight elongated cloud along the diagonal (r = 0.90), mostly between 0.2 and 1.5, lying slightly below the identity line (mean x 0.81, mean y 0.67), so gains are higher for the conscious face in state 1. (e) Scatter plot of 531 small black dots (one per cell): x is the mean response across stimuli in state 1 (normalised), y the same in state 2, both axes from 0 to 3 with ticks at 0 and 3, and a dashed identity diagonal. The dots lie tightly along the diagonal between about 0.3 and 2.9 (r = 0.95), slightly above it (mean x 1.70, mean y 1.76), so mean responses are slightly larger in state 2. (f) Line plot of decoding accuracy in percent (y from 0 to 100, ticks 0, 50, 100) for discriminating state 1 from state 2 against the number of cells (x from 0 to 1000). A thick black line with a grey band (s.d. over 100 iterations, about plus or minus 2 to 3 points) and a dashed line at the 50 percent chance level. Accuracy starts at 50.7 percent for 1 cell, rises steeply to about 60 percent at 50 cells and 63 percent at 70 to 100 cells, dips to 59 percent at 200 cells, then increases slowly to 64 percent at 500 and 66 percent at 1000 cells. (g) Four scatter plots in a 2 by 2 grid of modulation index binocular (y) against modulation index unambiguous (x), both from -1 to 1 with ticks at -1, 0 and 1; each has dashed lines at x = 0 and y = 0 and a dashed diagonal, and shows 2885 small black dots (cell-stimulus pairs; 1612 drawn in the last plot). Top left, All states: a flat cloud around y = 0 with a shallow positive trend (r = 0.45) and the identity diagonal. Top right, State 1 (title in orange): a cloud elongated along the identity diagonal (r = 0.66). Bottom left, State 2 (title in purple): a cloud elongated along the anti-diagonal from top left to bottom right (r = -0.44), mostly below y = 0. Bottom right, Slope-equalized: a round cloud centred near zero with no trend (r about 0) and the identity diagonal.}
\figuredata{a}{figures/fig4a.csv}{figures/fig4a.py}{vector paths in the PDF (tier 2, exact to the PDF's precision), axes calibrated on the tick marks whose values were read off the rendered panel; details in the CSV header}
\figuredata{a}{figures/fig4a_inset.csv}{figures/fig4a.py}{vector paths in the PDF (tier 2, exact to the PDF's precision), axes calibrated on the tick marks whose values were read off the rendered panel; details in the CSV header}
\figuredata{b}{figures/fig4b_state1_conscious.csv}{figures/fig4b.py}{embedded raster image in the PDF at native resolution (tier 3), pixel values mapped through the colour bar; details and uncertainty in the CSV header}
\figuredata{b}{figures/fig4b_state1_suppressed.csv}{figures/fig4b.py}{embedded raster image in the PDF at native resolution (tier 3), pixel values mapped through the colour bar; details and uncertainty in the CSV header}
\figuredata{b}{figures/fig4b_state2_conscious.csv}{figures/fig4b.py}{embedded raster image in the PDF at native resolution (tier 3), pixel values mapped through the colour bar; details and uncertainty in the CSV header}
\figuredata{b}{figures/fig4b_state2_suppressed.csv}{figures/fig4b.py}{embedded raster image in the PDF at native resolution (tier 3), pixel values mapped through the colour bar; details and uncertainty in the CSV header}
\figuredata{c}{figures/fig4c.csv}{figures/fig4c.py}{vector paths in the PDF (tier 2, exact to the PDF's precision), axes calibrated on the tick marks whose values were read off the rendered panel; details in the CSV header}
\figuredata{d}{figures/fig4d.csv}{figures/fig4d.py}{vector paths in the PDF (tier 2, exact to the PDF's precision), axes calibrated on the tick marks whose values were read off the rendered panel; details in the CSV header}
\figuredata{e}{figures/fig4e.csv}{figures/fig4e.py}{vector paths in the PDF (tier 2, exact to the PDF's precision), axes calibrated on the tick marks whose values were read off the rendered panel; details in the CSV header}
\figuredata{f}{figures/fig4f.csv}{figures/fig4f.py}{vector paths in the PDF (tier 2, exact to the PDF's precision), axes calibrated on the tick marks whose values were read off the rendered panel; details in the CSV header}
\figuredata{g}{figures/fig4g.csv}{figures/fig4g.py}{vector paths in the PDF (tier 2, exact to the PDF's precision), axes calibrated on the tick marks whose values were read off the rendered panel; details in the CSV header}

\pdfpage{39}
\begin{center}
\includegraphics[width=0.95\linewidth,height=0.8\textheight,keepaspectratio]{figures/fig5.png}
\end{center}
\noindent\textbf{Figure 5}
\alttext{Figure 5 (PDF page 39).  (a) Schematic of the Hidden Markov Model. Top left: two vertical columns of eight grey squares in shades from light to dark grey, the two response templates (one value per neuron), labelled 1 and 2 at the bottom with an orange square under column 1 (state 1) and a purple square under column 2 (state 2). Top right: a spike raster labelled Neurons on the y axis, grey tick marks for the spikes of about nine neurons over time, with one time bin highlighted by a translucent blue vertical bar and a small dotted box under it; an arrow leads from that bin down to the emission term. An arrow from the templates points down to the word Poisson (script font) beside a drawing of a die; an arrow from the die points to the emission probability P(x(t) given S(t)), written above a multiplication sign and the transition probability P(S(t) given S(t-1)). To the right a state-transition diagram: two rows of four squares over time, orange squares on top (state 1) and purple squares below (state 2), connected by arrows between consecutive time steps labelled p(1,1) (orange to orange), p(2,1) (purple to orange), p(1,2) (orange to purple) and p(2,2) (purple to purple), with a horizontal axis line labelled time; an arrow from the diagram points to the transition term. Below: the joint probability P(x, S) = product over t of P(x(t) given S(t)) times P(S(t) given S(t-1)). A downward arrow labelled argmax over all state sequences leads to the result: a horizontal bar of alternating orange and purple segments (state 1 and state 2 epochs over time), with a bracket over one purple segment labelled dwell. (b) Heatmap of the inferred HMM state for every trial of an example session: rows are trials (axis labelled 800 trials; about 849 rows, sorted by proportion of state 1, top rows have the most state 1), columns are time from flash in seconds from -1.02 to 1.03 (ticks -1, 0, 1), each cell orange for state 1 or purple for state 2. Before the flash (negative times) the plot is almost entirely purple (about 7 percent orange, as short orange streaks) apart from a thin orange strip at the very left edge near -1.0 s. After the flash orange appears from about 0.1 s on and dominates the upper trials: between 0.1 and 0.5 s about 82 percent of the cells in the top third of rows are orange, 55 percent in the middle third and 22 percent in the bottom third, so there is a sharp vertical boundary at about 0.1 s and a purple-to-orange gradient from the bottom rows to the top rows. Overall 47 percent of the cells after the flash and 28 percent of all cells are orange. (c) Two stacked line plots against HMM bin size in ms on a logarithmic x axis with ticks 2, 3, 5, 10, 20, 50 and 100; data points with circle markers at 2, 3, 5, 10, 20, 50, 75 and 100 ms and shaded s.e.m. bands across sessions. Top: normalised conditional log-likelihood ratio (black), y from 0.6 to 1: 0.85 at 2 ms, 0.93 at 3 ms, a flat peak of 0.97 at 5 and 10 ms, 0.94 at 20 ms, 0.86 at 50 ms, 0.76 at 75 ms and 0.69 at 100 ms; the band is about +/- 0.02 to 0.04 wide. Bottom: median dwell time in ms, y from 0 to 200, state 1 in orange and state 2 in purple: both rise monotonically from 55 and 52 ms at 2 ms bins, 60 to 68 ms at 3 to 10 ms, 80 and 77 ms at 20 ms, 126 and 110 ms at 50 ms, 155 and 146 ms at 75 ms, to 200 and 195 ms at 100 ms (the state 1 value sits on the 200 ms axis limit); state 1 dwells slightly longer than state 2 at every bin size and the bands are about +/- 5 ms wide. (d) Line plot of the normalised conditional log-likelihood ratio (y, 0.5 to 1) against median dwell time in ms (x, 0 to 250, binned) for state 1 (orange) and state 2 (purple), with circle markers at the bin centres 15, 35, 60, 85, 150 and 250 ms and shaded s.e.m. bands across sessions. Both curves rise steeply from 0.59 at 15 ms (identical for the two states) through 0.87 (state 1) and 0.90 (state 2) at 35 ms to a peak of 0.96 to 0.97 at 60 ms, stay at 0.95 at 85 ms, then decline to 0.84 and 0.83 at 150 ms and 0.70 and 0.69 at 250 ms. The bands are narrow near the peak (about +/- 0.01) and widen at the extremes (about +/- 0.065 at 15 ms and +/- 0.045 at 250 ms). (e) Two line plots of explained variance of face features in percent (y, -100 to 100) against feature number 1 to 50 (x ticks 10, 20, 30, 40, 50), left titled ML -> AM (states inferred from ML, features reconstructed from AM) and right AM -> ML. Five series per plot, legend below: priming average (black solid), binocular conscious state 1 (orange solid), binocular conscious state 2 (purple solid), binocular suppressed state 1 (orange dashed) and binocular suppressed state 2 (purple dashed). Left plot: the priming average is high, about 20 percent on average over features 1 to 20 (range -19 to 72) and 73 percent over features 21 to 50 (range 18 to 94); conscious state 1 fluctuates around zero (range -31 to 31); conscious state 2 is around zero for features 1 to 20 but mostly negative for 21 to 50 (mean -20, minimum -38 at feature 31); suppressed state 1 stays near zero to slightly negative (mean -7 over 21 to 50); suppressed state 2 is positive over features 21 to 50 (mean 15, maximum 24). Right plot: priming average 33 percent over features 1 to 20 and 76 percent over 21 to 50 (maximum 98 at feature 23); conscious state 1 is mostly positive over 21 to 50 (mean 20, maximum 46 at feature 24); conscious state 2 hovers around zero (range -24 to 41 over 21 to 50); suppressed state 1 is strongly negative over 21 to 50 (mean -36, minimum -84 at feature 50); suppressed state 2 is slightly negative (mean -8 over 21 to 50). All series dip around feature 15 in both plots.}
\figuredata{b}{figures/fig5b_state1_fraction.csv}{figures/fig5b.py}{vector paths in the PDF (tier 2, exact to the PDF's precision), axes calibrated on the tick marks; raster parts from the embedded images (tier 3); details in the CSV header}
\figuredata{c}{figures/fig5c.csv}{figures/fig5c.py}{vector paths in the PDF (tier 2, exact to the PDF's precision), axes calibrated on the tick marks whose values were read off the rendered panel; details in the CSV header}
\figuredata{d}{figures/fig5d.csv}{figures/fig5d.py}{vector paths in the PDF (tier 2, exact to the PDF's precision), axes calibrated on the tick marks whose values were read off the rendered panel; details in the CSV header}
\figuredata{e}{figures/fig5e_AM_to_ML.csv}{figures/fig5e.py}{vector paths in the PDF (tier 2, exact to the PDF's precision), axes calibrated on the tick marks whose values were read off the rendered panel; details in the CSV header}
\figuredata{e}{figures/fig5e_ML_to_AM.csv}{figures/fig5e.py}{vector paths in the PDF (tier 2, exact to the PDF's precision), axes calibrated on the tick marks whose values were read off the rendered panel; details in the CSV header}

\pdfpage{40}
\begin{center}
\includegraphics[width=0.95\linewidth,height=0.8\textheight,keepaspectratio]{figures/fig6.png}
\end{center}
\noindent\textbf{Figure 6}
\alttext{Figure 6 (PDF page 40).  (a) Row 1 ground truth conscious face: frontal grayscale macaque face on a mid-grey square, dark fur cap and side fringe, bright cheeks and muzzle, dark eyes looking straight ahead. Row 1 reconstruction from state 1 activity: smoother, slightly blurred macaque face that closely resembles the conscious face of row 1 (same dark fur cap, bright cheeks, straight gaze). Row 1 ground truth suppressed face: broader frontal macaque face with lighter, less contrasted fur, wider-set dark eye sockets and a pale muzzle. Row 1 reconstruction from state 2 activity: blurred, smooth macaque face with a broad outline, wide-set eyes and pale muzzle, resembling the suppressed face of row 1. Row 2 ground truth conscious face: wide macaque face with a broad dark fur cap, bright central face, small close-set eyes and a pale muzzle. Row 2 reconstruction from state 1 activity: smooth wide macaque face with a dark cap and bright central face, resembling the conscious face of row 2. Row 2 ground truth suppressed face: macaque face with a bright shaggy brow ruff, large round dark eyes and a pale, rather flat face. Row 2 reconstruction from state 2 activity: smooth macaque face with large round eyes and a pale face, resembling the suppressed face of row 2. Row 3 ground truth conscious face: macaque face with a bright shaggy fur cap, very large round dark eyes with bright rims and a dark mouth. Row 3 reconstruction from state 1 activity: smooth macaque face with a dark cap and large dark eyes, a blurred version of the conscious face of row 3. Row 3 ground truth suppressed face: darker, narrower macaque face with dark fur all around, dark eyes and a grey muzzle. Row 3 reconstruction from state 2 activity: smooth macaque face with a dark cap, moderately large eyes and a pale muzzle, closer to the suppressed face of row 3. Row 4 ground truth conscious face: narrower macaque face with a dark fur cap, dark eyes and a slightly darker muzzle region. Row 4 reconstruction from state 1 activity: smooth macaque face with a dark cap and pale cheeks, resembling the conscious face of row 4. Row 4 ground truth suppressed face: broad bright macaque face with a prominent bright brow ridge, wide-set dark eyes and a short pale muzzle. Row 4 reconstruction from state 2 activity: smooth broad macaque face with a flat dark cap and wide-set eyes, resembling the suppressed face of row 4. Row 5 ground truth conscious face: wide macaque face with a dark fur cap, bright cheeks, deep-set dark eyes and a pale muzzle. Row 5 reconstruction from state 1 activity: smooth wide macaque face with a dark cap and bright cheeks, resembling the conscious face of row 5. Row 5 ground truth suppressed face: macaque face with light shaggy fur, big round dark eyes and a startled expression. Row 5 reconstruction from state 2 activity: smooth macaque face with large round eyes and light fur, resembling the suppressed face of row 5. (b) Line plot of explained variance of face features in percent (y axis from -100 to 100) against feature number (x axis 1 to 50; brackets above mark features 1-20 as shape and 21-50 as appearance features), prefrontal cortex. Five series: priming average (black solid), binocular conscious state 1 (orange solid), binocular conscious state 2 (purple solid), binocular suppressed state 1 (orange dashed), binocular suppressed state 2 (purple dashed). For shape features all series fluctuate around zero (series means between -5 and -19 percent). For appearance features the priming average (mean 49 percent, peaks near 80), conscious state 1 (mean 43 percent, up to 75) and suppressed state 2 (mean 33 percent, up to 62) are high and follow each other closely, whereas conscious state 2 and suppressed state 1 are strongly negative: they dive below the -100 axis limit at many appearance features and come back up to between -18 and -94 percent in between (the purple solid line is visible at 16 of the 30 appearance features, the orange dashed line at only 5). (c) Line plot of decoding accuracy in percent (y axis 0 to 100) against number of faces (x axis 0 to 60) with shaded bands, prefrontal cortex; same five series and colours as panel b. All curves start at 100 percent for 1 face and decay. Priming average (black solid), binocular conscious state 1 (orange solid) and binocular suppressed state 2 (purple dashed) decay slowly: about 44, 44 and 38 percent at 5 faces, 15, 14 and 14 percent at 20 faces, and 5, 7 and 10 percent at 60 faces. Binocular conscious state 2 (purple solid) and binocular suppressed state 1 (orange dashed) collapse to near zero: about 4 percent at 5 faces and 0 to 1 percent beyond 10 faces. The shaded bands (spread across stimulus pairs) are widest, about plus or minus 20 to 30 percent, for 2 to 3 faces and narrow to a few percent at 60 faces. (d) Line plot of explained variance of face features in percent (y axis -100 to 100) against feature number (1 to 50, shape 1-20 and appearance 21-50), computed with the held-out half of neurons, prefrontal cortex; same five series and colours as panel b. Shape features are near zero for all series (means -1 to -8 percent). For appearance features the priming average (black) stays high (mean 48 percent, peaks near 80); binocular conscious state 1 (orange solid) is moderately positive (mean 12 percent, up to 25); binocular suppressed state 2 (purple dashed) is around zero (mean 2 percent, range -17 to 12); binocular conscious state 2 (purple solid) is negative (mean -19 percent); binocular suppressed state 1 (orange dashed) is the most negative (mean -35 percent, down to -55). No series leaves the axes. (e) Line plot of decoding accuracy in percent (0 to 100) against number of faces (0 to 60) with shaded bands, using the held-out half of neurons, prefrontal cortex; same five series and colours as panel c. All curves start at 100 percent for 1 face. The priming average (black) decays from 77 percent at 2 faces to 23 at 10, 12 at 20 and 3 at 60 faces. Binocular conscious state 1 (orange solid) and binocular suppressed state 2 (purple dashed) are next: 13 to 15 percent at 10 faces, 5 to 10 at 20 and 0 to 3 at 60. Binocular conscious state 2 (purple solid) and binocular suppressed state 1 (orange dashed) are lowest: about 7 percent at 10 faces, 4 at 20 and 2 at 60. The separation between the two pairs is much smaller than in panel c; bands overlap broadly. (f) Four square scatter plots side by side of modulation index binocular (y axis -1 to 1) against modulation index unambiguous (x axis -1 to 1), one small black dot per cell-stimulus pair (n = 283), each with dashed horizontal and vertical zero lines and a dashed diagonal reference line. All states (left): points form a cloud around the origin, elongated along the x axis, with a weak positive trend (Pearson r = 0.56); the reference diagonal is the identity line y = x and the cloud is much flatter than it. State 1 (title in orange): points scatter along the identity line with a clear positive relation (r = 0.63), extending from about (-0.8, -0.8) to (0.6, 0.6). State 2 (title in purple): the reference line is the anti-diagonal y = -x and the points follow it with a negative relation (r = -0.32), mostly in the upper-left and lower-right quadrants. Slope-equalized (right): 126 visible points around the origin with little relation to the identity line (r = 0.14 for the visible points; the legend reports r = 0.24 for all 283 pairs; the remaining points lie outside the plotted range).}
\figuredata{b}{figures/fig6b.csv}{figures/fig6b.py}{vector paths in the PDF (tier 2, exact to the PDF's precision), axes calibrated on the tick marks whose values were read off the rendered panel; details in the CSV header}
\figuredata{c}{figures/fig6c.csv}{figures/fig6c.py}{vector paths in the PDF (tier 2, exact to the PDF's precision), axes calibrated on the tick marks whose values were read off the rendered panel; details in the CSV header}
\figuredata{d}{figures/fig6d.csv}{figures/fig6d.py}{vector paths in the PDF (tier 2, exact to the PDF's precision), axes calibrated on the tick marks whose values were read off the rendered panel; details in the CSV header}
\figuredata{e}{figures/fig6e.csv}{figures/fig6e.py}{vector paths in the PDF (tier 2, exact to the PDF's precision), axes calibrated on the tick marks whose values were read off the rendered panel; details in the CSV header}
\figuredata{f}{figures/fig6f_all_states.csv}{figures/fig6f.py}{vector paths in the PDF (tier 2, exact to the PDF's precision), axes calibrated on the tick marks whose values were read off the rendered panel; details in the CSV header}
\figuredata{f}{figures/fig6f_slope_equalized.csv}{figures/fig6f.py}{vector paths in the PDF (tier 2, exact to the PDF's precision), axes calibrated on the tick marks whose values were read off the rendered panel; details in the CSV header}
\figuredata{f}{figures/fig6f_state1.csv}{figures/fig6f.py}{vector paths in the PDF (tier 2, exact to the PDF's precision), axes calibrated on the tick marks whose values were read off the rendered panel; details in the CSV header}
\figuredata{f}{figures/fig6f_state2.csv}{figures/fig6f.py}{vector paths in the PDF (tier 2, exact to the PDF's precision), axes calibrated on the tick marks whose values were read off the rendered panel; details in the CSV header}

\pdfpage{41}
\begin{center}
\includegraphics[width=0.95\linewidth,height=0.8\textheight,keepaspectratio]{figures/edfig1.png}
\end{center}
\noindent\textbf{Extended Data Figure 1}
\alttext{Extended Data Figure 1 (PDF page 41).  (a) Coronal T1-weighted MRI slice through the temporal lobe of a macaque, grey brain on a black background, dorsal at the top, with the faces-versus-objects fMRI contrast overlaid in colour. A bright yellow-orange blob lies on the lateral surface of the right inferotemporal cortex, at the lower right of the brain (about 92 percent of the image width from the left and 62 percent of the height from the top). A white upward-pointing arrowhead just below it is labelled ML (middle lateral face patch) and marks the recording target. A smaller orange spot sits at the mirror position on the left lateral surface. Top left: a small vertical colourbar running from yellow (top, p = 1e-10) through orange and red to dark red (bottom, p = 1e-3), labelled p. Left middle: a white vertical scale bar labelled 1 cm; it is 12.8 points long in an 84.6-point-wide image, so the slice is about 6.6 cm across. (b) Coronal MRI slice at a more anterior level of the temporal lobe, grey brain on black, dorsal at the top, with the faces-versus-objects contrast overlay. A yellow-orange blob on the ventral surface of the right anterior temporal lobe (lower right, about 79 percent of the width from the left and 84 percent of the height from the top) has a white upward-pointing arrowhead beneath it labelled AM (anterior medial face patch), the recording target. A larger bright yellow blob sits at the mirror position on the ventral left temporal lobe, and a small blue spot (the opposite contrast) lies on the right lateral surface, up and to the right of the AM blob. Top left: colourbar from yellow (top, p = 1e-10) to dark red (bottom, p = 1e-3), labelled p. Left middle: white vertical 1 cm scale bar (14.2 points in the 84.6-point image, so the slice is about 6.0 cm across). (c) Coronal MRI slice through the prefrontal cortex, dorsal at the top, with two bright textured orbital structures below the brain, and the faces-versus-objects contrast overlay. A small yellow-orange spot sits on the ventral prefrontal surface of the right hemisphere (about 71 percent of the width from the left and 68 percent of the height from the top); a white upward-pointing arrowhead below it is labelled PO (prefrontal orbital face patch), the recording target. A few further small orange spots lie near the ventral midline, at the far left edge and at the lower right edge of the brain. Top left: colourbar from yellow (top, p = 1e-5) to dark red (bottom, p = 1e-2), labelled p. Left: white vertical 1 cm scale bar (20.1 points in the 85.9-point image, so the slice is about 4.3 cm across, a higher magnification than panels a and b). (d) Coronal MRI slice through the prefrontal cortex at a level similar to panel c, dorsal at the top, bright textured orbital structures below the brain and a bright vertical streak along the midline. Faces-versus-objects contrast overlay: an orange spot sits on the lateral prefrontal surface of the right hemisphere (about 82 percent of the width from the left and 73 percent of the height from the top); a white arrowhead pointing right at it is labelled PVL (prefrontal ventrolateral face patch), the recording target. Two further small orange spots lie near the centre of the slice and at the lower left of the brain. Top left: colourbar from yellow (top, p = 1e-5) to dark red (bottom, p = 1e-2), labelled p. Left: white vertical 1 cm scale bar (20.1 points in the 85.4-point image, so the slice is about 4.3 cm across).}

\pdfpage{42}
\begin{center}
\includegraphics[width=0.95\linewidth,height=0.8\textheight,keepaspectratio]{figures/edfig2.png}
\end{center}
\noindent\textbf{Extended Data Figure 2}
\alttext{Extended Data Figure 2 (PDF page 42).  (a) Row 1, ground-truth conscious face: frontal grayscale macaque face photograph on a uniform mid-grey square; pale muzzle and cheeks, dark eye rings, darker fur on the crown. Row 1, reconstruction from average activity: a smooth, blurred, symmetric generic macaque face on a mid-grey square, low contrast, no fur texture; it does not clearly resemble either ground-truth face of the pair. Row 1, ground-truth suppressed face: frontal grayscale macaque face photograph on a mid-grey square; fluffy dark fur framing the face, lighter brow, pale muzzle. Row 2, ground-truth conscious face: frontal grayscale macaque face photograph; dark fur around the face, narrow pale muzzle, dark eyes. Row 2, reconstruction from average activity: smooth, blurred generic macaque face on a mid-grey square, similar to the other reconstructions in this column. Row 2, ground-truth suppressed face: frontal grayscale macaque face photograph; light brow and forehead, broad pale muzzle, large dark eyes. Row 3, ground-truth conscious face: frontal grayscale macaque face photograph; very light, bright fur on the crown and cheeks, dark eye patches. Row 3, reconstruction from average activity: smooth, blurred generic macaque face on a mid-grey square. Row 3, ground-truth suppressed face: frontal grayscale macaque face photograph; darker fur, pale muzzle, dark eyes. Row 4, ground-truth conscious face: frontal grayscale macaque face photograph; dark fur, wide face, pale muzzle. Row 4, reconstruction from average activity: smooth, blurred generic macaque face on a mid-grey square. Row 4, ground-truth suppressed face: frontal grayscale macaque face photograph; bright white brow band, lighter fur, large eyes. Row 5, ground-truth conscious face: frontal grayscale macaque face photograph; dark fur, rounded head, pale muzzle. Row 5, reconstruction from average activity: smooth, blurred generic macaque face on a mid-grey square. Row 5, ground-truth suppressed face: frontal grayscale macaque face photograph; pale, bright fur with a darker central face and dark eyes. (b) Line plot of explained variance of features in percent (y axis, -100 to 100) against feature number 1 to 50 (x axis); features 1-20 are labelled Shape and 21-50 Appearance by brackets above the plot. Four series. Black solid, priming average: about 76-86 percent for the first three shape features, falling to 10-20 percent by features 15-20 (minimum -7 at feature 15), then 60-98 percent for the appearance features (mean 79). Black dashed, priming average predicting the flashed face: -66 to -84 percent at features 1-3 rising to about -10 to 0 at features 17-20; for the appearance features it lies below -100 and is off scale, apart from two brief spikes up to -58 (feature 24) and -29 (feature 40). Grey solid, binocular conscious average: starts at 38 percent, mostly 0-25 percent over shape features, 10-30 percent over appearance features with one peak of 50 at feature 24 (mean 14). Grey dashed, binocular suppressed average: 32 percent at feature 1 falling to around 0 by feature 15-20 and mostly between -22 and +15 over the appearance features (mean -8). (c) Line plot of decoding accuracy in percent (y axis, 0 to 100) against number of faces 1 to 60 (x axis), each line with a lighter shaded band of the same colour. All lines start at 100 percent for one face. Black solid, priming average: decays slowly, 89 percent at 10 faces, 82 at 20, 73 at 40 and 70 at 60 faces, band roughly plus or minus 5-10 points. Grey solid, binocular conscious average: 72 percent at 2 faces, 27 at 10, 16 at 20, 7 at 40 and 3 at 60 faces, with a wide band (for example 13-41 at 10 faces). Grey dashed, binocular suppressed average: slightly below the grey solid line at first (62 at 2 faces, 20 at 10, 11 at 20) and about the same later (7 at 40, 5 at 60). Black dashed, priming average predicting the flashed face: collapses immediately, 23 percent at 2 faces, 1 percent at 5 faces and 0 from about 9 faces onward, with a narrow band near the bottom axis.}
\figuredata{b}{figures/edfig2b.csv}{figures/edfig2b.py}{vector paths in the PDF (tier 2, exact to the PDF's precision), axes calibrated on the tick marks whose values were read off the rendered panel; details in the CSV header}
\figuredata{c}{figures/edfig2c.csv}{figures/edfig2c.py}{vector paths in the PDF (tier 2, exact to the PDF's precision), axes calibrated on the tick marks whose values were read off the rendered panel; details in the CSV header}

\pdfpage{43}
\begin{center}
\includegraphics[width=0.95\linewidth,height=0.8\textheight,keepaspectratio]{figures/edfig3.png}
\end{center}
\noindent\textbf{Extended Data Figure 3}
\alttext{Extended Data Figure 3 (PDF page 43).  (a) Monkey A, row 1, ground-truth conscious face: frontal grayscale macaque face photograph on a uniform mid-grey square; pale muzzle and cheeks, dark eye rings, darker crown. Monkey A, row 1, reconstruction from state 1 activity: smooth, blurred macaque face on a mid-grey square with pale muzzle and dark eye rings, resembling the conscious face of this row. Monkey A, row 1, ground-truth suppressed face: frontal grayscale macaque face photograph; fluffy dark fur framing the face, lighter brow, pale muzzle. Monkey A, row 1, reconstruction from state 2 activity: smooth, blurred macaque face on a mid-grey square, resembling the suppressed face of this row more than the conscious one. Monkey A, row 2, ground-truth conscious face: frontal grayscale macaque face photograph; dark fur around the face, narrow pale muzzle, dark eyes. Monkey A, row 2, reconstruction from state 1 activity: smooth, blurred face with dark surrounding fur and a pale muzzle, like the conscious face of this row. Monkey A, row 2, ground-truth suppressed face: frontal grayscale macaque face photograph; light brow and forehead, broad pale muzzle, large dark eyes. Monkey A, row 2, reconstruction from state 2 activity: smooth, blurred face with a light brow and broad pale muzzle, like the suppressed face of this row. Monkey A, row 3, ground-truth conscious face: frontal grayscale macaque face photograph; very light, bright fur on the crown and cheeks, dark eye patches. Monkey A, row 3, reconstruction from state 1 activity: smooth, blurred face with light fur around the face, like the conscious face of this row. Monkey A, row 3, ground-truth suppressed face: frontal grayscale macaque face photograph; darker fur, pale muzzle, dark eyes. Monkey A, row 3, reconstruction from state 2 activity: smooth, blurred face with darker surrounding fur, like the suppressed face of this row. Monkey A, row 4, ground-truth conscious face: frontal grayscale macaque face photograph; dark fur, wide face, pale muzzle. Monkey A, row 4, reconstruction from state 1 activity: smooth, blurred face with dark fur and a pale muzzle, like the conscious face of this row. Monkey A, row 4, ground-truth suppressed face: frontal grayscale macaque face photograph; bright white brow band, lighter fur, large eyes. Monkey A, row 4, reconstruction from state 2 activity: smooth, blurred face with a pale brow band and lighter fur, like the suppressed face of this row. Monkey A, row 5, ground-truth conscious face: frontal grayscale macaque face photograph; dark fur, rounded head, pale muzzle. Monkey A, row 5, reconstruction from state 1 activity: smooth, blurred, darker rounded face with a pale muzzle, like the conscious face of this row. Monkey A, row 5, ground-truth suppressed face: frontal grayscale macaque face photograph; pale, bright fur with a darker central face and dark eyes. Monkey A, row 5, reconstruction from state 2 activity: smooth, blurred face with pale fur and a darker centre, like the suppressed face of this row. (b) Monkey A. Line plot of explained variance of features in percent (y axis, -100 to 100) against feature number 1 to 50 (x axis); features 1-20 are labelled Shape and 21-50 Appearance by brackets above the plot. Five series. Black solid, priming average: 68-79 percent for the first three shape features, falling to about 0-20 percent by features 12-20, then 24-90 percent for the appearance features (mean 65). Orange solid, binocular conscious state 1: 43-56 percent at features 1-3, then mostly between -20 and +20 (mean 8 over shape, 13 over appearance features, maximum 41). Purple dashed, binocular suppressed state 2: follows the orange solid line closely, 36-52 percent at features 1-3 then mostly -15 to +30 (mean 5 over shape, 12 over appearance features). Purple solid, binocular conscious state 2: -26 to -43 percent at features 1-3, around -20 to +10 over the rest of the shape features, and strongly negative over the appearance features, between -72 and -10 (mean -47). Orange dashed, binocular suppressed state 1: lowest of all, -45 to -55 percent at features 1-3, -10 to -40 over the rest of the shape features, and -14 to -96 over the appearance features (mean -60). (c) Monkey A. Line plot of decoding accuracy in percent (y axis, 0 to 100) against number of faces 1 to 60 (x axis), each line with a lighter shaded band of the same colour; all lines start at 100 percent for one face. Black solid, priming average: 75 percent at 10 faces, 63 at 20, 51 at 40, 45 at 60 (band about plus or minus 5-14). Orange solid, binocular conscious state 1, and purple dashed, binocular suppressed state 2, run close together in the middle: 42 and 38 percent at 10 faces, 29 and 24 at 20, 21 and 16 at 40, 18 and 12 at 60 faces, bands roughly plus or minus 4-14. Purple solid, binocular conscious state 2, and orange dashed, binocular suppressed state 1, drop fastest and run together near the bottom: 8 and 6 percent at 10 faces, 4 and 4 at 20, 1 and 2 at 40, 0 and 2 at 60 faces, with bands reaching down to 0. (d) Monkey M, row 1, ground-truth conscious face: frontal grayscale macaque face photograph on a uniform mid-grey square; pale muzzle and cheeks, dark eye rings, darker crown. Monkey M, row 1, reconstruction from state 1 activity: smooth, blurred macaque face on a mid-grey square with pale muzzle and dark eye rings, resembling the conscious face of this row. Monkey M, row 1, ground-truth suppressed face: frontal grayscale macaque face photograph; fluffy dark fur framing the face, lighter brow, pale muzzle. Monkey M, row 1, reconstruction from state 2 activity: smooth, blurred macaque face on a mid-grey square, resembling the suppressed face of this row more than the conscious one. Monkey M, row 2, ground-truth conscious face: frontal grayscale macaque face photograph; dark fur around the face, narrow pale muzzle, dark eyes. Monkey M, row 2, reconstruction from state 1 activity: smooth, blurred face with dark surrounding fur and a pale muzzle, like the conscious face of this row. Monkey M, row 2, ground-truth suppressed face: frontal grayscale macaque face photograph; light brow and forehead, broad pale muzzle, large dark eyes. Monkey M, row 2, reconstruction from state 2 activity: smooth, blurred face with a light brow and broad pale muzzle, like the suppressed face of this row. Monkey M, row 3, ground-truth conscious face: frontal grayscale macaque face photograph; very light, bright fur on the crown and cheeks, dark eye patches. Monkey M, row 3, reconstruction from state 1 activity: smooth, blurred face with light fur around the face, like the conscious face of this row. Monkey M, row 3, ground-truth suppressed face: frontal grayscale macaque face photograph; darker fur, pale muzzle, dark eyes. Monkey M, row 3, reconstruction from state 2 activity: smooth, blurred face with darker surrounding fur, like the suppressed face of this row. Monkey M, row 4, ground-truth conscious face: frontal grayscale macaque face photograph; dark fur, wide face, pale muzzle. Monkey M, row 4, reconstruction from state 1 activity: smooth, blurred face with dark fur and a pale muzzle, like the conscious face of this row. Monkey M, row 4, ground-truth suppressed face: frontal grayscale macaque face photograph; bright white brow band, lighter fur, large eyes. Monkey M, row 4, reconstruction from state 2 activity: smooth, blurred face with a pale brow band and lighter fur, like the suppressed face of this row. Monkey M, row 5, ground-truth conscious face: frontal grayscale macaque face photograph; dark fur, rounded head, pale muzzle. Monkey M, row 5, reconstruction from state 1 activity: smooth, blurred, darker rounded face with a pale muzzle, like the conscious face of this row. Monkey M, row 5, ground-truth suppressed face: frontal grayscale macaque face photograph; pale, bright fur with a darker central face and dark eyes. Monkey M, row 5, reconstruction from state 2 activity: smooth, blurred face with pale fur and a darker centre, like the suppressed face of this row. (e) Monkey M. Line plot of explained variance of features in percent (y axis, -100 to 100) against feature number 1 to 50 (x axis); features 1-20 labelled Shape and 21-50 Appearance. Five series. Black solid, priming average: 63-77 percent at features 1-3, decreasing to about 0-15 percent by features 12-20, then 23-98 percent over the appearance features (mean 80). Orange solid, binocular conscious state 1: 15-38 percent at features 1-3, around -30 to +20 over the rest of the shape features, and clearly positive over the appearance features, mostly 15-55 percent (mean 30). Purple dashed, binocular suppressed state 2: 9-33 percent at the first shape features, then around -25 to +20 over shape features and 0-31 over appearance features (mean 8), lying between the orange solid and purple solid lines. Purple solid, binocular conscious state 2: 8 to -20 over shape features, mostly negative over the appearance features, between -56 and +21 (mean -26). Orange dashed, binocular suppressed state 1: -8 to -32 over shape features and strongly negative over the appearance features, mostly -40 to -100 (mean -63), touching the bottom of the axis near features 24-25. (f) Monkey M. Line plot of decoding accuracy in percent (y axis, 0 to 100) against number of faces 1 to 60 (x axis), each line with a lighter shaded band; all lines start at 100 percent for one face. Black solid, priming average: 82 percent at 10 faces, 72 at 20, 61 at 40, 55 at 60 (band about plus or minus 5-12). Orange solid, binocular conscious state 1: 42 percent at 10 faces, 33 at 20, 25 at 40, 22 at 60, clearly the second-highest line. Purple dashed, binocular suppressed state 2: 28 at 10 faces, 18 at 20, 8 at 40, 3 at 60. Purple solid, binocular conscious state 2: 21 at 10 faces, 13 at 20, 7 at 40, 5 at 60. Orange dashed, binocular suppressed state 1: 18 at 10 faces, 10 at 20, 5 at 40, 2 at 60. The three lower lines run close together with overlapping bands (roughly plus or minus 5-15 points) and separate only slightly, the purple dashed line lying above the purple solid and orange dashed lines for most of the range.}
\figuredata{b}{figures/edfig3b.csv}{figures/edfig3b.py}{vector paths in the PDF (tier 2, exact to the PDF's precision), axes calibrated on the tick marks whose values were read off the rendered panel; details in the CSV header}
\figuredata{c}{figures/edfig3c.csv}{figures/edfig3c.py}{vector paths in the PDF (tier 2, exact to the PDF's precision), axes calibrated on the tick marks whose values were read off the rendered panel; details in the CSV header}
\figuredata{e}{figures/edfig3e.csv}{figures/edfig3e.py}{vector paths in the PDF (tier 2, exact to the PDF's precision), axes calibrated on the tick marks whose values were read off the rendered panel; details in the CSV header}
\figuredata{f}{figures/edfig3f.csv}{figures/edfig3f.py}{vector paths in the PDF (tier 2, exact to the PDF's precision), axes calibrated on the tick marks whose values were read off the rendered panel; details in the CSV header}

\pdfpage{44}
\begin{center}
\includegraphics[width=0.95\linewidth,height=0.8\textheight,keepaspectratio]{figures/edfig4.png}
\end{center}
\noindent\textbf{Extended Data Figure 4}
\alttext{Extended Data Figure 4 (PDF page 44).  Box plot of explained variance in percent (y axis from -60 to 80, dotted horizontal line at 0) for five conditions along the x axis: Priming (black), Conscious State 1 (orange), Suppressed State 1 (pale orange), Conscious State 2 (pale purple) and Suppressed State 2 (purple). Each box spans the interquartile range across 9 sessions with a horizontal median line, the whiskers extend to the full range, and the 9 individual sessions are drawn as small jittered dots. Medians with ranges: Priming 57.6 (33.2 to 69.3); Conscious State 1 10.8 (6.4 to 22.7); Suppressed State 1 -33.4 (-55.0 to -28.3); Conscious State 2 -23.0 (-37.0 to -12.6); Suppressed State 2 5.2 (-3.0 to 13.8). Interquartile ranges: 44.9 to 66.7, 7.6 to 20.2, -42.6 to -29.9, -28.4 to -19.3 and 1.8 to 9.7. So in every session state 1 explains the conscious face positively and the suppressed face strongly negatively, state 2 explains the suppressed face positively (8 of 9 sessions above 0) and the conscious face negatively, and priming activity explains the most variance.}
\figuredata{edfig4}{figures/edfig4_boxes.csv}{figures/edfig4.py}{vector paths in the PDF (tier 2, exact to the PDF's precision), axes calibrated on the tick marks whose values were read off the rendered panel; details in the CSV header}
\figuredata{edfig4}{figures/edfig4_points.csv}{figures/edfig4.py}{vector paths in the PDF (tier 2, exact to the PDF's precision), axes calibrated on the tick marks whose values were read off the rendered panel; details in the CSV header}

\pdfpage{45}
\begin{center}
\includegraphics[width=0.95\linewidth,height=0.8\textheight,keepaspectratio]{figures/edfig5.png}
\end{center}
\noindent\textbf{Extended Data Figure 5}
\alttext{Extended Data Figure 5 (PDF page 45).  (a) Ground-truth conscious face, row 1: grayscale frontal photograph-like macaque face on a uniform grey background; broad dark-furred face with a pale muzzle and pale brow, direct gaze. Reconstruction from state-1 activity (leave-one-neuron-out), row 1: smooth, blurred grayscale frontal macaque face on a grey background that resembles the ground-truth conscious face of this row in overall brightness and head shape, without fine fur detail. Ground-truth suppressed face, row 1: grayscale frontal photograph-like macaque face on a uniform grey background; narrower face with lighter fur, dark eyes and a pale muzzle. Reconstruction from state-2 activity (leave-one-neuron-out), row 1: smooth, blurred grayscale frontal macaque face on a grey background that resembles the ground-truth suppressed face of this row in overall brightness and head shape, without fine fur detail. Ground-truth conscious face, row 2: grayscale frontal photograph-like macaque face on a uniform grey background; broad face with a dark brow ridge and pale cheeks. Reconstruction from state-1 activity (leave-one-neuron-out), row 2: smooth, blurred grayscale frontal macaque face on a grey background that resembles the ground-truth conscious face of this row in overall brightness and head shape, without fine fur detail. Ground-truth suppressed face, row 2: grayscale frontal photograph-like macaque face on a uniform grey background; pale face with a dark cap of fur on top and dark eyes. Reconstruction from state-2 activity (leave-one-neuron-out), row 2: smooth, blurred grayscale frontal macaque face on a grey background that resembles the ground-truth suppressed face of this row in overall brightness and head shape, without fine fur detail. Ground-truth conscious face, row 3: grayscale frontal photograph-like macaque face on a uniform grey background; dark face framed by a pale fluffy fur ruff, with prominent dark eyes. Reconstruction from state-1 activity (leave-one-neuron-out), row 3: smooth, blurred grayscale frontal macaque face on a grey background that resembles the ground-truth conscious face of this row in overall brightness and head shape, without fine fur detail. Ground-truth suppressed face, row 3: grayscale frontal photograph-like macaque face on a uniform grey background; dark broad face with a lighter band below the eyes. Reconstruction from state-2 activity (leave-one-neuron-out), row 3: smooth, blurred grayscale frontal macaque face on a grey background that resembles the ground-truth suppressed face of this row in overall brightness and head shape, without fine fur detail. Ground-truth conscious face, row 4: grayscale frontal photograph-like macaque face on a uniform grey background; medium-toned, somewhat narrow face with dark eyes. Reconstruction from state-1 activity (leave-one-neuron-out), row 4: smooth, blurred grayscale frontal macaque face on a grey background that resembles the ground-truth conscious face of this row in overall brightness and head shape, without fine fur detail. Ground-truth suppressed face, row 4: grayscale frontal photograph-like macaque face on a uniform grey background; pale round face with large dark eyes and dark fur on top. Reconstruction from state-2 activity (leave-one-neuron-out), row 4: smooth, blurred grayscale frontal macaque face on a grey background that resembles the ground-truth suppressed face of this row in overall brightness and head shape, without fine fur detail. Ground-truth conscious face, row 5: grayscale frontal photograph-like macaque face on a uniform grey background; dark face with a pale muzzle and pale brow. Reconstruction from state-1 activity (leave-one-neuron-out), row 5: smooth, blurred grayscale frontal macaque face on a grey background that resembles the ground-truth conscious face of this row in overall brightness and head shape, without fine fur detail. Ground-truth suppressed face, row 5: grayscale frontal photograph-like macaque face on a uniform grey background; pale, round, juvenile-looking face with large dark eyes. Reconstruction from state-2 activity (leave-one-neuron-out), row 5: smooth, blurred grayscale frontal macaque face on a grey background that resembles the ground-truth suppressed face of this row in overall brightness and head shape, without fine fur detail. (b) Line plot of explained variance of face features (\%, y axis from -100 to 100) against feature index (x axis 1 to 50; features 1-20 are shape features and 21-50 appearance features, marked by brackets labelled Shape and Appearance above the plot). Five series: black solid = priming, average; orange solid = binocular conscious, state 1; purple solid = binocular conscious, state 2; orange dashed = binocular suppressed, state 1; purple dashed = binocular suppressed, state 2. Leave-one-neuron-out cross-validation. The black priming line is highest (mean 37 \% over shape features and 79 \% over appearance features, peaks near 98 \%). Orange solid averages 20 \% (shape) and 39 \% (appearance); purple dashed 19 \% and 25 \%. Purple solid is negative (-6 \% shape, -30 \% appearance) and orange dashed is the most negative (-16 \% shape, -58 \% appearance, minimum -91 \%). All lines jump up at feature 21. Legend below the plot. (c) Line plot of decoding accuracy (\%, y axis 0 to 100) against number of faces (x axis 0 to 60), each line with a shaded error band of its own colour (lighter; the legend does not state SD or SEM). Series as in panel b: black solid = priming average, orange solid = binocular conscious state 1, purple solid = binocular conscious state 2, orange dashed = binocular suppressed state 1, purple dashed = binocular suppressed state 2. Leave-one-neuron-out cross-validation. Black declines from 100 \% at 1 face to 84 \% at 20 faces and 72 \% at 60. Orange solid declines to 52 \% at 20 faces and 33 \% at 60; purple dashed runs just below it, 38 \% at 20 faces and 20 \% at 60. Purple solid drops steeply to 6 \% at 20 faces and 3 \% at 60; orange dashed drops to about 1 \% by 20 faces and stays at 0. Bands are widest between 5 and 30 faces (for orange solid 36-52 \% at 30 faces). (d) Ground-truth conscious face, row 1: grayscale frontal photograph-like macaque face on a uniform grey background; broad dark-furred face with a pale muzzle and pale brow, direct gaze. Reconstruction from state-1 activity (50\% high-confidence time points), row 1: smooth, blurred grayscale frontal macaque face on a grey background that resembles the ground-truth conscious face of this row in overall brightness and head shape, without fine fur detail. Ground-truth suppressed face, row 1: grayscale frontal photograph-like macaque face on a uniform grey background; narrower face with lighter fur, dark eyes and a pale muzzle. Reconstruction from state-2 activity (50\% high-confidence time points), row 1: smooth, blurred grayscale frontal macaque face on a grey background that resembles the ground-truth suppressed face of this row in overall brightness and head shape, without fine fur detail. Ground-truth conscious face, row 2: grayscale frontal photograph-like macaque face on a uniform grey background; broad face with a dark brow ridge and pale cheeks. Reconstruction from state-1 activity (50\% high-confidence time points), row 2: smooth, blurred grayscale frontal macaque face on a grey background that resembles the ground-truth conscious face of this row in overall brightness and head shape, without fine fur detail. Ground-truth suppressed face, row 2: grayscale frontal photograph-like macaque face on a uniform grey background; pale face with a dark cap of fur on top and dark eyes. Reconstruction from state-2 activity (50\% high-confidence time points), row 2: smooth, blurred grayscale frontal macaque face on a grey background that resembles the ground-truth suppressed face of this row in overall brightness and head shape, without fine fur detail. Ground-truth conscious face, row 3: grayscale frontal photograph-like macaque face on a uniform grey background; dark face framed by a pale fluffy fur ruff, with prominent dark eyes. Reconstruction from state-1 activity (50\% high-confidence time points), row 3: smooth, blurred grayscale frontal macaque face on a grey background that resembles the ground-truth conscious face of this row in overall brightness and head shape, without fine fur detail. Ground-truth suppressed face, row 3: grayscale frontal photograph-like macaque face on a uniform grey background; dark broad face with a lighter band below the eyes. Reconstruction from state-2 activity (50\% high-confidence time points), row 3: smooth, blurred grayscale frontal macaque face on a grey background that resembles the ground-truth suppressed face of this row in overall brightness and head shape, without fine fur detail. Ground-truth conscious face, row 4: grayscale frontal photograph-like macaque face on a uniform grey background; medium-toned, somewhat narrow face with dark eyes. Reconstruction from state-1 activity (50\% high-confidence time points), row 4: smooth, blurred grayscale frontal macaque face on a grey background that resembles the ground-truth conscious face of this row in overall brightness and head shape, without fine fur detail. Ground-truth suppressed face, row 4: grayscale frontal photograph-like macaque face on a uniform grey background; pale round face with large dark eyes and dark fur on top. Reconstruction from state-2 activity (50\% high-confidence time points), row 4: smooth, blurred grayscale frontal macaque face on a grey background that resembles the ground-truth suppressed face of this row in overall brightness and head shape, without fine fur detail. Ground-truth conscious face, row 5: grayscale frontal photograph-like macaque face on a uniform grey background; dark face with a pale muzzle and pale brow. Reconstruction from state-1 activity (50\% high-confidence time points), row 5: smooth, blurred grayscale frontal macaque face on a grey background that resembles the ground-truth conscious face of this row in overall brightness and head shape, without fine fur detail. Ground-truth suppressed face, row 5: grayscale frontal photograph-like macaque face on a uniform grey background; pale, round, juvenile-looking face with large dark eyes. Reconstruction from state-2 activity (50\% high-confidence time points), row 5: smooth, blurred grayscale frontal macaque face on a grey background that resembles the ground-truth suppressed face of this row in overall brightness and head shape, without fine fur detail. (e) Line plot of explained variance of face features (\%, y axis from -100 to 100) against feature index (x axis 1 to 50; features 1-20 are shape features and 21-50 appearance features, marked by brackets labelled Shape and Appearance above the plot). Five series: black solid = priming, average; orange solid = binocular conscious, state 1; purple solid = binocular conscious, state 2; orange dashed = binocular suppressed, state 1; purple dashed = binocular suppressed, state 2. Analysis restricted to the 50 \% high-confidence time points. Black priming line as in panel b (mean 34 \% shape, 77 \% appearance). Orange solid: 7 \% shape, 42 \% appearance (maximum 74 \%). Purple dashed: 8 \% shape, 27 \% appearance. Purple solid: -25 \% shape, -54 \% appearance (minimum -98 \%). Orange dashed is the most negative (-43 \% shape, -91 \% appearance) and repeatedly runs below the -100 \% axis limit between features 21 and 50. Legend below the plot. (f) Line plot of decoding accuracy (\%, y axis 0 to 100) against number of faces (x axis 0 to 60), each line with a shaded error band of its own colour (lighter; the legend does not state SD or SEM). Series as in panel b: black solid = priming average, orange solid = binocular conscious state 1, purple solid = binocular conscious state 2, orange dashed = binocular suppressed state 1, purple dashed = binocular suppressed state 2. Analysis restricted to the 50 \% high-confidence time points. Black declines from 100 \% to 82 \% at 20 faces and 72 \% at 60. Orange solid and purple dashed nearly coincide: 50 \% and 46 \% at 20 faces, both 33 \% at 60, with overlapping bands. Purple solid falls to 5 \% at 20 faces and 2 \% at 60; orange dashed to 0.2 \% at 20 faces and 0 thereafter.}
\figuredata{b}{figures/edfig5b.csv}{figures/edfig5b.py}{vector paths in the PDF (tier 2, exact to the PDF's precision), axes calibrated on the tick marks whose values were read off the rendered panel; details in the CSV header}
\figuredata{c}{figures/edfig5c.csv}{figures/edfig5c.py}{vector paths in the PDF (tier 2, exact to the PDF's precision), axes calibrated on the tick marks whose values were read off the rendered panel; details in the CSV header}
\figuredata{e}{figures/edfig5e.csv}{figures/edfig5e.py}{vector paths in the PDF (tier 2, exact to the PDF's precision), axes calibrated on the tick marks whose values were read off the rendered panel; details in the CSV header}
\figuredata{f}{figures/edfig5f.csv}{figures/edfig5f.py}{vector paths in the PDF (tier 2, exact to the PDF's precision), axes calibrated on the tick marks whose values were read off the rendered panel; details in the CSV header}

\pdfpage{46}
\begin{center}
\includegraphics[width=0.95\linewidth,height=0.8\textheight,keepaspectratio]{figures/edfig6.png}
\end{center}
\noindent\textbf{Extended Data Figure 6}
\alttext{Extended Data Figure 6 (PDF page 46).  (a) Ground-truth conscious face, row 1: grayscale frontal photograph-like macaque face on a uniform grey background; broad dark-furred face with a pale muzzle and pale brow, direct gaze. Reconstruction from state-1 activity (balanced-trial control), row 1: smooth, blurred grayscale frontal macaque face on a grey background that resembles the ground-truth conscious face of this row in overall brightness and head shape, without fine fur detail. Ground-truth suppressed face, row 1: grayscale frontal photograph-like macaque face on a uniform grey background; narrower face with lighter fur, dark eyes and a pale muzzle. Reconstruction from state-2 activity (balanced-trial control), row 1: smooth, blurred grayscale frontal macaque face on a grey background that resembles the ground-truth suppressed face of this row in overall brightness and head shape, without fine fur detail. Ground-truth conscious face, row 2: grayscale frontal photograph-like macaque face on a uniform grey background; broad face with a dark brow ridge and pale cheeks. Reconstruction from state-1 activity (balanced-trial control), row 2: smooth, blurred grayscale frontal macaque face on a grey background that resembles the ground-truth conscious face of this row in overall brightness and head shape, without fine fur detail. Ground-truth suppressed face, row 2: grayscale frontal photograph-like macaque face on a uniform grey background; pale face with a dark cap of fur on top and dark eyes. Reconstruction from state-2 activity (balanced-trial control), row 2: smooth, blurred grayscale frontal macaque face on a grey background that resembles the ground-truth suppressed face of this row in overall brightness and head shape, without fine fur detail. Ground-truth conscious face, row 3: grayscale frontal photograph-like macaque face on a uniform grey background; dark face framed by a pale fluffy fur ruff, with prominent dark eyes. Reconstruction from state-1 activity (balanced-trial control), row 3: smooth, blurred grayscale frontal macaque face on a grey background that resembles the ground-truth conscious face of this row in overall brightness and head shape, without fine fur detail. Ground-truth suppressed face, row 3: grayscale frontal photograph-like macaque face on a uniform grey background; dark broad face with a lighter band below the eyes. Reconstruction from state-2 activity (balanced-trial control), row 3: smooth, blurred grayscale frontal macaque face on a grey background that resembles the ground-truth suppressed face of this row in overall brightness and head shape, without fine fur detail. Ground-truth conscious face, row 4: grayscale frontal photograph-like macaque face on a uniform grey background; medium-toned, somewhat narrow face with dark eyes. Reconstruction from state-1 activity (balanced-trial control), row 4: smooth, blurred grayscale frontal macaque face on a grey background that resembles the ground-truth conscious face of this row in overall brightness and head shape, without fine fur detail. Ground-truth suppressed face, row 4: grayscale frontal photograph-like macaque face on a uniform grey background; pale round face with large dark eyes and dark fur on top. Reconstruction from state-2 activity (balanced-trial control), row 4: smooth, blurred grayscale frontal macaque face on a grey background that resembles the ground-truth suppressed face of this row in overall brightness and head shape, without fine fur detail. Ground-truth conscious face, row 5: grayscale frontal photograph-like macaque face on a uniform grey background; dark face with a pale muzzle and pale brow. Reconstruction from state-1 activity (balanced-trial control), row 5: smooth, blurred grayscale frontal macaque face on a grey background that resembles the ground-truth conscious face of this row in overall brightness and head shape, without fine fur detail. Ground-truth suppressed face, row 5: grayscale frontal photograph-like macaque face on a uniform grey background; pale, round, juvenile-looking face with large dark eyes. Reconstruction from state-2 activity (balanced-trial control), row 5: smooth, blurred grayscale frontal macaque face on a grey background that resembles the ground-truth suppressed face of this row in overall brightness and head shape, without fine fur detail. (b) Line plot of explained variance of face features (\%, y axis from -100 to 100) against feature index (x axis 1 to 50; features 1-20 are shape features and 21-50 appearance features, marked by brackets labelled Shape and Appearance above the plot). Five series: black solid = priming, average; orange solid = binocular conscious, state 1; purple solid = binocular conscious, state 2; orange dashed = binocular suppressed, state 1; purple dashed = binocular suppressed, state 2. Balanced-trial control (equal numbers of state-1 and state-2 time points per trial). Black priming line highest (mean 37 \% shape, 79 \% appearance, peaks near 98 \%). Orange solid averages 15 \% (shape) and 38 \% (appearance); purple dashed 12 \% and 21 \%. Purple solid is negative (-8 \% shape, -29 \% appearance); orange dashed most negative (-16 \% shape, -60 \% appearance, minimum -96 \%). Legend below the plot. (c) Line plot of decoding accuracy (\%, y axis 0 to 100) against number of faces (x axis 0 to 60), each line with a shaded error band of its own colour (lighter; the legend does not state SD or SEM). Series as in panel b: black solid = priming average, orange solid = binocular conscious state 1, purple solid = binocular conscious state 2, orange dashed = binocular suppressed state 1, purple dashed = binocular suppressed state 2. Balanced-trial control. Black declines from 100 \% to 82 \% at 20 faces and 70 \% at 60. Orange solid declines to 45 \% at 20 faces and 27 \% at 60. Purple dashed is lower: 40 \% at 10 faces, 15 \% at 40 and 8 \% at 60. Purple solid falls to 5 \% at 20 faces and 2 \% at 60; orange dashed to about 1 \% at 20 faces and 0 thereafter. (d) Three small line plots of single-trial decoding accuracy (\%, y axis 0 to 100) against time from flash (x axis ticks 0, 1, 2; the original axis label says ms but the values are seconds), arranged as Average (grey, top left), State 1 (orange, bottom left) and State 2 (purple, bottom right), with a legend at top right: solid = Conscious, dashed = Suppressed. Solid lines (decoded identity compared with the conscious stimulus) start near 0-10 \% before the flash and rise within 0.3 s to fluctuate around 36 \% (Average), 39 \% (State 1) and 20 \% (State 2) between 0.5 and 2 s. Dashed lines (compared with the suppressed stimulus) start high, 60-75 \% around the flash, and fall to fluctuate around 29 \% (Average), 19 \% (State 1) and 34 \% (State 2). In State 1 the solid line lies above the dashed one, in State 2 the dashed line lies above the solid one. Chance is 10 \%. Dashed traces are interpolated across short stretches that could not be recovered from the PDF. (e) Paired line plot of single-trial decoding accuracy (\%, y axis 0 to 100) at three x categories: Average, State 1, State 2. Thin grey lines are the ten sessions, thick black lines the mean across sessions; solid lines with filled circles = compared with the conscious stimulus, dashed lines with open circles = compared with the suppressed stimulus; circles are grey (Average), orange (State 1) and purple (State 2). Conscious mean: 38.4 \% (Average), 39.9 \% (State 1), 18.7 \% (State 2). Suppressed mean: 28.4 \%, 13.3 \%, 37.1 \%. The two mean lines cross between State 1 and State 2. Session values span 29-55 \% (conscious, Average), 32-56 \% (conscious, State 1), 4-30 \% (conscious, State 2), 15-40 \% (suppressed, Average), 2-43 \% (suppressed, State 1) and 29-47 \% (suppressed, State 2). Legend at top left.}
\figuredata{b}{figures/edfig6b.csv}{figures/edfig6b.py}{vector paths in the PDF (tier 2, exact to the PDF's precision), axes calibrated on the tick marks whose values were read off the rendered panel; details in the CSV header}
\figuredata{c}{figures/edfig6c.csv}{figures/edfig6c.py}{vector paths in the PDF (tier 2, exact to the PDF's precision), axes calibrated on the tick marks whose values were read off the rendered panel; details in the CSV header}
\figuredata{d}{figures/edfig6d.csv}{figures/edfig6d.py}{vector paths in the PDF (tier 2, exact to the PDF's precision), axes calibrated on the tick marks whose values were read off the rendered panel; details in the CSV header}
\figuredata{e}{figures/edfig6e.csv}{figures/edfig6e.py}{vector paths in the PDF (tier 2, exact to the PDF's precision), axes calibrated on the tick marks whose values were read off the rendered panel; details in the CSV header}

\pdfpage{47}
\begin{center}
\includegraphics[width=0.95\linewidth,height=0.8\textheight,keepaspectratio]{figures/edfig7.png}
\end{center}
\noindent\textbf{Extended Data Figure 7}
\alttext{Extended Data Figure 7 (PDF page 47).  (a) Ground-truth conscious face (shuffle-reassigned template A-prime), row 1: grayscale frontal macaque face on a uniform grey background, dark fur on the crown and cheeks, broad pale muzzle, eyes looking straight at the viewer. Reconstruction from state 1 activity, row 1: smooth, low-detail frontal macaque face on grey, dark eyes and dark fur outline around a pale face; it looks like an average face and does not resemble the reassigned template. Ground-truth suppressed face (shuffle-reassigned template B-prime), row 1: macaque with dark eye rings, a pale lower face and light fluffy fur on both sides, looking straight ahead. Reconstruction from state 2 activity, row 1: smooth, low-detail frontal macaque face on grey, nearly identical to the state 1 reconstruction, not resembling the reassigned template. Ground-truth conscious face, row 2: macaque with a wide pale face, dark brows and deep-set dark eyes, framed by fur on both sides. Reconstruction from state 1 activity, row 2: smooth average-looking frontal macaque face on grey, dark eyes, pale face, dark fur outline. Ground-truth suppressed face, row 2: macaque with a round face, large dark eyes, dark crown and thin fur at the sides, head slightly tilted. Reconstruction from state 2 activity, row 2: smooth average-looking frontal macaque face on grey, like the other reconstructions. Ground-truth conscious face, row 3: frontal macaque face with dark fur above the brow, small dark eyes and a soft pale muzzle. Reconstruction from state 1 activity, row 3: smooth average-looking frontal macaque face on grey. Ground-truth suppressed face, row 3: macaque with a light face, wide-set dark eyes, dark crown fur and a slightly elongated snout. Reconstruction from state 2 activity, row 3: smooth average-looking frontal macaque face on grey. Ground-truth conscious face, row 4: macaque with a pale bright face, dark crown, prominent dark eyes and a narrow chin. Reconstruction from state 1 activity, row 4: smooth average-looking frontal macaque face on grey. Ground-truth suppressed face, row 4: macaque with a dark facial mask, pale muzzle and a lighter fur cap, head tilted. Reconstruction from state 2 activity, row 4: smooth average-looking frontal macaque face on grey. Ground-truth conscious face, row 5: macaque with a large bright face, close-set dark eyes and sparse fur. Reconstruction from state 1 activity, row 5: smooth average-looking frontal macaque face on grey. Ground-truth suppressed face, row 5: macaque with a dark shaggy face, pale muzzle and small eyes, head tilted downward. Reconstruction from state 2 activity, row 5: smooth average-looking frontal macaque face on grey. (b) Line plot of explained variance of face features in percent (y axis, -100 to 100) against feature number 1 to 50 (x axis; features 1-20 bracketed as Shape, 21-50 as Appearance), for the shuffle-template control. Five series: priming average (black solid) stays high, about 30-85 percent for shape features, dropping to 0 at feature 15, and 60-97 percent for appearance features. Binocular conscious state 1 (orange solid) fluctuates around -20 to +10 percent; binocular conscious state 2 (purple solid) around -30 to -50 percent with dips to -80 and -90 percent at features 3 and 26; binocular suppressed state 1 (orange dashed) around -30 to -50 percent; binocular suppressed state 2 (purple dashed) around -10 to +20 percent. None of the binocular series is consistently positive. Legend below: Priming, average; Binocular conscious, state 1 and 2; Binocular suppressed, state 1 and 2. (c) Line plot of decoding accuracy in percent (y axis, 0 to 100) against number of faces 1 to 60 (x axis), shuffle-template control, with shaded bands around each line. The priming average (black) declines slowly from 100 percent at 1 face to 70 percent at 60 faces with a band of about +/-5 percent. All four binocular series start at 100 percent at 1 face and collapse within a few faces: binocular conscious state 1 (orange solid) and state 2 (purple solid) fall below 10 percent by 10 faces and to about 0 percent by 30 faces; binocular suppressed state 1 (orange dashed) and state 2 (purple dashed) fall to about 5 percent by 20 faces and 1-2 percent at 60 faces, with wide light bands. This is chance-level decoding. (d) Ground-truth conscious face (actually presented face A), row 1: grayscale frontal macaque face on grey, wide pale cheeks, dark eyes, dark crown. Reconstruction from state 1 activity with shuffled templates, row 1: smooth, low-detail frontal macaque face on grey (virtually identical to the panel a reconstruction), which does not match the presented face. Ground-truth suppressed face (actually presented face B), row 1: macaque with a pale slightly elongated face and big dark round eyes. Reconstruction from state 2 activity with shuffled templates, row 1: smooth, low-detail frontal macaque face on grey (virtually identical to the panel a reconstruction), not matching the presented face. Ground-truth conscious face, row 2: macaque with a light face, dark eye sockets and fur parted at the crown. Reconstruction from state 1 activity, row 2: smooth average-looking frontal macaque face on grey. Ground-truth suppressed face, row 2: macaque with a dark fur cap, pale face and large round eyes with strong dark rings. Reconstruction from state 2 activity, row 2: smooth average-looking frontal macaque face on grey. Ground-truth conscious face, row 3: macaque with a broad face, dark fur on the sides, small eyes and a pale muzzle. Reconstruction from state 1 activity, row 3: smooth average-looking frontal macaque face on grey. Ground-truth suppressed face, row 3: macaque with an oval pale face, dark eyes and dark fur at the sides. Reconstruction from state 2 activity, row 3: smooth average-looking frontal macaque face on grey. Ground-truth conscious face, row 4: macaque with a long pale face, dark fur cap and dark eyes, head tilted to the left. Reconstruction from state 1 activity, row 4: smooth average-looking frontal macaque face on grey. Ground-truth suppressed face, row 4: macaque with a bright face, dark eyes and a fur ruff on both sides, slightly tilted. Reconstruction from state 2 activity, row 4: smooth average-looking frontal macaque face on grey. Ground-truth conscious face, row 5: macaque with a pale broad face, dark eyes and thin fur. Reconstruction from state 1 activity, row 5: smooth average-looking frontal macaque face on grey. Ground-truth suppressed face, row 5: macaque with a dark fringe, pale muzzle and big eyes, head tilted. Reconstruction from state 2 activity, row 5: smooth average-looking frontal macaque face on grey. (e) Line plot of explained variance of face features in percent (y axis, -100 to 100) against feature number 1 to 50 (x axis; features 1-20 bracketed as Shape, 21-50 as Appearance), shuffle-template reconstructions compared with the actually presented faces. Priming average (black solid) is the same as in panel b: about 30-85 percent for shape features (dip to 0 at feature 15) and 60-97 percent for appearance features. Binocular conscious state 1 (orange solid) and state 2 (purple solid) are weakly positive, mostly 0 to +25 percent for shape features (with a dip to -40 percent at feature 15 for state 2) and 0 to +20 percent for appearance features. Binocular suppressed state 1 (orange dashed) hovers around 0 to -15 percent and binocular suppressed state 2 (purple dashed) around -10 to -35 percent for appearance features. Same legend as panel b. (f) Line plot of decoding accuracy in percent (y axis, 0 to 100) against number of faces 1 to 60 (x axis), shuffle-template reconstructions compared with the actually presented faces, with shaded bands. Priming average (black) declines slowly from 100 percent to 70 percent at 60 faces. The binocular series decay much faster but stay above the chance-level curves of panel c: binocular conscious state 2 (purple solid) is highest, about 50 percent at 5 faces, 20 percent at 20 faces and 7 percent at 60 faces; binocular conscious state 1 (orange solid) is slightly lower, about 5 percent at 60 faces; binocular suppressed state 1 (orange dashed) ends near 5 percent; binocular suppressed state 2 (purple dashed) is lowest, about 10 percent at 20 faces and 2 percent at 60 faces. Bands are about +/-10 to 15 percent wide at 10 faces.}
\figuredata{b}{figures/edfig7b.csv}{figures/edfig7b.py}{vector paths in the PDF (tier 2, exact to the PDF's precision), axes calibrated on the tick marks whose values were read off the rendered panel; details in the CSV header}
\figuredata{c}{figures/edfig7c.csv}{figures/edfig7c.py}{vector paths in the PDF (tier 2, exact to the PDF's precision), axes calibrated on the tick marks whose values were read off the rendered panel; details in the CSV header}
\figuredata{e}{figures/edfig7e.csv}{figures/edfig7e.py}{vector paths in the PDF (tier 2, exact to the PDF's precision), axes calibrated on the tick marks whose values were read off the rendered panel; details in the CSV header}
\figuredata{f}{figures/edfig7f.csv}{figures/edfig7f.py}{vector paths in the PDF (tier 2, exact to the PDF's precision), axes calibrated on the tick marks whose values were read off the rendered panel; details in the CSV header}

\pdfpage{48}
\begin{center}
\includegraphics[width=0.95\linewidth,height=0.8\textheight,keepaspectratio]{figures/edfig8.png}
\end{center}
\noindent\textbf{Extended Data Figure 8}
\alttext{Extended Data Figure 8 (PDF page 48).  (left) Overlaid step histograms of the cosine angle (x axis from -1 to 1, thin grey vertical line at 0) between each cell's state-1 feature-tuning axis and its priming axis; y axis Cell count from 0 to 50; 40 bins of width 0.05; 548 cells per series. Orange, State 1 conscious: occupies -0.40 to 1, peaks at 45 cells in the 0.55 to 0.60 bin, 93 percent of cells above 0 and 55 percent above 0.5. Pale orange, State 1 suppressed: shifted to the left, occupies -0.65 to 0.85 with a broad peak of 41 cells at 0.15 to 0.20, only 66 percent above 0 and 10 percent above 0.5. Black, Priming half-trials (noise ceiling): occupies -0.55 to 1, peaks at 49 cells in the 0.75 to 0.80 bin, 61 percent above 0.5. The orange conscious distribution lies almost on top of the black noise-ceiling distribution, whereas the pale orange suppressed distribution is clearly shifted towards lower cosine values. Legend at right: State 1 conscious, State 1 suppressed, Priming half-trials. (right) Overlaid step histograms of the cosine angle (x axis from -1 to 1, thin grey vertical line at 0) between each cell's state-2 feature-tuning axis and its priming axis; y axis from 0 to 60 cells (no axis title); 40 bins of width 0.05; 548 cells per series. Pale purple, State 2 conscious: shifted to the left, occupies -0.55 to 0.90, peaks at 41 cells in the 0.25 to 0.30 bin, 73 percent above 0 and 21 percent above 0.5. Purple, State 2 suppressed: occupies -0.55 to 1, peaks at 55 cells in the 0.65 to 0.70 bin, 94 percent above 0 and 57 percent above 0.5. Black, Priming half-trials (noise ceiling, same data as in the left panel): peaks at 49 cells in the 0.75 to 0.80 bin, 61 percent above 0.5. The purple suppressed distribution overlaps the black noise-ceiling distribution, while the pale purple conscious distribution sits at clearly lower cosine values. Legend at right: State 2 conscious, State 2 suppressed, Priming half-trials.}
\figuredata{left}{figures/edfig8_left_state1.csv}{figures/edfig8_left_state1.py}{vector paths in the PDF (tier 2, exact to the PDF's precision), axes calibrated on the tick marks whose values were read off the rendered panel; details in the CSV header}
\figuredata{right}{figures/edfig8_right_state2.csv}{figures/edfig8_right_state2.py}{vector paths in the PDF (tier 2, exact to the PDF's precision), axes calibrated on the tick marks whose values were read off the rendered panel; details in the CSV header}

\pdfpage{49}
\begin{center}
\includegraphics[width=0.95\linewidth,height=0.8\textheight,keepaspectratio]{figures/edfig9.png}
\end{center}
\noindent\textbf{Extended Data Figure 9}
\alttext{Extended Data Figure 9 (PDF page 49).  (a) Spike raster of one example flash-suppression trial: 429 simultaneously recorded cells as rows (ML cells above a thin horizontal line, AM cells below, label 429 cells on the y axis) against time from the flash in seconds from -1 to 1 on the x axis; 5258 spikes drawn as small black ticks. A vertical black line marks the flash at 0 s. The background is shaded in four colours giving the state inferred for each time bin by the four-state model: pale orange = conscious (state 1) and On, dark tan = conscious and Off, pale lavender = suppressed (state 2) and On, dark grey-purple = suppressed and Off. The shading alternates rapidly (219 epochs, median 4 ms, longest 139 ms). Before the flash the suppressed colours cover 0.89 s of the 1.01 s shown and the conscious colours 0.12 s; after the flash the conscious colours cover 0.43 s of 1.0 s, mostly from about 0.1 s after the flash on, interleaved with suppressed epochs (0.58 s). Spikes are denser in the ML block and are concentrated in a few highly active rows. (b) Row 1 ground-truth conscious face: sharp grey-scale frontal macaque face on a uniform mid-grey background, narrow face with a dark fur crown, pale muzzle and dark eyes. Row 1 reconstruction from state 1: smooth, blurred grey-scale macaque face with the same narrow outline and dark crown as the row 1 conscious face. Row 1 ground-truth suppressed face: sharp grey-scale frontal macaque face, broader and rounder than the conscious face, with pale cheeks and fur spreading sideways. Row 1 reconstruction from state 2: blurred grey-scale macaque face with the broader, rounder outline of the row 1 suppressed face. Row 2 ground-truth conscious face: sharp grey-scale frontal macaque face with wide dark fur on both sides, pale central face and dark eye region. Row 2 reconstruction from state 1: blurred grey-scale macaque face with the wide dark side fur and pale centre of the row 2 conscious face. Row 2 ground-truth suppressed face: sharp grey-scale frontal macaque face with a flat dark fur cap, wide-set bright eyes and a squarish outline. Row 2 reconstruction from state 2: blurred grey-scale macaque face with a squarish outline and flat dark cap like the row 2 suppressed face. Row 3 ground-truth conscious face: sharp grey-scale frontal macaque face with a tall dark fur crown, narrow pale muzzle and dark surrounding fur. Row 3 reconstruction from state 1: blurred grey-scale macaque face, narrow with a dark crown, resembling the row 3 conscious face. Row 3 ground-truth suppressed face: sharp grey-scale frontal macaque face with a broad dark fur mass around a bright central face. Row 3 reconstruction from state 2: blurred grey-scale macaque face, broad with dark surrounding fur like the row 3 suppressed face. Row 4 ground-truth conscious face: sharp grey-scale frontal macaque face with a light fur crown, pale face and small dark eyes. Row 4 reconstruction from state 1: blurred grey-scale macaque face, pale with a light crown, resembling the row 4 conscious face. Row 4 ground-truth suppressed face: sharp grey-scale frontal macaque face with a wide, dark, shaggy fur outline and bright wide-set eyes. Row 4 reconstruction from state 2: blurred grey-scale macaque face, wide with dark side fur like the row 4 suppressed face. Row 5 ground-truth conscious face: sharp grey-scale frontal macaque face with a dark fur cap, pale cheeks and a broad muzzle. Row 5 reconstruction from state 1: blurred grey-scale macaque face with a dark cap and broad muzzle resembling the row 5 conscious face. Row 5 ground-truth suppressed face: sharp grey-scale frontal macaque face with a tall shaggy fur crown and bright eyes on a darker face. Row 5 reconstruction from state 2: blurred grey-scale macaque face with a tall crown and darker face like the row 5 suppressed face. (c) Line plot of explained variance of features in percent (y axis, -100 to 100) against feature number 1 to 50 (x axis; features 1-20 bracketed as Shape, 21-50 as Appearance), five series. Priming average (black solid) is highest: mean 37 pct over shape features and 79 pct over appearance features, peaking at 98 pct at feature 21. Binocular suppressed state 2 (purple dashed) follows it at a lower level: 28 pct (shape) and 55 pct (appearance). Binocular conscious state 1 (orange solid) is intermediate: 19 pct (shape) and 35 pct (appearance). Binocular suppressed state 1 (orange dashed) is negative: -9 pct (shape) and -45 pct (appearance). Binocular conscious state 2 (purple solid) is lowest: -22 pct over shape features and below -50 pct for most appearance features, leaving the axes below -100 pct at 9 of the 30 appearance features. All series are jagged from feature to feature; the four binocular series show a sharp step at the shape/appearance boundary (feature 20/21). Legend below the axes. (d) Line plot of decoding accuracy in percent (y axis, 0 to 100) against number of faces 1 to 60 (x axis), five series with shaded bands of the same colour. All start at 100 pct for one face and fall as more faces are added. Priming average (black solid) stays highest: 89 pct at 10 faces, 82 pct at 20, 70 pct at 60 (band about 74-91 pct at 20 faces). Binocular suppressed state 2 (purple dashed) is next: 76 pct at 10, 66 pct at 20, 48 pct at 60 faces (band 55-76 pct at 20). Binocular conscious state 1 (orange solid): 52 pct at 10, 39 pct at 20, 23 pct at 60 faces (band 27-51 pct at 20). Binocular suppressed state 1 (orange dashed) drops quickly: 45 pct at 2 faces, 7 pct at 10, about 2 pct at 20 and near 0 from 40 faces on. Binocular conscious state 2 (purple solid) collapses fastest: 27 pct at 2 faces, 2 pct at 5, 0 pct from about 12 faces on, with its band also at 0. Legend below the axes. (e) Three stacked histograms of single-cell tuning gain (x axis Gain, -0.5 to 2.5, ticks at 0 and 2; y axis number of cells, ticks at 0 and 100), each overlaying two series as 30 pct transparent fills with step outlines; legend at right. Top: State 1 conscious (orange, 0.2-wide bins, 542 cells) peaks at 128 cells in the 0.4-0.6 bin with a long tail to 2.5, mean gain about 0.75; State 1 suppressed (pale orange, 0.1-wide bins, 543 cells) is narrower and shifted left, peaking at 108 cells in 0.2-0.3, mean about 0.31. Middle: State 2 conscious (pale purple, 0.1-wide bins, 543 cells) peaks at 100 cells in 0.3-0.4, mean about 0.46; State 2 suppressed (purple, 0.2-wide bins, 541 cells) peaks at 156 cells in 0.4-0.6, mean about 0.63. Bottom: State 1 conscious (orange) and State 2 suppressed (purple) overlaid; they nearly coincide, the orange one slightly wider with a heavier tail at gains above 1. (f) Scatter plot of 539 cells: gain for the suppressed face in state 2 (y axis, -0.5 to 2.5, ticks 0 and 2) against gain for the conscious face in state 1 (x axis, same range), with a dashed identity line and dashed lines at x = 0 and y = 0. The points form a tight elongated cloud along the diagonal from about 0.1 to 1.5, mostly just below the identity line (76 pct of points below it; mean x 0.74, mean y 0.62; correlation about 0.88 computed from the extracted points), with a sparse tail reaching 2.5. (g) Scatter plot of 446 cells: normalised mean response in state 2 (y axis, 0 to 3) against normalised mean response in state 1 (x axis, 0 to 3), with a dashed identity line. Points lie tightly along the diagonal between about 0.3 and 2.7, slightly above the identity line (72 pct of points above it; mean x 1.56, mean y 1.63; correlation about 0.96 computed from the extracted points). (h) Line plot of decoding accuracy in percent for discriminating state 1 from state 2 (y axis, 0 to 100, ticks 0, 50, 100) against number of cells (x axis, 0 to 1000), black line with a grey band about +/- 2 to 3 pct wide, and a dashed horizontal chance line at 50 pct. Accuracy rises from 51 pct with 1 cell to 56 pct at 10 cells and 64 pct at 100 cells, dips to 61 pct at 200 cells, then climbs slowly to 65 pct at 500 and 67 pct at 1000 cells. (h, second panel with this letter) Four scatter plots in a 2 x 2 grid of modulation index during binocular rivalry (y axis, -1 to 1) against modulation index during unambiguous presentation (x axis, -1 to 1), one black dot per cell-stimulus pair, each with dashed lines at x = 0 and y = 0 and a dashed diagonal. Top left, All states (2906 points): a flattened cloud centred near y = 0 that rises only gently with x (slope about 0.2 relative to the identity diagonal; correlation about 0.46 from the extracted points). Top right, State 1 (orange title, 2906 points): the cloud is stretched along the identity diagonal (correlation about 0.62). Bottom left, State 2 (purple title, 2906 points): the dashed diagonal is the anti-diagonal y = -x and the cloud tilts weakly downward (correlation about -0.18). Bottom right, Slope-equalized (1256 points): a roughly round cloud centred near the origin with little relation to x (correlation about 0.14).}
\figuredata{a}{figures/edfig9a_spikes.csv}{figures/edfig9a.py}{vector paths in the PDF (tier 2, exact to the PDF's precision), axes calibrated on the tick marks whose values were read off the rendered panel; details in the CSV header}
\figuredata{a}{figures/edfig9a_state_epochs.csv}{figures/edfig9a.py}{vector paths in the PDF (tier 2, exact to the PDF's precision), axes calibrated on the tick marks whose values were read off the rendered panel; details in the CSV header}
\figuredata{c}{figures/edfig9c.csv}{figures/edfig9c.py}{vector paths in the PDF (tier 2, exact to the PDF's precision), axes calibrated on the tick marks whose values were read off the rendered panel; details in the CSV header}
\figuredata{d}{figures/edfig9d.csv}{figures/edfig9d.py}{vector paths in the PDF (tier 2, exact to the PDF's precision), axes calibrated on the tick marks whose values were read off the rendered panel; details in the CSV header}
\figuredata{e}{figures/edfig9e.csv}{figures/edfig9e.py}{vector paths in the PDF (tier 2, exact to the PDF's precision), axes calibrated on the tick marks whose values were read off the rendered panel; details in the CSV header}
\figuredata{f}{figures/edfig9f.csv}{figures/edfig9f.py}{vector paths in the PDF (tier 2, exact to the PDF's precision), axes calibrated on the tick marks whose values were read off the rendered panel; details in the CSV header}
\figuredata{g}{figures/edfig9g.csv}{figures/edfig9g.py}{vector paths in the PDF (tier 2, exact to the PDF's precision), axes calibrated on the tick marks whose values were read off the rendered panel; details in the CSV header}
\figuredata{h}{figures/edfig9h_decoding_vs_cells.csv}{figures/edfig9h_decoding_vs_cells.py}{vector paths in the PDF (tier 2, exact to the PDF's precision), axes calibrated on the tick marks whose values were read off the rendered panel; details in the CSV header}
\figuredata{h}{figures/edfig9h_modulation_index.csv}{figures/edfig9h_modulation_index.py}{vector paths in the PDF (tier 2, exact to the PDF's precision), axes calibrated on the tick marks whose values were read off the rendered panel; details in the CSV header}
\decompnote{Extended Data Figure 9 has two panels labelled h in the PDF; both are described above under (h).}

\pdfpage{50}
\begin{center}
\includegraphics[width=0.95\linewidth,height=0.8\textheight,keepaspectratio]{figures/edfig10.png}
\end{center}
\noindent\textbf{Extended Data Figure 10}
\alttext{Extended Data Figure 10 (PDF page 50).  (a) Table of five models (columns: Real, 2-state, C-only, Average, Uncoord) against four tests (rows), each cell a green check mark (pass) or a red cross (fail). Reconstruction: Real pass; 2-state pass; C-only fail; Average pass; Uncoord pass. Reconstruction (coord): Real pass; 2-state pass; C-only fail; Average fail; Uncoord fail. Reversal: Real pass; 2-state pass; C-only fail; Average fail; Uncoord fail. Log likelihood: Real pass; 2-state pass; C-only fail; Average fail; Uncoord fail. (b) Grouped bar chart of reconstruction explained variance in percent (y axis, -100 to 100, zero line drawn) for the recorded data (Real) and three generative models (2-state, Average, Uncoord) on the x axis, four bars per group with bootstrap error bars: state 1 conscious (orange), state 1 suppressed (light orange), state 2 conscious (light purple), state 2 suppressed (purple). Real: 63, -99, -77 and 59 percent. 2-state: 80, below -100 (bar clipped, printed -183), below -100 (printed -182), 80 percent. Average: 74, -108, -106, 74 percent. Uncoord: 75, -109, -105, 74 percent. So in every case the face that is conscious in state 1 and the face that is suppressed in state 2 are reconstructed well (60-80 percent) while the other face has strongly negative explained variance. Error bars are about +/-1 to 3 percent, +/-9 percent for the Real suppressed and conscious bars. (c) Line plot with circular markers of conscious-face reconstruction explained variance in state 1 in percent (y axis, 45 to 65) against achieved mean HMM dwell time in milliseconds (x axis, 0 to 300), for the recorded data (red) and the averaging model (blue). Each of the nine points per series is labelled with the imposed HMM self-transition probability: 0.50, 0.60, 0.70, 0.80, 0.87, 0.92, 0.95, 0.97, 0.99, with the dwell time growing with the probability. Red (real data): from 63.4 percent at 66 ms (p 0.50) falling gently to 55.8 percent at 223 ms (p 0.99). Blue (average model): from 62.8 percent at 39 ms (p 0.50) falling steeply to 46.1 percent at 258 ms (p 0.99). The blue curve lies below the red one at every dwell time above 60 ms and the gap widens with dwell time. Small horizontal error bars (up to +/-3 ms) on the dwell times. Legend upper right: Real data, Average model. (d) Grouped bar chart of held-out (cross-validated) reconstruction explained variance in percent (y axis, -100 to 100) for Real, 2-state, Average and Uncoord, four bars per group in the same colours as panel b (state 1 conscious orange, state 1 suppressed light orange, state 2 conscious light purple, state 2 suppressed purple) with error bars. Real: 46, -60, -29 and 32 percent. 2-state: 74, below -100 (printed -124), below -100 (printed -128), 73 percent. Average: 14, 16, 13, 15 percent. Uncoord: 15, 13, 16, 13 percent. Only the recorded data and the two-state model show the state-dependent pattern (conscious face positive and suppressed face negative in state 1, reversed in state 2); the averaging and uncoordinated models give four similar small positive bars of about 13-16 percent. Error bars about +/-2 to 4 percent, +/-7 percent for the Real suppressed bars. (e) Bar chart of held-out tuning correlation r (y axis, -1 to 1, zero line drawn) for Real, 2-state, Average and Uncoord on the x axis, two bars per group: state 1 (orange) and state 2 (purple), with bootstrap error bars. Real: state 1 +0.60, state 2 -0.42. 2-state: +0.65 and -0.66. Average: +0.01 and +0.02. Uncoord: +0.01 and +0.01. Thus only the recorded data and the two-state model show a strongly negative state 2 correlation (tuning reversal); the averaging and uncoordinated models give correlations near zero. Error bars about +/-0.01 to 0.015. Legend on top: State 1, State 2. (f) Bar chart of held-out log likelihood difference between a two-state and a one-state hidden Markov model (y axis, -0.15 to 0.1, zero line drawn) for Real, 2-state, Average and Uncoord on the x axis, one dark grey bar per model with s.e.m. error bars across sessions. Real +0.080 (error bar 0.067 to 0.093) and 2-state +0.053 (0.039 to 0.066) are above zero, meaning two states fit better; Average -0.122 and Uncoord -0.122 (both with error bars from -0.136 to -0.108) are below zero, meaning one state fits better for those models.}
\figuredata{a}{figures/edfig10a.csv}{(table; no plot script)}{the check and cross marks are vector paths; pass/fail read from their colours, rows and columns by position (tier 2)}
\figuredata{b}{figures/edfig10b.csv}{figures/edfig10b.py}{vector paths in the PDF (tier 2, exact to the PDF's precision), axes calibrated on the tick marks whose values were read off the rendered panel; details in the CSV header}
\figuredata{c}{figures/edfig10c.csv}{figures/edfig10c.py}{vector paths in the PDF (tier 2, exact to the PDF's precision), axes calibrated on the tick marks whose values were read off the rendered panel; details in the CSV header}
\figuredata{d}{figures/edfig10d.csv}{figures/edfig10d.py}{vector paths in the PDF (tier 2, exact to the PDF's precision), axes calibrated on the tick marks whose values were read off the rendered panel; details in the CSV header}
\figuredata{e}{figures/edfig10e.csv}{figures/edfig10e.py}{vector paths in the PDF (tier 2, exact to the PDF's precision), axes calibrated on the tick marks whose values were read off the rendered panel; details in the CSV header}
\figuredata{f}{figures/edfig10f.csv}{figures/edfig10f.py}{vector paths in the PDF (tier 2, exact to the PDF's precision), axes calibrated on the tick marks whose values were read off the rendered panel; details in the CSV header}

\end{document}
